\documentclass[12pt,letterpaper,reqno]{amsart}
\usepackage{tikz}
\usetikzlibrary{positioning, shapes.geometric, arrows.meta}
\usepackage{amssymb}
\usepackage{amsmath}
\usepackage{amsthm}
\usepackage{amsfonts}
\usepackage{bm}
\IfFileExists{bbm.sty}{\usepackage{bbm}}{}
\usepackage{enumitem}
\usepackage{pgfplots}
\pgfplotsset{compat=1.18}
\usepackage{booktabs,tabularx,array}
\usepackage{graphicx}
\usepackage[T1]{fontenc}
\usepackage{doi}
\usepackage{float}

\usepackage[dvipsnames]{xcolor}
\usepackage{bookmark}
\usepackage{hyperref}
\hypersetup{colorlinks=true,
linkcolor=RoyalBlue,
citecolor=ForestGreen!65!white,
urlcolor=BrickRed}
\addtolength{\hoffset}{-2.00cm}
\addtolength{\textwidth}{3.9cm}
\addtolength{\voffset}{-1.45cm}
\addtolength{\textheight}{2.9cm}
\setlength{\parindent}{12pt}
\setlength{\parskip}{0pt}

\allowdisplaybreaks
\newtheorem{thm}{Theorem}[section]
\newtheorem{lem}[thm]{Lemma}
\newtheorem{prop}[thm]{Proposition}
\newtheorem{cor}[thm]{Corollary}
\newtheorem{conj}[thm]{Conjecture}

\newtheorem{thmA}{Theorem}
\renewcommand{\thethmA}{A}
\theoremstyle{definition}
\newtheorem{rem}[thm]{Remark}
\newtheorem{eg}[thm]{Example}
\numberwithin{equation}{section}
\newcommand{\C}{\mathbb C}
\newcommand{\R}{\mathbb R}
\newcommand{\Z}{\mathbb Z}
\newcommand{\PP}{\mathbb P}
\newcommand{\dd}{\,\mathrm d}
\newcommand{\dA}{\,\mathrm dA}
\newcommand{\loc}{\mathrm{loc}}
\newcommand{\NN}{\mathcal N}
\newcommand{\TN}{T_{\mathrm N}}
\newcommand{\mean}[2]{\mathcal M_{#1}\!\left[#2\right]}
\newcommand{\norm}[1]{\left\lVert#1\right\rVert}
\newcommand{\convmeas}{\xrightarrow{\,\mathrm{meas}\,}}
\newcommand{\av}{\frac{1}{2\pi}\int_{-\pi}^{\pi}}
\DeclareMathOperator{\Wr}{W}
\DeclareMathOperator{\supp}{supp}
\DeclareMathOperator{\sgn}{sgn}
\DeclareMathOperator{\spanop}{span}
\DeclareMathOperator{\Ree}{Re}
\DeclareMathOperator{\ord}{ord}
\DeclareMathOperator{\meas}{area}
\newcommand{\Gn}{\mathcal R_{n+1}}
\newcommand{\fv}{\bm f}
\newcommand{\gv}{\bm g}
\newcommand{\pv}{\bm p}
\newcommand{\hv}{\bm h}
\newcommand{\yv}{\bm y}
\newcommand{\ev}{\bm  e}

\begin{document}
\title[Curves with few inflection points]{Holomorphic curves of finite lower order with few inflection points}
\author[A.~Eremenko]{Alexandre Eremenko}

\address{Mathematics Department, Purdue University,
	West Lafayette, IN 47907, USA}
\email{eremenko@purdue.edu}

\author[T.~Zhang]{Teng Zhang}
\address{School of Mathematics and Statistics, Xi'an Jiaotong University,
Xi'an 710049, P. R. China}
\email{teng.zhang@stu.xjtu.edu.cn}
\date{}
\subjclass[2020]{Primary 30D35; Secondary 30D15, 31A05, 34M05, 14M15}
\keywords{Holomorphic curve, ramification, Wronskian, regular variation,
value distribution}
\begin{abstract}
We prove a conjecture of the first-named author posed in 1998.
Let $f:\C\to\mathbb P^{n}$ be a transcendental linearly non-degenerate
holomorphic curve of finite lower order with small ramification,
that is,
$
N_1(r,f)=o(T(r,f)),
$
where $N_1$ is the counting function of the zeros of the Wronskian and
$T$ is the Cartan--Nevanlinna characteristic.
Then the order and lower order of $f$ coincide, their common value belongs to
$
\left\{1+k/q: k\in\mathbb Z_{\ge0},\ 2\le q\le n+1\right\},
$
and $T(r,f)$ is regularly varying.

\end{abstract}
\maketitle

\tableofcontents

\section{Introduction}\label{sec:intro}

We use the standard terminology and facts of Nevanlinna theory in
one complex variable; see, for example, \cite{Hay64}. The main definitions
will be recalled below for holomorphic curves in projective space.
The lower order and order of a meromorphic function $f$ are defined by
\[
 \lambda:=\liminf_{r\to\infty}\frac{\log T(r,f)}{\log r}
 \leq\limsup_{r\to\infty}\frac{\log T(r,f)}{\log r}=: \rho.
\]
Thus $0\leq\lambda\leq\rho\leq+\infty$ for a nonconstant function.
The simplest prototype of the results considered here is the following
classical elementary fact. If $f$ is a transcendental entire function of
finite lower order such that $f'(z)\ne0$ for every $z\in\C$, then
\[
 f(z)=\int_0^z\exp P(\zeta)\,\mathrm d\zeta+C,
\]
where $P$ is a nonconstant polynomial. In particular,
$
 \rho(f)=\lambda(f)=\deg P\geq1.
$

The linear case corresponds to constant $P$ and has order zero.
 For meromorphic functions, the
corresponding problem is somewhat harder. If $f$ is a transcendental
meromorphic function of finite lower order with no critical points, then
its Schwarzian derivative is entire, and the logarithmic derivative
lemma shows that the Schwarzian derivative is a polynomial. Consequently $f=w_1/w_0$, where
$w_0,w_1$ are linearly independent entire solutions of
\[
 w''+Pw=0
\]
with a polynomial coefficient $P$. In the transcendental case $P\not\equiv0$.
The asymptotic integration theory of such equations \cite{Was65} shows
that the lower order and order coincide and equal
$(\deg P+2)/2\geq1$.
This observation is due to Frithiof Nevanlinna \cite{Nev30}, who conjectured
that finite order together with the weaker condition
$N_1(r)=o(T(r))$ on critical points suffices for the same conclusion about
regularity of growth. When discussing a fixed function, we omit it from
the notation when no confusion can arise.

For the scalar discussion that follows, $f$ is transcendental.
Recall the Second Fundamental Theorem of Nevanlinna:
\begin{equation}\label{1}
 \sum_{j=1}^q m(r,a_j)+N_1(r)\leq(2+o(1))T(r),\qquad r\to\infty.
\end{equation}
Here the $a_j$ are distinct points of $\PP^1=\C\cup\{\infty\}$,
$m(r,a_j)$ are proximity functions, $T(r)$ is the Nevanlinna
characteristic, and
\[
 N_1(r):=N(r,0,f')-N(r,f')+2N(r,f)
\]
is the counting function of critical points. For finite-order functions,
\eqref{1} holds as written; in general one excludes a set of radii of
finite length.

The \emph{deficiency} is
\[
 \delta(a,f)=\liminf_{r\to\infty}\frac{m(r,a)}{T(r)}
 =1-\limsup_{r\to\infty}\frac{N(r,a)}{T(r)}.
\]
The equality follows from the First Main Theorem, and
$0\leq\delta(a,f)\leq1$. For functions of finite order, \eqref{1} gives
the defect relation
\begin{equation}\label{2}
 \sum_{a\in\PP^1}\delta(a,f)
 +\limsup_{r\to\infty}\frac{N_1(r)}{T(r)}\leq2.
\end{equation}
In particular, the maximal-deficiency condition
\begin{equation}\label{sum}
 \sum_a\delta(a,f)=2
\end{equation}
implies F.~Nevanlinna's small-ramification condition
\begin{equation}\label{smal}
 N_1(r)=o(T(r)).
\end{equation}

There is a substantial literature on the consequences of \eqref{sum}
for functions of finite order
\cite{Pfl46,Wei69,Dra81,Ere89a,Ere89b,ES91}; for further discussion, see
\cite{Ere93,Ere02}. The final result was proved by Drasin \cite{Dra87}.

\medskip\noindent
\textbf{Drasin's Theorem.} \emph{Let $f$ be a transcendental meromorphic
function of finite order $\rho$ satisfying \eqref{sum}. Then:
\begin{itemize}
\item[(i)] $2\rho$ is an integer not less than $2$;
\item[(ii)] $\delta(a,f)=p(a)/\rho$, where the $p(a)$ are nonnegative integers.
\end{itemize}}

Drasin's theorem also has a third conclusion that is not needed here.
The first-named author, in collaboration with M.~Sodin, developed an
alternative approach in \cite{Ere89a,Ere89b,ES91}, culminating in the
following theorem from \cite[p.~1196]{Ere93}.

A positive continuous function $\ell:(0,\infty)\to(0,\infty)$ is called
\emph{slowly varying}, in the sense of Karamata, if
$\ell(cr)/\ell(r)\to1$ uniformly for $c\in[1,2]$ as $r\to\infty$.

\begin{thmA}\label{thm:A}
Let $f$ be a transcendental meromorphic function
of finite lower order $\lambda$ satisfying \eqref{smal}. Then:
\begin{itemize}
\item[(a)] $\rho=\lambda=N/2$, where $N\geq2$ is an integer, and
$T(r)\sim r^\rho\ell(r)$ for a slowly varying function $\ell$;
\item[(b)] $\delta(a,f)=p(a)/\lambda$, where the $p(a)$ are nonnegative
integers satisfying
\[
 \sum_{a\in\PP^1}p(a)=2\lambda,
\]
so that \eqref{sum} holds.
\end{itemize}
\end{thmA}

Since conclusion (b) implies the hypothesis of Drasin's theorem,
Theorem~\ref{thm:A} is both a strengthening and a generalization of that theorem.
Its proof is based on potential theory and does not use Drasin's theorem.
One consequence is that, for transcendental meromorphic functions of
finite lower order, \eqref{sum} and \eqref{smal} are equivalent. This
consequence follows from Theorem~\ref{thm:A}, rather than from a direct comparison
of the two conditions. As explained below, the equivalence is specific
to the scalar case. The paper \cite{Ere93} also describes the asymptotic
behavior of these functions in more detail.

We now pass to holomorphic curves. Throughout the rest of the paper,
$n\geq1$ denotes the projective dimension. Thus the target is
$\PP^{n}$, and a reduced representation has \(n+1\) homogeneous
coordinates; the associated differential equations have order \(n+1\).
A \emph{reduced homogeneous representation} of a curve
$f:\C\to\PP^{n}$ is a vector
\[
 \fv=(f_0,\ldots,f_{n}),\qquad f=[\fv],
\]
of entire functions with no common zero. Two reduced representations of
the same curve differ by multiplication by a zero-free entire function.
A curve is \emph{linearly non-degenerate} if its image is not contained in
a hyperplane, equivalently if its coordinate functions are linearly
independent. A curve is \emph{rational} if it extends holomorphically to
$\PP^1$; otherwise it is \emph{transcendental}.

The Cartan--Nevanlinna characteristic is
\[
 T(r,f)=T(r,\fv):=
 \frac1{2\pi}\int_{-\pi}^{\pi}\log\norm{\fv(re^{it})}\,\mathrm dt
 -\log\norm{\fv(0)},
\]
where $\norm{\cdot}$ is the Euclidean norm. This definition does not
depend on the reduced representation. The order and lower order of the
curve are defined from this characteristic as in the scalar case.

The analogue of the critical-point counting function is the counting
function of inflection points. Following the traditional terminology,
we call it the \emph{ramification term}. More precisely, let
\[
 W_{\fv}=W(f_0,\ldots,f_{n})
 :=\det\bigl(f_j^{(i)}\bigr)_{0\leq i,j\le n}
\]
be the Wronskian of a reduced representation. If $n_1(r)$ counts its zeros
in $|z|\leq r$, with multiplicity, then
\[
 \begin{split}
 N_1(r,f)&=\int_0^r\bigl(n_1(t)-n_1(0)\bigr)\frac{\mathrm dt}{t}
                 +n_1(0)\log r\\
 &=\frac1{2\pi}\int_{-\pi}^{\pi}\log|W_{\fv}(re^{it})|\,\mathrm dt+c(\fv).
 \end{split}
\]
This counting function is independent of the reduced representation.
Indeed, if $\widetilde{\fv}=g\fv$ is another reduced representation,
where $g$ is an entire function without zeros, then
\[
W_{\widetilde{\fv}}=g^{n+1}W_{\fv}.
\]
Thus $W_{\widetilde{\fv}}$ and $W_{\fv}$ have the same zero divisor.
Except in the paragraph below explicitly discussing degenerate curves, all
curves in this paper are linearly non-degenerate; thus
$W_{\fv}\not\equiv0$.

Proximity functions and deficiencies are defined for hyperplanes in
$\PP^{n}$; see \cite{Car33,Ru21}. Cartan's generalization of \eqref{2} is
the following form \cite[Theorem~7.1 and Corollary~7.4, pp.~445--446]{GH04};
see also \cite{Car33,Ru21}:
\begin{equation}\label{cartan}
 \sum_{a\in A}\delta(a,f)
 +\limsup_{r\to\infty}\frac{N_1(r,f)}{T(r,f)}\leq n+1
\end{equation}
for a transcendental curve of finite order. The finite-order form used here is
also discussed in \cite{GH04,Ru21}. Here $A$ is any finite set of
hyperplanes in general position: the intersection of any \(n+1\) of them is
empty. Such a system is also called \emph{admissible}. The bound
\eqref{cartan} is the form used in the discussion below. Without the
finite-order assumption, the underlying second main theorem requires
an exceptional set of radii of finite length.

The following regularity part of the conjecture stated in \cite[pp.~4--5]{Ere98}
was later highlighted in the problem note \cite{Ere15b}.

\begin{conj}\label{conj:regularity}
If a linearly non-degenerate holomorphic curve of finite lower order
satisfies \eqref{smal}, then its order and lower order coincide and their
common value is rational.
\end{conj}

Conclusion (b) of Theorem~\ref{thm:A} does not extend to target dimension at least
two, that is, to \(n\geq2\) in the present convention. Examples in \cite{Ere98}
(see also \cite{Ere14,Ere15a}) have finite order and $N_1\equiv0$, but no
admissible system $A$ satisfies
\[
 \sum_{a\in A}\delta(a,f)=n+1.
\]
If linear non-degeneracy is dropped, the defect relation becomes
\[
 \sum_{a\in A}\delta(a,f)\leq2n,
\]
provided the curve is nonconstant and its image is not contained in the
union of the hyperplanes of the admissible system $A$.
Nochka proved this bound \cite{Noc83}; complete proofs are also given in
\cite{Voj97,Ru21}, and a different method appears in \cite{ES92}.
In this setting, \cite{Ere98} describes the curves with maximal
deficiency sum $2n$, in a form analogous to conclusions (a) and (b)
of Theorem~\ref{thm:A}.

We concentrate here on the regularity conclusion (a). Our main result
settles Conjecture~\ref{conj:regularity} in the transcendental case; the
rational case is included in Proposition~\ref{prop:small-order} below.

\begin{thm}\label{thm:main}
Let $f=[\fv]:\C\to\PP^{n}$ be a transcendental linearly non-degenerate
holomorphic curve of finite lower order satisfying
\begin{equation}\label{eq:mainhyp}
 N_1(r,f)=o\bigl(T(r,f)\bigr).
\end{equation}
Then its order and lower order coincide, and their common value $\rho$
belongs to
\begin{equation}\label{ord}
 \mathcal R_{n+1}:=\left\{1+\frac{k}{q}:\ k\in\Z_{\geq0},\
                           q\in\Z,\ 2\leq q\leq n+1\right\}.
\end{equation}
Moreover,
\begin{equation}\label{reg}
 T(r,f)=r^\rho\ell(r),
\end{equation}
where $\ell$ is slowly varying. Equivalently,
$T(cr,f)/T(r,f)\to c^\rho$ locally uniformly for $c\in(0,\infty)$.
\end{thm}


The order set in \eqref{ord} is sharp: for every
\(\rho\in\mathcal R_{n+1}\) there is a transcendental linearly
non-degenerate curve with \(W_{\fv}\equiv1\) and
\[
 c r^\rho\le T(r,\fv)\le C r^\rho
\]
for all sufficiently large \(r\) (Theorem~\ref{thm:realization}).
If one keeps the finite-lower-order and linear-nondegeneracy hypotheses
but drops transcendence, the only additional possibility under
\eqref{eq:mainhyp} is the rational normal curve
\[
 [\fv(z)]=[(1,z,\ldots,z^{n})A],
 \qquad A\in\operatorname{GL}_{n+1}(\mathbb C),
\]
which has order zero (Proposition~\ref{prop:small-order}).
We also have the following corollary.

% We concentrate here on the regularity conclusion (a). Our main result completely resolves Conjecture~\ref{conj:regularity}, including the rational case, and gives a sharp classification of all possible orders.

% \begin{thm}\label{thm:main}
% Let $f=[\fv]:\C\to\PP^{n}$ be a linearly non-degenerate holomorphic curve of finite lower order satisfying
% \begin{equation}\label{eq:mainhyp}
% N_1(r,f)=o\bigl(T(r,f)\bigr).
% \end{equation}
% Then:
% \begin{itemize}
% \item[(i)] The order and lower order of $f$ coincide. If their common value is denoted by $\rho$, then
% $$
%  T(r,f)=r^\rho\ell(r),
% $$
% where $\ell$ is slowly varying; equivalently,
% $
% {T(cr,f)}/{T(r,f)}\longrightarrow c^\rho
% $
% locally uniformly for $c\in(0,\infty)$.
% \item[(ii)] One has
% $$
%  \rho\in\{0\}\cup\mathcal R_{n+1},
%  \qquad
%  \mathcal R_{n+1}=
%  \left\{1+\frac{k}{q}: k\in\Z_{\ge0},\ 2\le q\le n+1\right\}.
% $$
% Moreover, $\rho=0$ if and only if $f$ is rational, in which case
% $$
%  [\fv(z)]=[(1,z,\ldots,z^{n})A],
%  \qquad A\in\operatorname{GL}_{n+1}(\C).
% $$
% Thus every transcendental curve satisfying the hypotheses has
% $\rho\in\mathcal R_{n+1}$.
% \item[(iii)] The classification in \textup{(ii)} is sharp. The value $\rho=0$ is attained by the rational normal curve above. Conversely, for every $\rho\in\mathcal R_{n+1}$, there exists a transcendental linearly non-degenerate holomorphic curve $f=[\fv]$ with $W_{\fv}\equiv1$
% and
% \[
% c r^\rho\leq T(r,f)\leq C r^\rho
% \]
% for all sufficiently large $r$, where $c,C>0$. Hence every value in
% ${0}\cup\mathcal R_{n+1}$ actually occurs under the hypotheses of the theorem, and the set of possible orders cannot be reduced.
% \end{itemize}
% \end{thm}
% The zero-order case admits a quantitative strengthening. Indeed, Theorem~\ref{thm:main} shows that a transcendental linearly non-degenerate curve of order zero cannot satisfy the small-ramification condition \eqref{eq:mainhyp}. In fact, the ramification term must be comparable to the characteristic in the following precise sense.


\begin{cor}\label{cor:zero-ratio}
If $f:\C\to\PP^{n}$ is a transcendental linearly non-degenerate curve
of order zero, then
\[
 \limsup_{r\to\infty}\frac{N_1(r,f)}{T(r,f)}\geq1.
\]
\end{cor}

The proof of Corollary~\ref{cor:zero-ratio} is given in
Proposition~\ref{prop:zero-order-ramification} below.
For earlier scalar results on the frequency of multiple values at small
orders, see \cite{She85}.

The argument of F.~Nevanlinna described above was generalized to
holomorphic curves by V.~P.~Petrenko \cite{Pet84}. We recall it because it
motivates the proof of Theorem~\ref{thm:main}. Suppose
$\fv=(f_0,\ldots,f_{n})$ is reduced and $W_{\fv}$ is zero-free. Expanding
\[
 \begin{vmatrix}
 w&f_0&\cdots&f_{n}\\
 w'&f_0'&\cdots&f_{n}'\\
 \vdots&\vdots&&\vdots\\
 w^{(n+1)}&f_0^{(n+1)}&\cdots&f_{n}^{(n+1)}
 \end{vmatrix}=0
\]
along the first column, the coefficient of $w^{(n+1)}$ is
$(-1)^{{n+1}+2}W_{\fv}$. Division by this coefficient gives a monic equation
\[
 w^{(n+1)}+a_1w^{(n)}+\cdots+a_nw'+a_{n+1}w=0
\]
with entire coefficients and fundamental system $f_0,\ldots,f_{n}$.
For a finite-order curve one can choose a reduced representation whose
coordinates have finite order. Frei's theorem \cite{Fre53} then implies
that the coefficients are polynomials. Passing to
$W_{\fv}^{-1/(n+1)}\fv$ removes the coefficient of the derivative of order
$n$, as proved in Section~2. Asymptotic integration \cite{Was65}
then gives coincident order and lower order, with a value in
$\mathcal R_{n+1}$ in the transcendental case; polynomial solution spaces
also permit the rational normal curve of order zero.

Our proof extends this differential-equation approach to a Wronskian
with infinitely many zeros. Three features are important:
\begin{enumerate}[label=(\arabic*)]
\item
The poles of the differential coefficients must be controlled collectively,
without estimates that deteriorate when poles coalesce.

\item
We rescale at P\'olya peaks and study the resulting distributional limits
of logarithms, following the potential-theoretic approach developed in
\cite{Ere89a,Ere89b,ES91,ES92,Ere93}. A further ingredient is the
convergence in measure of suitably rescaled logarithmic derivatives,
proved in Lemma~\ref{lem:logderivlimit}.

\item
The algebraic part of the argument relies on universal Pl\"ucker
coordinates developed by Karp--Purbhoo. We use the results in \cite{KP26},
together with Purbhoo's fundamental-operator formalism
\cite[Proposition~2.2]{Pur23}. These results are closely related
to the Bethe-algebra correspondence developed by
Mukhin--Tarasov--Varchenko \cite{MTV13}. They are used in
Proposition~\ref{prop:localcompact} to obtain the uniform estimates needed
near the poles of the differential coefficients.
\end{enumerate}


\medskip
\noindent\textbf{Organization of the paper.}
Section~\ref{sec:analytic} collects the analytic preliminaries needed later.
In Section~\ref{sec:algebra}, we establish two estimates for polynomial
solution spaces, which are then used in Section~\ref{sec:compactness} to
derive the local estimates and compactness results. Section~\ref{sec:indices}
is devoted to P\'olya peaks and growth indices, while
Section~\ref{sec:smallorder} treats the case of order less than one.
In Section~\ref{sec:regular}, we prove regular variation and complete the
proof of Theorem~\ref{thm:main}. Finally, Section~\ref{sec:sharpness}
constructs examples showing the sharpness of the theorem and derives the
radial area asymptotics stated in Corollary~\ref{cor:area}.


\medskip
\noindent\textbf{Acknowledgments and AI tools disclosure.}
We thank Mikhail Sodin for helpful comments.

Teng Zhang is supported by the China Scholarship Council, the Young Elite
Scientists Sponsorship Program for PhD Students (China Association for Science
and Technology), and the Fundamental Research Funds for the Central
Universities at Xi'an Jiaotong University (Grant No.~xzy022024045).

ChatGPT was used for English-language editing, proofreading, and grammatical
correction, and as an exploratory tool for discussing possible approaches to
selected results under the authors' mathematical supervision and guidance.
The authors take full responsibility for all mathematical arguments and for
the accuracy and correctness of the final manuscript.


\section{Analytic preliminaries}\label{sec:analytic}

This section collects the analytic tools used in the rescaling argument. We
first normalize the characteristic and the Wronskian, and then introduce the
differential operator attached to a local frame. The remaining lemmas control
logarithmic derivatives and the compactness of normalized logarithms. Each
statement is given in the form in which it will be used later.

Throughout, we write
\[
D_R=\{z\in\C:|z|<R\},\qquad
D(a,R)=\{z\in\C:|z-a|<R\},
\]
and \(\mathrm dA\) for planar area measure. We also use
\[
\Delta=\partial_x^2+\partial_y^2,\qquad
M(r,h)=\max_{|z|\le r}|h(z)|,
\]
for a holomorphic function \(h\) on a neighborhood of \(\overline D_r\), and
\[
\log^+x=\max\{\log x,0\},\qquad
\log^-x=\max\{-\log x,0\}\qquad(x>0).
\]
For a function integrable on \(|z|=r\), set
\[
\mean r u:=\frac1{2\pi}\int_{-\pi}^{\pi}u(re^{i\theta})\,\mathrm d\theta.
\]
Convergence in measure always means convergence with respect to area measure,
locally in the indicated domain. Functions appearing in such statements are
identified a.e.; in particular, values at isolated poles are irrelevant.


\subsection{Jensen's formula and basic characteristic estimates}

Let \(H\not\equiv0\) be entire, and write
\[
H(z)=cz^{m_0}+O(z^{m_0+1}),\qquad c\ne0.
\]
Let $\ord_a(H)$ denote the multiplicity of the zero of $H$ at $a$.
We define the associated integrated zero-counting function by
\begin{equation}\label{eq:zero-counting}
 \NN_H(r):=m_0\log r+
 \sum_{0<|a|<r}\ord_a(H)\log\frac r{|a|}.
\end{equation}
Jensen's formula gives
\begin{equation}\label{eq:zerojensen}
 \mean r{\log|H|}=\NN_H(r)+\log|c|.
\end{equation}
Zeros on \(|z|=r\) contribute zero to the sum in
\eqref{eq:zero-counting}. Hence the small-ramification condition
\eqref{eq:mainhyp} is equivalent to
\[
\NN_{W_{\fv}}(r)=o\bigl(T(r,\fv)\bigr).
\]

If \(n_H(r)\) denotes the number of zeros of \(H\) in \(D_r\), counted with
multiplicity, then for all sufficiently large \(r\),
\begin{equation}\label{eq:countbound}
 n_H(r)\log2\leq \NN_H(2r).
\end{equation}
Indeed, every nonzero zero in \(D_r\) contributes at least \(\log2\) to
\(\NN_H(2r)\), while a zero of multiplicity \(m_0\) at the origin contributes
\(m_0\log(2r)\ge m_0\log2\).

A constant unitary change of homogeneous coordinates leaves both \(T\) and
\(|W_{\fv}|\) unchanged. We therefore make, once and for all, a unitary change
of coordinates such that
\[
\frac{\fv(0)}{\norm{\fv(0)}}={n+1}^{-1/2}(1,\ldots,1).
\]
In particular, $
f_j(0)\ne0
$ for all $0\le j\le n$.

For a meromorphic function \(h\), let
\[
\TN(r,h)=m(r,h)+N(r,h)
\]
be the usual Nevanlinna characteristic, where \(N\) counts poles and
\(m(r,h)=\mean r{\log^+|h|}\). If \(h=f_\ell/f_j\), then
\[
\log^+\left|\frac{f_\ell}{f_j}\right|
 \leq \log\norm{\fv}-\log|f_j|.
\]
Using \eqref{eq:zerojensen} and
\(N(r,f_\ell/f_j)\leq\NN_{f_j}(r)\), we obtain
\begin{equation}\label{eq:quotientbound}
 \TN\!\left(r,\frac{f_\ell}{f_j}\right)
 \leq T(r,\fv)+\log\norm{\fv(0)}-\log|f_j(0)|.
\end{equation}
We shall use this estimate repeatedly below.


\subsection{The fundamental operator and the canonical gauge}


We put \(D=\mathrm d/\mathrm dz\). We first give the elementary local form of the fundamental-operator construction that will be used later. For rational functions, the corresponding construction is discussed in
\cite[Section~2, Proposition~2.2]{Pur23}.

\begin{lem}\label{lem:fundamental-operator}
Let \(U\subset\C\) be a domain and let
\(\gv=(g_0,\ldots,g_{n})\) be holomorphic on \(U\), with
\(W(\gv)\not\equiv0\). Set
$
Z=\{z\in U:W(\gv)(z)=0\}.
$
Then:
\begin{enumerate}[label=\textup{(\roman*)}]
\item On \(U\setminus Z\), there is a unique monic differential operator
\begin{equation}\label{eq:fundamental-operator}
L_{\gv}=D^{n+1}+\sum_{q=1}^{n+1}b_qD^{n+1-q}
\end{equation}
such that
$
L_{\gv}g_j=0, 0\le j\le n.
$ We call \(L_{\gv}\) the \emph{fundamental operator associated with
\(\gv\)}. 
\item Each coefficient \(b_q\) in \eqref{eq:fundamental-operator} extends meromorphically to \(U\), with possible
poles only at points of \(Z\).

\item The coefficient of \(D^{n}\)  in \eqref{eq:fundamental-operator}  is
$$
b_1=-\frac{W(\gv)'}{W(\gv)}.
$$
\end{enumerate}
\end{lem}


\begin{proof}
Since \(W(\gv)\neq0\) on \(U\setminus Z\), the equations
\(L_{\gv}g_j=0\), \(0\le j\le n\), form a nonsingular linear system for
\(b_1,\ldots,b_{n+1}\). Cramer's rule therefore gives a unique solution, and
each \(b_q\) is a quotient of holomorphic functions with denominator
\(W(\gv)\). Hence the coefficients extend meromorphically to \(U\), with
possible poles only at the zeros of the Wronskian.

It remains to identify \(b_1\). Write
$$
W(\gv)=
\det
\begin{pmatrix}
g_0 & \cdots & g_{n}\\
g_0' & \cdots & g_{n}'\\
\vdots & & \vdots\\
g_0^{(n)} & \cdots & g_{n}^{(n)}
\end{pmatrix}.
$$
Differentiating the determinant row by row, for \(0\le k\le n-1\) the
derivative of the \(k\)-th row coincides with the next row, so the
corresponding determinant vanishes. Differentiating the last row gives
$$
W(\gv)'
=
\det
\begin{pmatrix}
g_0 & \cdots & g_{n}\\
g_0' & \cdots & g_{n}'\\
\vdots & & \vdots\\
g_0^{(n-1)} & \cdots & g_{n}^{(n-1)}\\
g_0^{(n+1)} & \cdots & g_{n}^{(n+1)}
\end{pmatrix}.
$$
Since \(L_{\gv}g_j=0\),
$$
g_j^{(n+1)}
=
-b_1g_j^{(n)}
-b_2g_j^{(n-1)}
-\cdots
-b_{n+1}g_j.
$$
By linearity of the determinant in its last row, all terms involving
\(b_2,\ldots,b_{n+1}\) vanish, since each of them makes the last row equal to
one of the preceding rows. Hence
$$
W(\gv)'=-b_1W(\gv),
$$
and therefore
$$
b_1=-\frac{W(\gv)'}{W(\gv)}.
$$
\end{proof}


Recall that \(b_1\) is the coefficient of \(D^{n}\) in
Lemma~\ref{lem:fundamental-operator}. We now remove this term by a scalar gauge
transformation.


\begin{lem}\label{lem:canonical-gauge}
Let \(\gv\) be as in Lemma~\ref{lem:fundamental-operator}, and let \(U_0\)
be simply connected with \(W(\gv)\ne0\) on \(U_0\). Choose a branch
\begin{equation}\label{eq:gauge-factor}
\eta=W(\gv)^{-1/(n+1)}
\end{equation}
on \(U_0\), and put \(\widehat{\gv}=\eta\gv\). Then
\begin{equation}\label{eq:wronskian-gauge}
W(\eta\gv)=\eta^{n+1}W(\gv),\qquad
W(\widehat{\gv})\equiv1.
\end{equation}
The fundamental operator of \(\widehat{\gv}\) has the form
\begin{equation}\label{eq:canonical}
D^{n+1}+\sum_{q=2}^{n+1}Q_q(z)D^{n+1-q}.
\end{equation}
In the application to \(\fv\), the coefficients \(Q_q\) are single-valued
meromorphic functions on \(\C\), holomorphic away from the zeros of
\(W_{\fv}\). Under the dilation \(z\mapsto tz\), they transform according to
\begin{equation}\label{eq:scalecoeff}
Q_q(z)\longmapsto t^qQ_q(tz),\qquad 2\le q\le n+1.
\end{equation}
\end{lem}

\begin{proof}
By Leibniz's rule, the Wronskian matrix of \(\eta\gv\) is obtained from that
of \(\gv\) by left multiplication by a lower triangular matrix whose
diagonal entries are all \(\eta\). This proves the first identity in
\eqref{eq:wronskian-gauge}. The second follows immediately from
\eqref{eq:gauge-factor}.

By Lemma~\ref{lem:fundamental-operator}~(iii), the coefficient of \(D^{n}\) in the fundamental
operator of \(\widehat{\gv}\) is
\[
\widehat b_1
=b_1-{n+1}\frac{\eta'}{\eta}
=-\frac{W(\gv)'}{W(\gv)}
-{n+1}\left(-\frac1{n+1}\frac{W(\gv)'}{W(\gv)}\right)
=0.
\]
Hence the normalized operator has the form \eqref{eq:canonical}.

Different choices of the branch in \eqref{eq:gauge-factor} differ by a
constant \({n+1}\)th root of unity. Multiplying all solutions by such a constant
does not change their fundamental operator. The locally defined coefficients
in \eqref{eq:canonical} therefore agree on overlaps. In the application to
\(\fv\), they define single-valued meromorphic functions on \(\C\), with
possible poles only at the zeros of \(W_{\fv}\).

Finally, let \(v\) be a solution of \eqref{eq:canonical} and set
\(u(z)=v(tz)\). Since
\[
 u^{(k)}(z)=t^k v^{(k)}(tz),
\]
substitution into \eqref{eq:canonical} shows that the coefficient of
\(D^{n+1-q}u\) is \(t^qQ_q(tz)\). This proves \eqref{eq:scalecoeff}.
\end{proof}



\subsection{Logarithmic derivatives at a prescribed radius}

The usual logarithmic derivative lemma is often stated outside an exceptional
set of radii. Later we need an estimate at a prescribed radius, so we use the
following fixed-radius form. It follows directly from the Poisson--Jensen
formula; compare \cite[Chapter~3]{Hay64}.

\begin{lem}\label{lem:logderivbound}
For every integer \(k\ge1\) there is a constant \(C_k\) such that, whenever
\(h\) is meromorphic in a neighborhood of \(\overline D_4\) and
\(0<|h(0)|<\infty\),
\begin{equation}\label{eq:logderivbound}
 m\!\left(1,\frac{h^{(k)}}h\right)
 \leq C_k\log\bigl(2+\TN(4,h)+|\log|h(0)||\bigr).
\end{equation}
\end{lem}

\begin{proof}
Set
$
A:=2+\TN(4,h)+|\log|h(0)||.
$

\medskip\noindent
\textbf{\underline{Step 1. Counting zeros and poles:}}
Since \(m(r,h)\ge0\), we have
$$
N(4,h)\le \TN(4,h),
\qquad
N(4,1/h)\le \TN(4,1/h).
$$
Together with
$
\TN(4,1/h)=\TN(4,h)-\log|h(0)|,
$
this gives
\begin{equation}\label{eq:counting-bound}
\begin{aligned}
N(4,h)+N(4,1/h)
&\le \TN(4,h)+\TN(4,1/h)\\
&=2\TN(4,h)-\log|h(0)|\\
&\le 2\TN(4,h)+|\log|h(0)||\\
&\le 2A.
\end{aligned}
\end{equation}
Every zero or pole \(a\) of \(h\) in \(D_3\), counted with multiplicity \(m\), contributes at least
$$
m\log\frac{4}{|a|}
\ge m\log\frac43
$$
to \(N(4,1/h)\) or \(N(4,h)\), respectively. Hence, if \(M\) denotes the total number of zeros and poles of \(h\) in \(D_3\), counted with multiplicity, then
\begin{equation}\label{eq}
M\log\frac43
\le N(4,h)+N(4,1/h)
\le 2A.
\end{equation}


\medskip\noindent
\textbf{\underline{Step 2. The Poisson--Jensen representation:}}
Choose \(R\in(2,3)\) so that \(|z|=R\) contains no zero or pole of \(h\).
Since
$
|\log|h||=\log^+|h|+\log^-|h|
$
and
$
\log^-|h|=\log^+|1/h|,
$
the monotonicity of the Nevanlinna characteristic and
\eqref{eq:counting-bound} give
\begin{equation}\label{eq:boundary-mean}
\mean R{|\log|h||}
=m(R,h)+m(R,1/h)
\leq \TN(4,h)+\TN(4,1/h)
\leq 2A.
\end{equation}

Put \(L=h'/h\), and let \(\nu_h(a)=\ord_a(h)\), with the convention that
\(\nu_h(a)>0\) at a zero and \(\nu_h(a)<0\) at a pole. Since \(h\) has no
zero or pole on \(|z|=R\), the Poisson--Jensen formula in \(D_R\) gives
\begin{equation}\label{eq:poisson-jensen}
\log |h(z)|=H_R(z)
+
\sum_{a\in D_R}\nu_h(a)
\log\left|
\frac{R(z-a)}{R^2-\overline a z}
\right|,
\end{equation}
where
$
H_R(z)=\frac1{2\pi}\int_{-\pi}^{\pi}
\frac{R^2-|z|^2}{|Re^{it}-z|^2}
\log|h(Re^{it})|\,\mathrm dt.
$
Set
$
\partial=\frac12(\partial_x-i\partial_y).
$
Away from the zeros and poles of \(h\),
$
2\partial\log|h|=h'/h=L.
$
Moreover,
$$
2\partial\log|z-a|=\frac1{z-a},
\qquad
2\partial\log|R^2-\overline a z|
=-\frac{\overline a}{R^2-\overline a z}.
$$
Applying \(2\partial\) to \eqref{eq:poisson-jensen} therefore gives
$$
L(z)
=
2\partial H_R(z)
+
\sum_{a\in D_R}\nu_h(a)
\left(
\frac1{z-a}
+
\frac{\overline a}{R^2-\overline a z}
\right).
$$
Differentiating this identity \(j\) times with respect to \(z\), for
\(j\geq0\), we obtain
\begin{equation}\label{eq:Lj-representation}
L^{(j)}(z)=
2\partial^{j+1}H_R(z)
+
j!\sum_{a\in D_R}\nu_h(a)
\left(
\frac{(-1)^j}{(z-a)^{j+1}}
+
\frac{\overline a^{j+1}}
{(R^2-\overline a z)^{j+1}}
\right).
\end{equation}

\medskip\noindent
\textbf{\underline{Step 3. Pointwise bounds on the unit circle:}}
We first estimate the harmonic term in \eqref{eq:Lj-representation}. Write
\begin{equation}\label{eq:poisson-kernel}
P_R(z,e^{it})
= \frac{R^2-|z|^2}{|Re^{it}-z|^2}
= \operatorname{Re}\frac{Re^{it}+z}{Re^{it}-z}.
\end{equation}
Differentiating the second identity in \eqref{eq:poisson-kernel}, for every
integer \(m\ge1\),
$$
2\partial^mP_R(z,e^{it})
=
\frac{2m!Re^{it}}{(Re^{it}-z)^{m+1}}.
$$
Since \(2<R<3\) and \(|z|=1\),
$$
|Re^{it}-z|\ge R-1>1,
$$
and hence
\begin{equation}\label{eq:poisson-kernel-derivative-bound}
\left|2\partial^mP_R(z,e^{it})\right|
\le
\frac{2m!R}{(R-1)^{m+1}}
\le 6m!.
\end{equation}
Differentiating under the integral defining \(H_R\), and using the triangle
inequality for the integral, we obtain, for \(|z|=1\),
\begin{equation}\label{eq:poisson-harmonic-estimate}
\begin{aligned}
\left|2\partial^{j+1}H_R(z)\right|
&=
\left|
\frac1{2\pi}\int_{-\pi}^{\pi}
2\partial^{j+1}P_R(z,e^{it})
\log|h(Re^{it})|\,\mathrm dt
\right|\\
&\le
\frac1{2\pi}\int_{-\pi}^{\pi}
\left|2\partial^{j+1}P_R(z,e^{it})\right|
|\log|h(Re^{it})||\,\mathrm dt\\
&\le
6(j+1)!
\frac1{2\pi}\int_{-\pi}^{\pi}
|\log|h(Re^{it})||\,\mathrm dt
\qquad\text{(by \eqref{eq:poisson-kernel-derivative-bound})}\\
&=
6(j+1)!\mean R{|\log|h||}\\
&\le
12(j+1)!A
\qquad\text{(by \eqref{eq:boundary-mean})}.
\end{aligned}
\end{equation}
Set
$
C_{1,j}:=12(j+1)!.
$

We next estimate the reflected terms in \eqref{eq:Lj-representation}.
For \(a\in D_R\) and \(|z|=1\),
$$
|R^2-\overline a z|
\ge R^2-|a|
>R(R-1),
$$
and therefore
\begin{equation}\label{eq:reflected-single-estimate}
\left|
\frac{\overline a^{j+1}}
{(R^2-\overline a z)^{j+1}}
\right|
\le
\left(\frac{|a|}{R^2-|a|}\right)^{j+1}
<
\frac1{(R-1)^{j+1}}
\le1.
\end{equation}
Since \(R<3\), all zeros and poles occurring in
\eqref{eq:Lj-representation} lie in \(D_3\). By \eqref{eq}, their total multiplicity satisfies
$$
M\le\frac{2A}{\log(4/3)}.
$$
By \eqref{eq:reflected-single-estimate}, the total contribution of the
reflected terms in \eqref{eq:Lj-representation} is bounded by
\begin{equation}\label{eq:reflected-total-estimate}
j!\sum_{a\in D_R}|\nu_h(a)|
\left|
\frac{\overline a^{j+1}}
{(R^2-\overline a z)^{j+1}}
\right|
\le
\frac{2j!}{\log(4/3)}A.
\end{equation}
Set
$
C_{2,j}:=\frac{2j!}{\log(4/3)}.
$

Combining \eqref{eq:Lj-representation},
\eqref{eq:poisson-harmonic-estimate}, and
\eqref{eq:reflected-total-estimate}, and listing the zeros and poles according to multiplicity, we obtain, for every \(z\) which is neither a zero nor a pole of \(h\),
\begin{equation}\label{eq:Lj-pointwise}
|L^{(j)}(z)|
\le
(C_{1,j}+C_{2,j})A+
j!\sum_a |z-a|^{-(j+1)},
\qquad |z|=1.
\end{equation}
Here and below, zeros and poles are
listed with multiplicity.

\medskip\noindent
\textbf{\underline{Step 4. The proximity estimate for \(L^{(j)}\):}}
Choose
$
0<\alpha<1/{(j+1)},
$
depending only on \(j\), and put
$
\beta=(j+1)\alpha<1.
$
We first note that
$
\int_{-\pi}^{\pi}|e^{i\theta}-a|^{-\beta}\,\mathrm d\theta
$
is bounded uniformly in \(a\in\C\). Indeed, if
\(|a|\le1/2\) or \(|a|\ge3/2\), then the distance from \(a\) to the unit
circle is at least \(1/2\). If \(1/2<|a|<3/2\), a rotation allows us to
write \(a=r>0\), and
$$
|e^{i\theta}-r|^2
=
(1-r)^2+2r(1-\cos\theta)
\ge
\frac{2}{\pi^2}\theta^2,
\qquad -\pi\le\theta\le\pi.
$$

Since \(\beta<1\), the function \(|\theta|^{-\beta}\) is integrable on
\([-\pi,\pi]\). Thus there is a constant \(C_\beta\), independent of
\(a\), such that
\begin{equation}\label{eq:circle-kernel-bound}
\int_{-\pi}^{\pi}|e^{i\theta}-a|^{-\beta}\,\mathrm d\theta
\le C_\beta.
\end{equation}
Since \(0<\alpha<1\), we have
$
(x+y)^\alpha\le x^\alpha+y^\alpha
$
for \(x,y\ge0\), and more generally
$
(\sum_i x_i)^\alpha\le\sum_i x_i^\alpha.
$
Applying these inequalities to \eqref{eq:Lj-pointwise}, we obtain
$$
|L^{(j)}(e^{i\theta})|^\alpha
\le
(C_{1,j}+C_{2,j})^\alpha A^\alpha
+
(j!)^\alpha\sum_a
|e^{i\theta}-a|^{-(j+1)\alpha}.
$$
Integrating over the unit circle and using \eqref{eq:circle-kernel-bound}, we obtain
$$
\begin{aligned}
\frac1{2\pi}\int_{-\pi}^{\pi}
|L^{(j)}(e^{i\theta})|^\alpha\,\mathrm d\theta
&\le
(C_{1,j}+C_{2,j})^\alpha A^\alpha
+
\frac{(j!)^\alpha}{2\pi}
\sum_a
\int_{-\pi}^{\pi}
|e^{i\theta}-a|^{-\beta}\,\mathrm d\theta\\
&\le
(C_{1,j}+C_{2,j})^\alpha A^\alpha
+
\frac{(j!)^\alpha C_\beta}{2\pi}M\\
&\le
\left(
(C_{1,j}+C_{2,j})^\alpha
+
\frac{(j!)^\alpha C_\beta}{\pi\log(4/3)}
\right)A,
\end{aligned}
$$
where in the last inequality we used \(A^\alpha\le A\), because \(A\ge2\), and
\eqref{eq}. Set
$$
C_{3,j}:=
(C_{1,j}+C_{2,j})^\alpha
+
\frac{(j!)^\alpha C_\beta}{\pi\log(4/3)}.
$$
Thus
\begin{equation}\label{eq:Lj-moment}
\frac1{2\pi}\int_{-\pi}^{\pi}
|L^{(j)}(e^{i\theta})|^\alpha\,\mathrm d\theta
\le C_{3,j}A.
\end{equation}
Finally,
$$
\log^+x
\le
\frac1\alpha\log(1+x^\alpha),
\qquad x\ge0.
$$
Therefore, by Jensen's inequality for the concave function \(\log x\) and
\eqref{eq:Lj-moment},
$$
\begin{aligned}
m(1,L^{(j)})
&\le
\frac1\alpha\frac1{2\pi}
\int_{-\pi}^{\pi}
\log\left(1+|L^{(j)}(e^{i\theta})|^\alpha\right)
\,\mathrm d\theta\\
&\le
\frac1\alpha
\log\left(
1+
\frac1{2\pi}\int_{-\pi}^{\pi}
|L^{(j)}(e^{i\theta})|^\alpha\,\mathrm d\theta
\right)\\
&\le
\frac1\alpha\log(1+C_{3,j}A)\\
&\le
C_j\log A,
\end{aligned}
$$
where
$
C_j:=
\left(
1+{\log(1+C_{3,j})}/{\log2}
\right)/\alpha.
$
Thus, for every fixed \(j\ge0\),
\begin{equation}\label{eq:Lj-proximity}
m(1,L^{(j)})\le C_j\log A.
\end{equation}

\medskip\noindent
\textbf{\underline{Step 5. Recovering \(h^{(k)}/h\):}}
Put
$
P_k={h^{(k)}}/h.
$
Recall from Step~2 that \(h'=Lh\). Thus \(P_1=L\), and differentiation gives the
recursion
$$
P_{k+1}
=
P_k'+LP_k.
$$
It follows inductively that \(P_k\) is a differential polynomial in
\(L,L',\ldots,L^{(k-1)}\), with coefficients depending only on \(k\).

More precisely, 
\begin{equation}\label{eq:Pk}
    P_k
=
\sum_{\substack{m_0,\ldots,m_{k-1}\ge0\\
\sum_{j=0}^{k-1}(j+1)m_j=k}}
C_{\mathbf m}
\prod_{j=0}^{k-1}
\bigl(L^{(j)}\bigr)^{m_j},
\end{equation}
where
$
C_{\mathbf m}
={k!}/
{\displaystyle\prod_{j=0}^{k-1}m_j!\,((j+1)!)^{m_j}}
$
depends only on the multi-index \(\mathbf m=(m_0,\ldots,m_{k-1})\); for details, see Lemma~\ref{lem:differential-polynomial} in Appendix~\ref{app:differential-polynomial}.
We also need the proximity estimate that follows from
\eqref{eq:Pk}. For each monomial occurring in the sum,
$$
m\!\left(
1,
C_{\mathbf m}
\prod_{j=0}^{k-1}
\bigl(L^{(j)}\bigr)^{m_j}
\right)
\le
\sum_{j=0}^{k-1}
m_j\,m(1,L^{(j)})
+O_k(1).
$$
Indeed,
$$
\log^+\left|
C_{\mathbf m}
\prod_{j=0}^{k-1}
\bigl(L^{(j)}\bigr)^{m_j}
\right|
\le
\log^+|C_{\mathbf m}|
+
\sum_{j=0}^{k-1}
m_j\log^+|L^{(j)}|,
$$
and integration over the unit circle gives the asserted estimate.
For fixed \(k\), only finitely many multi-indices
\(\mathbf m\) satisfy the constraint in \eqref{eq:Pk}.
Applying the standard proximity estimate for finite sums therefore gives
\begin{equation}\label{eq:differential-polynomial-proximity}
m\!\left(1,\frac{h^{(k)}}{h}\right)
\le
C_k\left(
1+\sum_{j=0}^{k-1}m(1,L^{(j)})
\right),
\end{equation}
where \(C_k\) depends only on \(k\).
Combining \eqref{eq:differential-polynomial-proximity} with
\eqref{eq:Lj-proximity} for \(0\le j\le k-1\), and using the definition of
\(A\), we obtain
\[
m\!\left(1,\frac{h^{(k)}}h\right)\le C'_k\log A,
\]
which is \eqref{eq:logderivbound}.
\end{proof}






\subsection{Limits of logarithmic derivatives}

We also need a limiting form of the logarithmic derivative estimate. After
normalization, the first logarithmic derivative has a nontrivial limit, while
its higher derivatives become negligible.

\begin{lem}\label{lem:logderivlimit}
Let \(\Omega\subset\C\) be a domain, let \(s_\nu\to\infty\), and let
\(h_\nu\not\equiv0\) be holomorphic on \(\Omega\). Suppose that
\begin{equation}\label{eq:log-limit-assumption}
 u_\nu:=s_\nu^{-1}\log|h_\nu|
 \longrightarrow u
 \qquad\text{in }L^1_{\loc}(\Omega),
\end{equation}
where \(u\) is subharmonic and not identically \(-\infty\). Then
\begin{equation}\label{eq:sobolev-limit}
 u\in W^{1,a}_{\loc}(\Omega)\qquad(1\leq a<2).
\end{equation}
With
$
 \partial=\tfrac12(\partial_x-i\partial_y),
$
and with \(\partial u\) represented by its a.e.-defined Sobolev function,
one has, for every fixed integer \(k\ge1\),
\begin{equation}\label{eq:logderivlimit}
 \frac{h_\nu^{(k)}}{s_\nu^kh_\nu}
 \convmeas (2\partial u)^k
 \qquad\text{locally in }\Omega.
\end{equation}
\end{lem}
\begin{proof}

\medskip\noindent
\textbf{\underline{Step 1. The Riesz measures and localization:}}
Set $
\mu_\nu:=\Delta u_\nu/(2\pi).$
Since
$
u_\nu=\log|h_\nu|/{s_\nu},
$
the measure \(\mu_\nu\) is the zero-counting measure of \(h_\nu\),
divided by \(s_\nu\). By \eqref{eq:log-limit-assumption},
\(u_\nu\to u\) in \(L^1_{\loc}(\Omega)\), hence also in the sense of
distributions. Therefore
\begin{equation}\label{eq:riesz-measure-limit}
 \mu_\nu\longrightarrow
 \mu:=\frac1{2\pi}\Delta u
\end{equation}
weakly on compact subsets of \(\Omega\). In particular, for every
\(K_1\Subset\Omega\),
\begin{equation}\label{eq:riesz-mass-bound}
 \sup_\nu \mu_\nu(K_1)<\infty.
\end{equation}

Fix \(K\Subset\Omega\). Choose
\(\chi\in C_c^\infty(\Omega)\) such that
\[
0\leq\chi\leq1,
\qquad
\chi\equiv1
\quad\text{on a neighborhood of }K,
\]
and define
\begin{equation}\label{eq:localized-potentials}
 v_\nu(z)
 :=
 \int\chi(a)\log|z-a|\,\mathrm d\mu_\nu(a),
 \qquad
 v(z)
 :=
 \int\chi(a)\log|z-a|\,\mathrm d\mu(a).
\end{equation}
Put
$
H_\nu:=u_\nu-v_\nu$
and $
H:=u-v.
$
Since
\[
\frac1{2\pi}\Delta v_\nu=\chi\mu_\nu,
\qquad
\frac1{2\pi}\Delta v=\chi\mu,
\]
and \(\chi=1\) near \(K\), the functions \(H_\nu\) and \(H\) are
harmonic on a fixed neighborhood of \(K\).

We shall use the decomposition
\begin{equation}\label{eq:potential-harmonic-decomposition}
 u_\nu=v_\nu+H_\nu,
 \qquad
 u=v+H
\end{equation}
in the proof below.

\medskip\noindent
\textbf{\underline{Step 2. The first logarithmic derivative:}}
We first prove convergence of the derivatives of the logarithmic
potentials. Let \(1\leq p<2\). Since
\[
2\partial\log|z-a|=\frac1{z-a},
\]
we have formally
\[
2\partial v_\nu(z)
 =
 \int\frac{\chi(a)}{z-a}\,\mathrm d\mu_\nu(a).
\]
To justify passage to the limit, truncate the kernel
\((z-a)^{-1}\) smoothly at distance \(\delta>0\) from the diagonal.
For fixed \(\delta\), the truncated kernel is bounded and continuous
on compact sets, so \eqref{eq:riesz-measure-limit} gives convergence
of the corresponding truncated integrals.

For the part removed by the truncation, Minkowski's inequality and
\eqref{eq:riesz-mass-bound} give
\begin{equation}\label{eq:kernel-truncation-bound}
 \left\|
 \int_{|z-a|<2\delta}
 \frac{\chi(a)}{z-a}\,\mathrm d\mu_\nu(a)
 \right\|_{L^p(K)}
 \leq C_{K,p}\delta^{\,2/p-1},
\end{equation}
uniformly in \(\nu\). The same estimate
\eqref{eq:kernel-truncation-bound} holds with \(\mu\) in place of
\(\mu_\nu\). Since \(2/p-1>0\), letting
\(\delta\downarrow0\) gives
\begin{equation}\label{eq:potential-derivative-limit}
 \partial v_\nu\longrightarrow\partial v
 \qquad\text{in }L^p(K),
 \qquad 1\leq p<2.
\end{equation}
As \(K\Subset\Omega\) was arbitrary, this is local convergence in
\(\Omega\).

The same truncation argument, now applied to the kernel
\(\log|z-a|\), gives
\[
v_\nu\longrightarrow v
\qquad\text{in }L^1_{\loc}(\Omega).
\]
Together with \eqref{eq:log-limit-assumption}, this yields
\begin{equation}\label{eq:harmonic-part-limit}
 H_\nu\longrightarrow H
 \qquad\text{in }L^1_{\loc}
\end{equation}
on a neighborhood of \(K\). Since all these functions are harmonic
there, the usual interior estimates imply convergence of all their
derivatives on smaller compact subsets. In particular,
\begin{equation}\label{eq:harmonic-derivative-limit}
 \partial H_\nu\longrightarrow\partial H
 \qquad\text{locally uniformly near }K.
\end{equation}

Away from the zeros of \(h_\nu\),
\[
2\partial u_\nu
 =
 \frac{h_\nu'}{s_\nu h_\nu}.
\]
The zeros form a set of area zero, so this identity holds almost
everywhere. Combining
\eqref{eq:potential-harmonic-decomposition},
\eqref{eq:potential-derivative-limit}, and
\eqref{eq:harmonic-derivative-limit}, we obtain
\begin{equation}\label{eq:first-logderiv-limit}
 \frac{h_\nu'}{s_\nu h_\nu}
 =
 2\partial u_\nu
 \longrightarrow
 2\partial u
 \qquad\text{in }L^p_{\loc}(\Omega),
 \qquad 1\leq p<2.
\end{equation}
Since \(u\) is real-valued, control of \(\partial u\) gives control of
the full weak gradient. Thus
\[
u\in W^{1,p}_{\loc}(\Omega),
\qquad
1\leq p<2,
\]
which proves \eqref{eq:sobolev-limit}.

\medskip\noindent
\textbf{\underline{Step 3. Higher derivatives of the logarithmic derivative:}}
Put $
L_\nu:={h_\nu'}/{h_\nu}.$
We claim the convergence in \eqref{eq:higher-logderiv-goal} for every
fixed \(j\geq1\):
\begin{equation}\label{eq:higher-logderiv-goal}
 \frac{L_\nu^{(j)}}{s_\nu^{j+1}}
 \convmeas0
 \qquad\text{locally in }\Omega.
\end{equation}

From \eqref{eq:localized-potentials} and the fact that
\(\mu_\nu\) is the zero-counting measure divided by \(s_\nu\), we have
near \(K\)
\[
2s_\nu\partial v_\nu(z)
 =
 \sum_a
 \frac{\chi(a)\ord_a(h_\nu)}{z-a}.
\]
Using
$
L_\nu=2s_\nu\partial u_\nu
     =2s_\nu\partial v_\nu+2s_\nu\partial H_\nu,
$
and differentiating \(j\) times, we obtain
\begin{equation}\label{eq:singular-logderivative}
 L_\nu^{(j)}(z)
 =
 (-1)^j j!
 \sum_a
 \frac{\chi(a)\ord_a(h_\nu)}
 {(z-a)^{j+1}}
 +
 R_{\nu,j}(z),
\end{equation}
where
\[
R_{\nu,j}(z)
 :=
 2s_\nu\partial^{j+1}H_\nu(z).
\]
By \eqref{eq:harmonic-part-limit} and the interior estimates for
harmonic functions,
\begin{equation}\label{eq:harmonic-remainder-bound}
 \sup_K|R_{\nu,j}|
 \leq C_{K,j}s_\nu.
\end{equation}

On the other hand, \eqref{eq:riesz-mass-bound} gives
\begin{equation}\label{eq:local-zero-number}
 \sum_{a\in\supp\chi}\ord_a(h_\nu)
 =
 s_\nu\mu_\nu(\supp\chi)
 =
 O(s_\nu).
\end{equation}

Choose
\begin{equation}\label{eq:alpha-choice}
 \frac1{j+1}
 <
 \alpha_j
 <
 \min\left\{1,\frac2{j+1}\right\}.
\end{equation}
Then
$
(j+1)\alpha_j<2,
$
so
$
|z-a|^{-(j+1)\alpha_j}
$
is locally integrable with respect to planar area measure. Moreover,
\(\alpha_j<1\), and therefore
\[
\left(\sum_\ell x_\ell\right)^{\alpha_j}
\leq
\sum_\ell x_\ell^{\alpha_j},
\qquad x_\ell\geq0.
\]
Listing the zeros with multiplicity, and using
\eqref{eq:singular-logderivative},
\eqref{eq:harmonic-remainder-bound}, and
\eqref{eq:local-zero-number}, we obtain
\begin{equation}\label{eq:higher-logderiv-Lalpha}
\begin{aligned}
 \int_K
 \left|
 \frac{L_\nu^{(j)}}{s_\nu^{j+1}}
 \right|^{\alpha_j}\mathrm dA
 &\leq
 C_{K,j}
 \frac1{s_\nu^{(j+1)\alpha_j}}
 \sum_a
 \int_K
 \frac{\mathrm dA(z)}
 {|z-a|^{(j+1)\alpha_j}}
 +C_{K,j}s_\nu^{-j\alpha_j}
 \\
 &\leq
 C_{K,j}s_\nu^{\,1-(j+1)\alpha_j}
 +C_{K,j}s_\nu^{-j\alpha_j}.
\end{aligned}
\end{equation}
By \eqref{eq:alpha-choice},
\[
1-(j+1)\alpha_j<0,
\]
and both terms on the right-hand side of
\eqref{eq:higher-logderiv-Lalpha} tend to zero. Consequently,
\[
\frac{L_\nu^{(j)}}{s_\nu^{j+1}}
\longrightarrow0
\qquad\text{in }L^{\alpha_j}_{\loc},
\]
and hence
\begin{equation}\label{eq:higher-logderiv-measure}
 \frac{L_\nu^{(j)}}{s_\nu^{j+1}}
 \convmeas0
 \qquad(j\geq1)
\end{equation}
locally in \(\Omega\).

\medskip\noindent
\textbf{\underline{Step 4. The differential-polynomial expansion:}}
As in Step 5 of the proof of Lemma~\ref{lem:logderivbound}
(see \eqref{eq:Pk}), repeated differentiation of
$
h_\nu'=L_\nu h_\nu
$
gives, for every \(k\geq1\),
\begin{equation}\label{eq:differential-polynomial}
 \frac{h_\nu^{(k)}}{h_\nu}
 =
 \sum_{\substack{m_0,\ldots,m_{k-1}\ge0\\
                  \sum_{j=0}^{k-1}(j+1)m_j=k}}
 C_{\mathbf m}
 \prod_{j=0}^{k-1}
 \bigl(L_\nu^{(j)}\bigr)^{m_j},
\end{equation}
where the constants \(C_{\mathbf m}\) depend only on the multi-index
$
\mathbf m=(m_0,\ldots,m_{k-1}).
$
The weight condition in \eqref{eq:differential-polynomial},
\(\sum_{j=0}^{k-1}(j+1)m_j=k\), gives
\begin{equation}\label{eq:normalized-monomial}
 \frac1{s_\nu^k}
 \prod_{j=0}^{k-1}
 \bigl(L_\nu^{(j)}\bigr)^{m_j}
 =
 \prod_{j=0}^{k-1}
 \left(
 \frac{L_\nu^{(j)}}{s_\nu^{j+1}}
 \right)^{m_j}.
\end{equation}

The term corresponding to
\[
m_0=k,
\qquad
m_1=\cdots=m_{k-1}=0,
\]
is \(L_\nu^k\). Every other term in
\eqref{eq:differential-polynomial} contains at least one factor
\(L_\nu^{(j)}\) with \(j\geq1\).

\medskip\noindent
\textbf{\underline{Step 5. Passage to the limit:}}
Define
\[
X_{\nu,0}:=\frac{L_\nu}{s_\nu},
\qquad
X_{\nu,j}
 :=
 \frac{L_\nu^{(j)}}{s_\nu^{j+1}}
 \quad(j\geq1).
\]
By \eqref{eq:first-logderiv-limit},
\begin{equation}\label{eq:X0-limit}
 X_{\nu,0}
 \convmeas
 2\partial u,
\end{equation}
while by \eqref{eq:higher-logderiv-measure},
\begin{equation}\label{eq:Xj-limit}
 X_{\nu,j}\convmeas0
 \qquad(j\geq1).
\end{equation}
In particular, every sequence \(X_{\nu,j}\) is bounded in measure.

Consider a monomial on the right-hand side of
\eqref{eq:normalized-monomial}. If it is the distinguished term
\(m_0=k\), then \eqref{eq:X0-limit} gives
\begin{equation}\label{eq:leading-monomial-limit}
 X_{\nu,0}^k
 \convmeas
 (2\partial u)^k.
\end{equation}
For every other monomial there is some \(j_0\geq1\) with
\(m_{j_0}>0\). The corresponding factor
\(X_{\nu,j_0}^{m_{j_0}}\) tends to zero in measure by
\eqref{eq:Xj-limit}, while all the remaining factors are bounded in
measure. Hence the whole monomial tends to zero in measure.

Since the sum in \eqref{eq:differential-polynomial} contains only
finitely many terms, \eqref{eq:normalized-monomial} and
\eqref{eq:leading-monomial-limit} yield
\[
 \frac{h_\nu^{(k)}}{s_\nu^kh_\nu}
 \convmeas
 (2\partial u)^k
\]
locally in \(\Omega\). This is precisely
\eqref{eq:logderivlimit}, and the proof is complete.
\end{proof}


\begin{rem}\label{rem:measurable-products}
All products in \eqref{eq:logderivlimit} and in the limiting differential
equations below are products of measurable functions, not of distributions.
In particular,
\[
(\partial u)^k\in L^p_{\loc}
 \qquad
 \left(0<p<\frac2k\right),
\]
but local \(L^1\)-integrability may fail for \(k\ge2\). For example, if
\(u(z)=\log|z|\), then
\[
 2\partial u=\frac1z
 \qquad\text{a.e.},
\]
whereas \(1/z^2\notin L^1_{\loc}\).

We shall also use the elementary fact that if \(a_\nu\to0\) in measure and
\(b_\nu\) is bounded in measure, then \(a_\nu b_\nu\to0\) in measure.
\end{rem}


\subsection{A Wronskian consequence and subharmonic compactness}

Set
\[
\kappa:=\sum_{i=0}^{n}i=\frac{n(n+1)}2.
\]
Thus \(\kappa\) is the total derivative order in an \({n+1}\times {n+1}\)
Wronskian.

\begin{lem}\label{lem:sum}
Let \(s_\nu\to\infty\), and let
\(\hv_\nu=(h_{\nu,0},\ldots,h_{\nu,n})\) be holomorphic on a domain
\(\Omega\). Suppose that, for every \(j\),
\begin{equation}\label{eq:component-log-limits}
 s_\nu^{-1}\log|h_{\nu,j}|
 \longrightarrow v_j
 \qquad\text{in }L^1_{\loc}(\Omega),
\end{equation}
where \(v_j\) is subharmonic and not identically \(-\infty\), and suppose that
\begin{equation}\label{eq:small-wronskian-local}
 s_\nu^{-1}\log|W(\hv_\nu)|
 \convmeas0
 \qquad\text{locally in }\Omega.
\end{equation}
Then
\begin{equation}\label{eq:sum-positive}
 \sum_{j=0}^{n}v_j\ge0
 \qquad\text{a.e.\ on }\Omega.
\end{equation}
\end{lem}

\begin{proof}

\medskip\noindent
\textbf{\underline{Step 1. Normalize the Wronskian:}}
Away from the zeros of the components \(h_{\nu,j}\), we may factor
\(h_{\nu,j}\) from the \(j\)-th column of the Wronskian. Since the
\(i\)-th row contains derivatives of order \(i\), we obtain
\begin{equation}\label{eq:normalized-wronskian}
 \frac{W(\hv_\nu)}
 {s_\nu^\kappa\prod_{j=0}^{n}h_{\nu,j}}
 =
 \det\left(
 \frac{h_{\nu,j}^{(i)}}
 {s_\nu^ih_{\nu,j}}
 \right)_{0\leq i,j\le n}.
\end{equation}
The identity holds almost everywhere, which is sufficient for what follows,
since the zeros of the components form a set of area zero.

Put
\[
D_\nu:=
\det\left(
 \frac{h_{\nu,j}^{(i)}}
 {s_\nu^ih_{\nu,j}}
 \right)_{0\leq i,j\le n}.
\]

\medskip\noindent
\textbf{\underline{Step 2. Control the normalized determinant:}}
For \(i=0\), the corresponding matrix entry is identically \(1\). For
\(i\geq1\), Lemma~\ref{lem:logderivlimit}, applied to
\(h_{\nu,j}\), gives
\[
 \frac{h_{\nu,j}^{(i)}}
 {s_\nu^ih_{\nu,j}}
 \convmeas
 (2\partial v_j)^i
\]
locally in \(\Omega\). In particular, every matrix entry is bounded in
measure. Since the determinant is a finite sum of products of these
entries, \(D_\nu\) is also bounded in measure.

We claim that
\begin{equation}\label{eq:determinant-log-small}
 \frac{\log^+|D_\nu|}{s_\nu}
 \convmeas0
 \qquad\text{locally in }\Omega.
\end{equation}
Indeed, fix \(K\Subset\Omega\) and \(\varepsilon>0\). Then
\[
 \left\{
 z\in K:
 \frac{\log^+|D_\nu(z)|}{s_\nu}>\varepsilon
 \right\}
 \subset
 \left\{
 z\in K:
 |D_\nu(z)|>e^{\varepsilon s_\nu}
 \right\}.
\]
Since \(s_\nu\to\infty\), the threshold \(e^{\varepsilon s_\nu}\) tends
to infinity. The boundedness of \(D_\nu\) in measure therefore implies
that the area of the set on the right tends to zero. This proves
\eqref{eq:determinant-log-small}.

\medskip\noindent
\textbf{\underline{Step 3. Take logarithms:}}
Taking absolute values and logarithms in
\eqref{eq:normalized-wronskian} gives, almost everywhere,
\[
 \frac{\log|W(\hv_\nu)|}{s_\nu}
 -
 \frac{\kappa\log s_\nu}{s_\nu}
 -
 \sum_{j=0}^{n}
 \frac{\log|h_{\nu,j}|}{s_\nu}
 =
 \frac{\log|D_\nu|}{s_\nu}.
\]
Hence
\[
 \sum_{j=0}^{n}
 \frac{\log|h_{\nu,j}|}{s_\nu}
 =
 \frac{\log|W(\hv_\nu)|}{s_\nu}
 -
 \frac{\kappa\log s_\nu}{s_\nu}
 -
 \frac{\log|D_\nu|}{s_\nu}.
\]
Using $
-\log|D_\nu|\geq-\log^+|D_\nu|,
$
we obtain
\begin{equation}\label{eq:wronskian-lower-bound}
 \sum_{j=0}^{n}
 \frac{\log|h_{\nu,j}|}{s_\nu}
 \geq
 \frac{\log|W(\hv_\nu)|}{s_\nu}
 -
 \frac{\kappa\log s_\nu}{s_\nu}
 -
 \frac{\log^+|D_\nu|}{s_\nu}.
\end{equation}

\medskip\noindent
\textbf{\underline{Step 4. Pass to the limit:}}
By \eqref{eq:component-log-limits}, after passing to a subsequence we
may assume that
\[
 \frac{\log|h_{\nu,j}|}{s_\nu}
 \longrightarrow v_j
 \qquad\text{a.e.}
\]
simultaneously for \(0\leq j\leq n\). By
\eqref{eq:small-wronskian-local} and
\eqref{eq:determinant-log-small}, we may pass to a further subsequence
such that
\[
 \frac{\log|W(\hv_\nu)|}{s_\nu}\longrightarrow0,
 \qquad
 \frac{\log^+|D_\nu|}{s_\nu}\longrightarrow0
\]
almost everywhere. Finally,
\[
\frac{\kappa\log s_\nu}{s_\nu}\longrightarrow0.
\]
Passing to the limit in \eqref{eq:wronskian-lower-bound} therefore gives
\[
 \sum_{j=0}^{n}v_j\geq0
 \qquad\text{a.e.\ on }\Omega,
\]
which is \eqref{eq:sum-positive}.
\end{proof}

If
$
U:=\max\limits_{0\leq j\le n}v_j,
$
then \eqref{eq:sum-positive} gives
\begin{equation}\label{eq:sandwich}
 U\ge0,\qquad
 -nU\leq v_j\leq U
 \qquad\text{a.e.}
\end{equation}
We shall use the bounds in \eqref{eq:sandwich} in the compactness
argument below.

We shall use the compactness theorem for subharmonic functions (Hartogs' lemma) in the following form.

\begin{lem}[{\cite[Theorem~4.1.9(a),(b)]{Hor03}}]\label{lem:subharmonic-compactness}
Let \((u_\nu)\) be a sequence of subharmonic functions on a domain
\(\Omega\), locally uniformly bounded above. Then every subsequence has a
further subsequence for which one of the following alternatives holds:
\begin{enumerate}[label=\textup{(\roman*)}]
\item \(u_\nu\to-\infty\) locally uniformly on \(\Omega\);
\item \(u_\nu\to u\) in \(L^1_{\loc}(\Omega)\), where \(u\) is subharmonic
and \(u\not\equiv-\infty\).
\end{enumerate}
In the second case, the corresponding Riesz measures converge weakly on
compact subsets. Moreover, if \(u_\nu\to u\) in \(L^1_{\loc}(\Omega)\) and
$
 K\Subset U_0\Subset\Omega,
$
then
\[
\limsup_{\nu\to\infty}\sup_Ku_\nu
 \leq\sup_{U_0}u.
\]
\end{lem}

For vector-valued logarithms we shall repeatedly use
\[
0\leq
 \log\norm{\hv}
 -\max_{0\leq j\le n}\log|h_j|
 \leq\frac12\log {n+1}.
\]
Thus the normalized logarithm of the Euclidean norm and the maximum of the
normalized component logarithms have the same subsequential limits.

Finally, the growth of the characteristic satisfies the following simple dichotomy.

\begin{lem}\label{lem:logdichotomy}
For a nonconstant holomorphic curve, exactly one of the following two
possibilities occurs:
\begin{equation}\label{eq:logdichotomy}
 T(r,\fv)=O(\log r)
 \qquad\text{or}\qquad
 \log r=o\bigl(T(r,\fv)\bigr).
\end{equation}
Moreover, there is a constant \(c>0\) such that
\begin{equation}\label{eq:characteristic-lower-log}
 T(r,\fv)\geq c\log r
\end{equation}
for all sufficiently large \(r\).
\end{lem}

\begin{proof}

Put $
\Phi(x):=T(e^x,\fv).$
Since \(T(r,\fv)\) is nondecreasing and convex as a function of
\(\log r\), the function \(\Phi\) is convex and nondecreasing. It is
also nonconstant because the curve is nonconstant.

\medskip\noindent
\textbf{\underline{Step 1. A logarithmic lower bound:}}
Choose \(x_0<x_1\) such that
\[
\Phi(x_1)>\Phi(x_0).
\]
For \(x\geq x_1\), convexity implies that the secant slope from \(x_0\)
to \(x\) is at least the secant slope from \(x_0\) to \(x_1\):
\begin{equation}\label{eq:positive-secant}
 \frac{\Phi(x)-\Phi(x_0)}{x-x_0}
 \geq
 \frac{\Phi(x_1)-\Phi(x_0)}{x_1-x_0}
 =:c_0>0.
\end{equation}
Hence
\[
\Phi(x)\geq c_0(x-x_0)+\Phi(x_0).
\]
After increasing \(x_1\), if necessary, this gives
$
\Phi(x)\geq c x
$
for some \(c>0\) and all sufficiently large \(x\). Returning to
\(r=e^x\), we obtain
\[
T(r,\fv)\geq c\log r
\]
for all sufficiently large \(r\), which is
\eqref{eq:characteristic-lower-log}.

\medskip\noindent
\textbf{\underline{Step 2. The limiting secant slope:}}
For \(x>x_0\), set
\[
q(x):=
\frac{\Phi(x)-\Phi(x_0)}{x-x_0}.
\]
By convexity of \(\Phi\), \(q(x)\) is nondecreasing in \(x\). Moreover,
\eqref{eq:positive-secant} shows that \(q(x)\geq c_0>0\) for
\(x\geq x_1\). Therefore the limit
$
L:=\lim\limits_{x\to\infty}q(x)
$
exists in \((0,\infty]\).

We distinguish the two possible cases.

\medskip\noindent
\textbf{\underline{Step 3. The two alternatives:}}
Suppose first that \(L<\infty\). Since \(q(x)\) is bounded for large
\(x\), there is a constant \(C>0\) such that
\[
\Phi(x)-\Phi(x_0)\leq C(x-x_0).
\]
Thus
\[
\Phi(x)=O(x),
\qquad x\to\infty.
\]
With \(x=\log r\), this is exactly
\[
T(r,\fv)=O(\log r).
\]

Suppose now that \(L=\infty\). From the definition of \(q(x)\),
\[
\Phi(x)
=
\Phi(x_0)+(x-x_0)q(x),
\]
and therefore
\[
\frac{\Phi(x)}{x}
=
\frac{\Phi(x_0)}{x}
+
\left(1-\frac{x_0}{x}\right)q(x).
\]
Since \(q(x)\to\infty\), it follows that
\[
\frac{\Phi(x)}{x}\longrightarrow\infty.
\]
Substituting \(x=\log r\) gives
\[
\frac{T(r,\fv)}{\log r}\longrightarrow\infty,
\]
or equivalently,
\[
\log r=o\bigl(T(r,\fv)\bigr).
\]

The two cases \(L<\infty\) and \(L=\infty\) are mutually exclusive and
exhaust all possibilities. This proves \eqref{eq:logdichotomy}.
\end{proof}



\section{Estimates for polynomial solution spaces}\label{sec:algebra}

In this section we derive two estimates for finite-dimensional spaces of
polynomials. The first controls the coefficients of the fundamental
differential operator in terms of the zeros of the Wronskian. The second
controls a basis normalized by prescribed initial derivatives.

Our only use of the representation-theoretic results of Karp--Purbhoo
\cite{KP26} is through their universal Pl\"ucker-coordinate formula. We first translate the part of their
result that is needed here into the language of derivative minors. After
this translation, the two estimates follow from elementary linear algebra
and Taylor expansion.

\subsection{Derivative minors and Pl\"ucker coordinates}
\label{sec:derivative-minors}

Let \(V\subset\C[z]\) be an \({n+1}\)-dimensional space of polynomials, and
let \(v_0,\ldots,v_{n}\) be a basis of \(V\). We use polynomial
coordinates in the factorial basis
\[
1,\ z,\ \frac{z^2}{2!},\ \frac{z^3}{3!},\ldots.
\]
The advantage of this normalization is that the coefficient of
\(z^r/r!\) in \(v_j(a+z)\) is exactly \(v_j^{(r)}(a)\).

Let
$
\tau=(\tau_1\geq\cdots\geq\tau_{n+1}\geq0)
$
be a partition, where trailing zeros are allowed. Set
\[
i_k=k-1+\tau_{n+2-k},
\qquad 1\leq k\leq n+1,
\]
and define the derivative minor
\begin{equation}\label{eq:derivative-minor}
 \Delta^\tau(V;a)
 :=
 \det\bigl(v_j^{(i_k)}(a)\bigr)_{
                  1\leq k\leq n+1,\ 0\leq j\leq n}.
\end{equation}
For the empty partition,
\[
\Delta^\varnothing(V;a)
=
W(v_0,\ldots,v_{n})(a).
\]
A change of basis of \(V\) multiplies every determinant
\(\Delta^\tau(V;a)\) by the same nonzero constant. Hence the ratios
\[
\frac{\Delta^\tau(V;a)}
     {\Delta^\varnothing(V;a)}
\]
are independent of the chosen basis.

To relate the derivative minors at \(a\) to the Pl\"ucker coordinates
used in \cite{KP26}, we translate the point \(a\) to the origin. Set
\[
V_a:=\{p(a+z):p\in V\}\subset\C[z].
\]
For each basis element \(v_j\), its Taylor expansion at \(a\) is
\[
v_j(a+z)
=
\sum_{r\geq0}v_j^{(r)}(a)\frac{z^r}{r!}.
\]
Hence, with respect to the factorial basis
$
1,\ z,\ {z^2}/{2!},\ldots,
$
the coefficient of \(z^r/r!\) in \(v_j(a+z)\) is exactly
\(v_j^{(r)}(a)\).

The minor \(\Delta^\tau(V;a)\) in
\eqref{eq:derivative-minor} uses the derivative orders
\[
i_k=k-1+\tau_{n+2-k},
\qquad 1\leq k\leq n+1.
\]
Thus, in the coefficient matrix of the translated space \(V_a\), it
corresponds to the columns
\[
i_k+1=k+\tau_{n+2-k},
\qquad 1\leq k\leq n+1.
\]
By the partition convention in \cite[equation~(2.5)]{KP26}, this is
precisely the column set indexed by \(\tau\). Therefore
\(\Delta^\tau(V;a)\) is the Pl\"ucker-coordinate minor of \(V_a\)
indexed by \(\tau\), up to transposition of the matrix. Since
transposition does not change the determinant, the two minors agree
exactly; in particular, no additional sign or factorial factor occurs.

Let \(P\) be the monic polynomial proportional to the Wronskian of
\(V\), and write
\begin{equation}\label{eq:monic-wronskian}
 P(z)=\prod_{\ell=1}^M(z-a_\ell),
\end{equation}
where the roots are repeated according to multiplicity. The same roots
appear in the universal minors identity \eqref{eq:universal-minors}.
We use the normalization \(P\) in \eqref{eq:monic-wronskian}
throughout this section\footnote{The notation of \cite{KP26} uses separate symbols for the dimension of
the polynomial space and the degree of its Wronskian. In the present
paper these quantities are \(n+1\) and \(M\), respectively; we keep this
notation throughout.}.




\subsection{The universal minors estimate}
\label{sec:universal-minors}

We now state the only consequence of the Karp--Purbhoo theory that will
be needed below. The point of the result is that a derivative minor of
weight \(s\), divided by the Wronskian, is a sum of products of \(s\)
simple factors \((a-a_\ell)^{-1}\), with coefficients bounded only in
terms of \(s\). In particular, the estimate remains uniform when
Wronskian zeros collide.

\begin{lem}\label{lem:universal-minors}
If \(M=0\), the assertion is immediate, with
\(\mathfrak S_0=\{e\}\); hence assume \(M\geq1\) below.
There exist a finite-dimensional unitary representation
\[
\pi:\mathfrak S_M\longrightarrow U(\mathcal H)
\]
and a unit vector \(\ev\in\mathcal H\). For a fixed \(V\), choose this
vector once and for all, independently of the center \(a\), of the
partition below, and of the point at which the estimate is evaluated.
The following identity then holds.

Suppose that \(P(a)\ne0\), let
\(\tau=(\tau_1,\ldots,\tau_{n+1})\) be a partition (trailing zeros
allowed), and put
$
s=|\tau|.
$
Assume first that \(s\leq M\). If \(s>M\), then
\(\tau\not\subseteq\nu\), because \(|\nu|=M\). The corresponding
Pl\"ucker coordinate is zero, and
\cite[Theorem~1.3(v)]{KP26} gives \(\beta^\tau=0\) on the relevant
eigenspace. Thus both sides of \eqref{eq:universal-minors} vanish when
\(s>M\). Identify \(\mathfrak S_I\) with \(\mathfrak S_s\) through the
increasing ordering of \(I\).
Then
\begin{equation}\label{eq:universal-minors}
 \frac{\Delta^\tau(V;a)}{\Delta^\varnothing(V;a)}
 =
 \sum_{\substack{I\subset\{1,\ldots,M\}\\|I|=s}}
 \frac{\langle C_I^\tau\ev,\ev\rangle}
      {\prod_{\ell\in I}(a-a_\ell)},
\end{equation}
where
\begin{equation}\label{eq:CI-definition}
 C_I^\tau
 :=
 \sum_{\sigma\in\mathfrak S_I}
 \chi^\tau(\sigma)\pi(\sigma).
\end{equation}
Here \(\mathfrak S_I\subset\mathfrak S_M\) is the subgroup which
permutes \(I\) and fixes its complement, and \(\chi^\tau\) is the
irreducible character indexed by \(\tau\).

Moreover,
\begin{equation}\label{eq:projectionbound}
 \bigl|\langle C_I^\tau\ev,\ev\rangle\bigr|
 \leq s!,
\end{equation}
and consequently
\begin{equation}\label{eq:minor-es-bound}
 \left|
 \frac{\Delta^\tau(V;a)}{\Delta^\varnothing(V;a)}
 \right|
 \leq
 s!\,
 e_s\bigl(
 |a-a_1|^{-1},\ldots,|a-a_M|^{-1}
 \bigr),
\end{equation}
where
$
e_s(x_1,\ldots,x_M)
=
\sum_{\substack{I\subset\{1,\ldots,M\}\\|I|=s}}
\prod_{\ell\in I}x_\ell,
e_0=1,
$
and \(e_s=0\) for \(s>M\).
\end{lem}

\begin{proof}

\medskip\noindent
\textbf{\underline{Step 1. Place \(V\) in a fibre of the Wronski map:}}
Choose \(m\) larger than every degree occurring in a basis of \(V\),
put \(d=n+1\), and regard \(V\) as a point of
\(\operatorname{Gr}(d,\mathbb C_{m-1}[u])\), where
\(\mathbb C_{m-1}[u]\) denotes the polynomials of degree at most
\(m-1\). Let \(\nu\) be the Schubert-cell index of \(V\). In the
factorial-basis convention of \cite{KP26},
\[
|\nu|=\deg\operatorname{Wr}(V)=M.
\]
Set \(z_\ell=-a_\ell\), \(1\leq\ell\leq M\). By
\cite[Theorem~1.3(v)]{KP26}, there is a joint eigenspace
\(E\subseteq\mathsf M^\nu\) of the Bethe algebra
\(\mathcal B_M(z_1,\ldots,z_M)\) corresponding to \(V\), and the
eigenvalues of the operators \(\beta^\tau(0)\) on \(E\) are the
normalized Pl\"ucker coordinates of \(V\). Thus
\[
P(u)=\prod_{\ell=1}^M(u+z_\ell).
\]

Equip the Specht module containing \(E\) with an invariant Hermitian
inner product, and let
\[
\pi:\mathfrak S_M\longrightarrow U(\mathcal H)
\]
be the resulting unitary representation. Choose a unit vector
$
\ev\in E.
$

\medskip\noindent
\textbf{\underline{Step 2. Translate the center:}}
For \(a\in\C\), recall
\[
V_a=\{p(a+z):p\in V\}\subset\C[z].
\]
The raw derivative minors \(\Delta^\tau(V;a)\) and the normalized
Pl\"ucker coordinates of \(V_a\) differ by a nonzero scalar \(c(a)\)
which is independent of \(\tau\). Consequently, their ratios are
equal. By \cite[Proposition~2.9]{KP26}, the normalized Pl\"cker
coordinates of the translated space \(V_a\) are obtained from those of
\(V\) by a triangular translation formula. By
\cite[Theorem~1.3(ii)]{KP26}, the operators \(\beta^\tau(t)\) satisfy
the same transformation law. Hence \(E\) remains a joint eigenspace,
and the eigenvalue of \(\beta^\tau(a)\) on \(E\) is the normalized
\(\tau\)-coordinate of \(V_a\).

It follows that, whenever \(P(a)\ne0\),
\begin{equation}\label{eq:minor-beta-ratio}
 \frac{\Delta^\tau(V;a)}
      {\Delta^\varnothing(V;a)}
 =
 \frac{\langle\beta^\tau(a)\ev,\ev\rangle}
      {\langle\beta^\varnothing(a)\ev,\ev\rangle}.
\end{equation}

\medskip\noindent
\textbf{\underline{Step 3. Expand the universal Pl\"ucker operator:}}
By \cite[equation~(1.2)]{KP26},
\[
\beta^\tau(a)
=
\sum_{\substack{I\subset\{1,\ldots,M\}\\|I|=s}}
\left(
 \sum_{\sigma\in\mathfrak S_I}
 \chi^\tau(\sigma)\pi(\sigma)
\right)
\prod_{\ell\notin I}(z_\ell+a).
\]
With the notation \eqref{eq:CI-definition} and
\(z_\ell=-a_\ell\), this becomes
\begin{equation}\label{eq:beta-tau-expanded}
 \beta^\tau(a)
 =
 \sum_{|I|=s}
 C_I^\tau
 \prod_{\ell\notin I}(a-a_\ell).
\end{equation}
For the empty partition,
\[
\beta^\varnothing(a)
=
\left(
\prod_{\ell=1}^M(a-a_\ell)
\right)\operatorname{Id}
=
P(a)\operatorname{Id}.
\]
Substituting this and \eqref{eq:beta-tau-expanded} into
\eqref{eq:minor-beta-ratio} gives
\[
\begin{aligned}
\frac{\Delta^\tau(V;a)}
     {\Delta^\varnothing(V;a)}
&=
\frac1{P(a)}
\sum_{|I|=s}
\langle C_I^\tau\ev,\ev\rangle
\prod_{\ell\notin I}(a-a_\ell)\\
&=
\sum_{|I|=s}
\frac{\langle C_I^\tau\ev,\ev\rangle}
     {\prod_{\ell\in I}(a-a_\ell)}.
\end{aligned}
\]
This proves \eqref{eq:universal-minors}.

\medskip\noindent
\textbf{\underline{Step 4. Bound the coefficients:}}
Let \(f^\tau\) denote the number of standard Young tableaux of shape
\(\tau\). By \cite[Proposition~2.16]{KP26}, applied to the restriction
of \(\pi\) to \(\mathfrak S_I\),
$
{f^\tau\,C_I^\tau}/{s!}
$
is an orthogonal projection. Therefore
\[
\|C_I^\tau\|
\leq
\frac{s!}{f^\tau}
\leq s!.
\]
Since \(\|\ev\|=1\),
\[
\bigl|\langle C_I^\tau\ev,\ev\rangle\bigr|
\leq
\|C_I^\tau\|
\leq s!,
\]
which proves \eqref{eq:projectionbound}.

\medskip\noindent
\textbf{\underline{Step 5. Sum over the Wronskian zeros:}}
Taking absolute values in \eqref{eq:universal-minors} and using
\eqref{eq:projectionbound}, we obtain
\[
\begin{aligned}
\left|
\frac{\Delta^\tau(V;a)}
     {\Delta^\varnothing(V;a)}
\right|
&\leq
s!
\sum_{|I|=s}
\prod_{\ell\in I}|a-a_\ell|^{-1}\\
&=
s!\,
e_s\bigl(
|a-a_1|^{-1},\ldots,|a-a_M|^{-1}
\bigr).
\end{aligned}
\]
This is \eqref{eq:minor-es-bound}. The identity is first obtained
when the roots \(a_1,\ldots,a_M\) are pairwise distinct. After
multiplication by \(P(a)\), both sides are polynomial expressions in the
roots, so the identity extends to repeated roots by continuity.
\end{proof}

\begin{rem}\label{rem:universal-minors-use}
Lemma~\ref{lem:universal-minors} is the only place in the proof where
the representation-theoretic description of the Wronski map is used.
For the rest of the paper, the important point is simply the following:
a derivative minor of weight \(s\), divided by the Wronskian, is a sum
of products of \(s\) factors \((a-a_\ell)^{-1}\), and the coefficients
of these products are bounded by \(s!\). The bound is independent of
the positions of the zeros \(a_\ell\).
\end{rem}


\subsection{Coefficients of the fundamental differential operator}
\label{sec:polynomial-coefficients}

We first apply Lemma~\ref{lem:universal-minors} to the differential
equation whose solution space is \(V\).

\begin{prop}\label{prop:polynomialcoeff}
Let \(V\subset\C[z]\) have dimension \({n+1}\), and let
\(a_1,\ldots,a_M\) be the zeros of its Wronskian, repeated according
to multiplicity. Write its monic fundamental operator as
\[
L_V
=
D^{n+1}+\sum_{q=1}^{n+1} b_q(z)D^{n+1-q}.
\]
Then
\begin{equation}\label{eq:productformula}
 b_q(z)
 =
 (-1)^q
 \sum_{\substack{I\subset\{1,\ldots,M\}\\|I|=q}}
 \frac{\gamma_{q,I}}
      {\prod_{i\in I}(z-a_i)},
 \qquad
 |\gamma_{q,I}|\leq q!,
 \qquad
 1\leq q\leq n+1,
\end{equation}
where the constants \(\gamma_{q,I}\) are independent of \(z\).
\end{prop}

\begin{proof}

\medskip\noindent
\textbf{\underline{Step 1. Express \(b_q\) by Cramer's rule:}}
Let \(v_0,\ldots,v_{n}\) be a basis of \(V\). At a point where the
Wronskian does not vanish, the equations
\[
v_j^{(n+1)}
+
b_1v_j^{(n)}
+\cdots+
b_{n+1}v_j
=0,
\qquad
0\leq j\leq n,
\]
form a nonsingular linear system for
\(b_1,\ldots,b_{n+1}\).

For \(b_q\), Cramer's rule replaces the derivative row of order
\({n+1}-q\) by the row of order \({n+1}\). After arranging the derivative
orders increasingly, they are
\[
0,1,\ldots,{n+1}-q-1,
{n+1}-q+1,\ldots,{n+1}.
\]
By the convention in \eqref{eq:derivative-minor}, this is the
partition
\[
(1^q)
=
(\underbrace{1,\ldots,1}_{q\text{ times}}).
\]
Keeping track of the row permutation gives
\begin{equation}\label{eq:cramer-minor}
 b_q(z)
 =
 (-1)^q
 \frac{\Delta^{(1^q)}(V;z)}
      {\Delta^\varnothing(V;z)}.
\end{equation}

\medskip\noindent
\textbf{\underline{Step 2. Apply Lemma~\ref{lem:universal-minors}:}}
The partition \((1^q)\) has size \(q\). Hence
\eqref{eq:universal-minors} gives
\begin{equation}\label{eq:q-minor-expansion}
\frac{\Delta^{(1^q)}(V;z)}
     {\Delta^\varnothing(V;z)}
=
\sum_{|I|=q}
\frac{\gamma_{q,I}}
     {\prod_{i\in I}(z-a_i)},
\end{equation}
where
$
\gamma_{q,I}
=
\langle C_I^{(1^q)}\ev,\ev\rangle.
$
By \eqref{eq:projectionbound},
\[
|\gamma_{q,I}|\leq q!.
\]
Substituting \eqref{eq:q-minor-expansion} into
\eqref{eq:cramer-minor} proves \eqref{eq:productformula}. The coefficients \(\gamma_{q,I}\) are
independent of \(z\), because the representation and the vector
\(\ev\) in Lemma~\ref{lem:universal-minors} are fixed once \(V\) is
fixed.

\medskip\noindent
\textbf{\underline{Step 3. The case \(M=0\):}}
Suppose that the Wronskian is constant. Choose a basis of \(V\) with
distinct polynomial degrees
$
d_0<d_1<\cdots<d_{n}.
$
The degree of the Wronskian is
\[
\sum_{j=0}^{n}d_j-\frac{n(n+1)}2.
\]
Since this degree is zero and \(d_j\geq j\), we must have
\[
d_j=j,
\qquad
0\leq j\leq n.
\]
Thus \(V\) is the space of all polynomials of degree at most \(n\),
and hence \(L_V=D^{n+1}\).
This is exactly the empty-sum case of \eqref{eq:productformula}.
\end{proof}

\begin{rem}\label{rem:product-form}
The product form in \eqref{eq:productformula} is important. One could
decompose the terms further into partial fractions when the roots are
distinct, but the resulting coefficients need not remain bounded when
roots approach one another. In contrast,
\[
|\gamma_{q,I}|\leq q!
\]
is uniform and does not depend on the mutual distances between the
Wronskian zeros.
\end{rem}


\subsection{A basis with prescribed initial derivatives}
\label{sec:initial-basis}

We next apply Lemma~\ref{lem:universal-minors} to the solutions
themselves. The resulting estimate is again independent of the degree
of the ambient polynomial space.

\begin{prop}\label{prop:initial-basis}
Let \(V\) and \(a_1,\ldots,a_M\) be as in
Proposition~\ref{prop:polynomialcoeff}. Suppose that \(a\) is not a
Wronskian zero, and let
\(\psi_0,\ldots,\psi_{n}\in V\) satisfy
\begin{equation}\label{eq:initial-normalization}
 \psi_j^{(i)}(a)=\delta_{ij},
 \qquad
 0\leq i,j\leq n.
\end{equation}
Then
\begin{equation}\label{eq:basismajorant}
 |\psi_j(a+z)|
 \leq
 \frac{|z|^j}{j!}
 \prod_{\ell=1}^M
 \left(
 1+\frac{|z|}{|a-a_\ell|}
 \right),
 \qquad
 0\leq j\leq n.
\end{equation}
\end{prop}

\begin{proof}

\medskip\noindent
\textbf{\underline{Step 1. Express the higher derivatives as minors:}}
By \eqref{eq:initial-normalization}, the matrix
\[
\bigl(\psi_j^{(i)}(a)\bigr)_{0\leq i,j\leq n}
\]
is the identity matrix. In particular,
\begin{equation}\label{eq:initial-wronskian-one}
 \Delta^\varnothing(V;a)=1.
\end{equation}

Fix \(j\in\{0,\ldots,n\}\) and \(k\geq n+1\). Replace derivative row
\(j\) of this identity matrix by the derivative row of order \(k\).
The determinant of the resulting matrix is
\(\psi_j^{(k)}(a)\).

After arranging the derivative orders increasingly, they become
\[
0,\ldots,j-1,j+1,\ldots,n,k.
\]
The corresponding partition is
\begin{equation}\label{eq:derivative-partition}
 \tau=
 \bigl(k-n,1^{\,n-j}\bigr),
\end{equation}
and therefore \eqref{eq:derivative-partition} has size
\[
|\tau|
=
(k-n)+(n-j)
=
k-j.
\]
The row permutation used to put the derivative orders in increasing
order contributes only a sign. Hence
\begin{equation}\label{eq:higher-derivative-minor}
 |\psi_j^{(k)}(a)|
 =
 |\Delta^\tau(V;a)|.
\end{equation}

\medskip\noindent
\textbf{\underline{Step 2. Apply the universal minors estimate:}}
Set
\[
x_\ell:=|a-a_\ell|^{-1},
\qquad
1\leq\ell\leq M.
\]
Using \eqref{eq:initial-wronskian-one},
\eqref{eq:higher-derivative-minor}, and
\eqref{eq:minor-es-bound} with \(s=k-j\), we obtain
\[
|\psi_j^{(k)}(a)|
\leq
(k-j)!\,
e_{k-j}(x_1,\ldots,x_M).
\]
Therefore
\begin{equation}\label{eq:higher-derivative-bound}
 \frac{|\psi_j^{(k)}(a)|}{k!}
 \leq
 \frac{(k-j)!}{k!}
 e_{k-j}(x_1,\ldots,x_M).
\end{equation}
We use \eqref{eq:higher-derivative-bound} together with
\(k\geq n+1>j\), so that
\[
\frac{(k-j)!}{k!}
=
\frac1{k(k-1)\cdots(k-j+1)}
\leq
\frac1{j!}.
\]
Thus
\begin{equation}\label{eq:higher-derivative-simple}
 \frac{|\psi_j^{(k)}(a)|}{k!}
 \leq
 \frac1{j!}
 e_{k-j}(x_1,\ldots,x_M),
 \qquad k\geq n+1.
\end{equation}

\medskip\noindent
\textbf{\underline{Step 3. Sum the Taylor series:}}
The initial conditions \eqref{eq:initial-normalization} give
\[
\psi_j(a+z)
=
\frac{z^j}{j!}
+
\sum_{k={n+1}}^{\infty}
\frac{\psi_j^{(k)}(a)}{k!}z^k.
\]
Using \eqref{eq:higher-derivative-simple},
\[
\begin{aligned}
|\psi_j(a+z)|
&\leq
\frac{|z|^j}{j!}
+
\frac1{j!}
\sum_{k={n+1}}^{\infty}
e_{k-j}(x_1,\ldots,x_M)|z|^k\\
&=
\frac{|z|^j}{j!}
\left(
1+
\sum_{k={n+1}}^{\infty}
e_{k-j}(x_1,\ldots,x_M)|z|^{k-j}
\right)\\
&\leq
\frac{|z|^j}{j!}
\sum_{s\geq0}
e_s(x_1,\ldots,x_M)|z|^s.
\end{aligned}
\]
The generating identity for the elementary symmetric functions is
\[
\sum_{s\geq0}
e_s(x_1,\ldots,x_M)t^s
=
\prod_{\ell=1}^M(1+t x_\ell).
\]
Taking \(t=|z|\) yields
\[
|\psi_j(a+z)|
\leq
\frac{|z|^j}{j!}
\prod_{\ell=1}^M
\left(
1+\frac{|z|}{|a-a_\ell|}
\right),
\]
which is \eqref{eq:basismajorant}.
\end{proof}


\subsection{The initial-value estimate}
\label{sec:initial-value-bound}

Proposition~\ref{prop:initial-basis} immediately yields the form that
will be used in the local compactness argument.

Let \(h\in V\). Since the basis
\(\psi_0,\ldots,\psi_{n}\) satisfies
\eqref{eq:initial-normalization}, we have
\[
h
=
\sum_{i=0}^{n}
h^{(i)}(a)\psi_i.
\]
Applying \eqref{eq:basismajorant} gives
\[
|h(a+z)|
\leq
\left(
\sum_{i=0}^{n}
\frac{|h^{(i)}(a)|\,|z|^i}{i!}
\right)
\prod_{\ell=1}^M
\left(
1+\frac{|z|}{|a-a_\ell|}
\right).
\]
Using \(1+x\leq e^x\), we obtain
\begin{equation}\label{eq:initial-value-bound}
 |h(a+z)|
 \leq
 \left(
 \sum_{i=0}^{n}
 \frac{|h^{(i)}(a)|\,|z|^i}{i!}
 \right)
 e^{|z|B_V(a)},
 \qquad
 B_V(a)
 =
 \sum_{\ell=1}^M|a-a_\ell|^{-1}.
\end{equation}

Both \eqref{eq:productformula} and
\eqref{eq:initial-value-bound} are uniform with respect to collisions
among the Wronskian zeros. This is the feature of the polynomial
estimates that will be used in Section~\ref{sec:compactness}.

\section{Local compactness}\label{sec:compactness}

The estimates of Section~\ref{sec:algebra} apply to polynomial solution
spaces. In this section we transfer them to local holomorphic frames by
Taylor approximation.

We proceed in three stages. First, we prove a local compactness result
for the coefficients of the fundamental differential operator. The main
point is to separate the \(o(s_\nu)\) Wronskian zeros that remain in a
fixed disk from the \(O(s_\nu)\) zeros lying outside that disk. Second,
we construct, after rescaling the original curve, a local representation
to which the compactness result applies. Finally, we retain the
polynomial approximants and the estimates on their Wronskians for later
use.


\subsection{Compactness of the local differential coefficients}
\label{sec:local-coefficient-compactness}

We begin with a local statement. The frame is holomorphic, while its
Wronskian is assumed to be a polynomial whose degree is small compared
with the normalization \(s_\nu\).

\begin{prop}\label{prop:localcompact}
Let \(s_\nu\to\infty\), and let
\(\gv_\nu=(g_{\nu,0},\ldots,g_{\nu,n})\) be holomorphic on \(D_{32}\).
Suppose that, for a fixed constant \(C\),
\begin{equation}\label{eq:localhyp}
 \begin{gathered}
 \sup_{D_{32}}\log\norm{\gv_\nu}\leq Cs_\nu,
 \qquad
 W(\gv_\nu)=P_\nu,\\
 P_\nu\text{ monic},
 \qquad
 m_\nu:=\deg P_\nu=o(s_\nu).
 \end{gathered}
\end{equation}
Let
\(L_\nu=D^{n+1}+\sum_{q=1}^{n+1}b_{q,\nu}D^{n+1-q}\)
be the fundamental operator of \(\gv_\nu\).
Then the sequences \(b_{q,\nu}/s_\nu^q\), \(1\leq q\leq n+1\), are
simultaneously relatively compact locally in measure on \(D_4\).
Every subsequential limit is holomorphic and is bounded on \(D_1\) by
a constant depending only on \({n+1}\) and \(C\).

If \(\widehat b_{q,\nu}\) denotes the coefficient of \(D^{n+1-q}\) in
the fundamental operator of \(P_\nu^{-1/(n+1)}\gv_\nu\), interpreted
locally away from the zeros of \(P_\nu\) and then extended
meromorphically, then
\begin{equation}\label{eq:gaugecomparison}
 \frac{\widehat b_{q,\nu}-b_{q,\nu}}{s_\nu^q}
 \convmeas0
 \qquad\text{locally on }D_4.
\end{equation}
Moreover, \(\widehat b_{1,\nu}=0\).
\end{prop}

\begin{proof}

\medskip\noindent
\textbf{\underline{Step 1. Polynomial approximation:}}
Fix \(A>0\), to be chosen at the end of Step~2. For
\(0\leq j\leq n\), write
\(g_{\nu,j}(z)=\sum_{m=0}^{\infty}c_{\nu,j,m}z^m\).
By \eqref{eq:localhyp},
\(\sup_{D_{32}}|g_{\nu,j}|\leq e^{Cs_\nu}\).
Cauchy's estimate on the circle \(|z|=24\) therefore gives
\begin{equation}\label{eq:taylor-coefficient-bound}
 |c_{\nu,j,m}|
 \leq
 e^{Cs_\nu}24^{-m},
 \qquad m\geq0.
\end{equation}

For an integer \(N\geq n+1\), let
\(p_{\nu,j}(z)=\sum_{m=0}^{N}c_{\nu,j,m}z^m\).
If \(0\leq k\leq n+1\) and \(|z|\leq16\), then
\[
\begin{aligned}
 |g_{\nu,j}^{(k)}(z)-p_{\nu,j}^{(k)}(z)|
 &\leq
 \sum_{m>N}
 \frac{m!}{(m-k)!}|c_{\nu,j,m}|\,16^{m-k}  \\
 &\leq
 C_{n+1} e^{Cs_\nu}
 \sum_{m>N}m^{n+1}\left(\frac23\right)^m,
\end{aligned}
\]
where \eqref{eq:taylor-coefficient-bound} was used in the second
inequality. Since
\(\sum_{m>N}m^{n+1}(2/3)^m
 \leq C_{n+1}(N+1)^{n+1}(2/3)^N\),
we obtain
\begin{equation}\label{eq:taylor-tail-bound}
 \max_{0\leq k\leq n+1}
 \sup_{D_{16}}
 |g_{\nu,j}^{(k)}-p_{\nu,j}^{(k)}|
 \leq
 C_{n+1} e^{Cs_\nu}(N+1)^{n+1}
 \left(\frac23\right)^N.
\end{equation}

Choose \(L=L({n+1},C,A)>0\) so that
\(L\log(3/2)>A+C+1\), and put
\(N_\nu=\lceil Ls_\nu\rceil\).
Since \(\log(N_\nu+1)=o(s_\nu)\), equation
\eqref{eq:taylor-tail-bound} gives, for all sufficiently large \(\nu\),
\begin{equation}\label{eq:taylorerror}
 \max_{0\leq k\leq n+1}
 \sup_{D_{16}}
 |p_{\nu,j}^{(k)}-g_{\nu,j}^{(k)}|
 \leq e^{-As_\nu}.
\end{equation}
Thus the \(p_{\nu,j}\) have a common degree
\(N_\nu=O_{n+1,C,A}(s_\nu)\).

Put
\[
 \pv_\nu=(p_{\nu,0},\ldots,p_{\nu,n}),
 \qquad
 R_\nu:=W(\pv_\nu).
\]
Cauchy's estimates applied to \eqref{eq:localhyp}, followed by
\eqref{eq:taylorerror}, give a constant \(C_1=C_1({n+1},C)\) such that
\begin{equation}\label{eq:derivative-growth}
 \max_{0\leq j\leq n}
 \max_{0\leq k\leq n+1}
 \sup_{D_{16}}
 \left(
 |g_{\nu,j}^{(k)}|
 +
 |p_{\nu,j}^{(k)}|
 \right)
 \leq e^{C_1s_\nu}
\end{equation}
for all sufficiently large \(\nu\).

Each term in the difference
\(W(\pv_\nu)-W(\gv_\nu)\) contains one factor bounded by
\(e^{-As_\nu}\) by \eqref{eq:taylorerror}, while all remaining factors
are bounded by \eqref{eq:derivative-growth}. Hence, after increasing
\(A\) if necessary,
\begin{equation}\label{eq:wronskiapprox}
 \deg R_\nu\leq (n+1)N_\nu=O(s_\nu),
 \qquad
 \sup_{D_{12}}|R_\nu-P_\nu|
 \leq e^{-Bs_\nu},
\end{equation}
where \(B=B(A,{n+1},C)\) and \(B\to\infty\) as \(A\to\infty\).

Since \(P_\nu\) is monic of degree \(m_\nu\), Cauchy's coefficient
estimate gives
\[
 \sup_{|z|\leq12}|P_\nu(z)|\geq12^{m_\nu}\geq1.
\]
Hence, for all sufficiently large \(\nu\),
\[
 \sup_{|z|\leq12}|R_\nu(z)-P_\nu(z)|<1,
\]
and therefore \(R_\nu\not\equiv0\). Thus the fundamental operator of
\(\pv_\nu\) is well defined.


Let \(\mathcal C_{q,\nu}^p\) and \(\mathcal C_{q,\nu}^g\) denote the
Cramer numerators defining, respectively, the coefficients of the
fundamental operators of \(\pv_\nu\) and \(\gv_\nu\). These
determinants involve derivatives of order at most \({n+1}\). The same
determinant expansion, using \eqref{eq:taylorerror} and
\eqref{eq:derivative-growth}, gives
\begin{equation}\label{eq:cramer-numerator-approx}
 \sup_{D_{12}}
 |\mathcal C_{q,\nu}^p-\mathcal C_{q,\nu}^g|
 \leq e^{-Bs_\nu},
 \qquad
 \sup_{D_{12}}|\mathcal C_{q,\nu}^g|
 \leq e^{C_2s_\nu},
\end{equation}
after decreasing \(B\) by a constant depending only on \({n+1}\) and
\(C\). Here \(C_2=C_2({n+1},C)\).


\medskip\noindent
\textbf{\underline{Step 2. Comparison of the two fundamental operators:}}
We first remove the set on which the denominator \(P_\nu\) is
exponentially small.

For a monic polynomial
\(P(z)=\prod_{i=1}^{m}(z-a_i)\), subadditivity of \(\log^-\) gives
\begin{equation}\label{eq:polynomial-negative-area}
 \int_{D_{12}}\log^-|P(z)|\,\dA
 \leq
 \sum_{i=1}^{m}
 \int_{\C}\log^-|z-a_i|\,\dA
 =
 \frac{\pi m}{2}.
\end{equation}
Define
\begin{equation}\label{eq:smallvalues}
 E_\nu
 :=
 \{z\in D_{12}:|P_\nu(z)|<e^{-s_\nu}\}.
 \qquad
 \meas(E_\nu)
 \leq
 \frac{\pi m_\nu}{2s_\nu}
 \longrightarrow0.
\end{equation}
Indeed, \(\log^-|P_\nu|>s_\nu\) on \(E_\nu\), so the measure estimate
follows directly from \eqref{eq:polynomial-negative-area} and
\(m_\nu=o(s_\nu)\).

On \(D_{12}\setminus E_\nu\), one has
\(|P_\nu|\geq e^{-s_\nu}\). If \(B>1\), then
\eqref{eq:wronskiapprox} gives, for all sufficiently large \(\nu\),
\begin{equation}\label{eq:denominator-lower-bounds}
 |P_\nu|\geq e^{-s_\nu},
 \qquad
 |R_\nu|\geq\frac12e^{-s_\nu}
 \quad\text{on }D_{12}\setminus E_\nu.
\end{equation}

Let \(b_{q,\nu}^p\) be the coefficient of \(D^{n+1-q}\) in the
fundamental operator of \(\pv_\nu\). Cramer's rule gives
\(b_{q,\nu}^p=\mathcal C_{q,\nu}^p/R_\nu\) and
\(b_{q,\nu}=\mathcal C_{q,\nu}^g/P_\nu\). Hence
\begin{equation}\label{eq:cramer-difference-bound}
 \left|b_{q,\nu}^p-b_{q,\nu}\right|
 \leq
 \frac{|\mathcal C_{q,\nu}^p-\mathcal C_{q,\nu}^g|}
      {|R_\nu|}
 +
 \frac{|\mathcal C_{q,\nu}^g|\,|R_\nu-P_\nu|}
      {|P_\nu|\,|R_\nu|}.
\end{equation}
Equations \eqref{eq:wronskiapprox},
\eqref{eq:cramer-numerator-approx}, and
\eqref{eq:denominator-lower-bounds} show that the two terms on the
right-hand side of \eqref{eq:cramer-difference-bound} are bounded,
respectively, by
\(2e^{-(B-1)s_\nu}\) and
\(2e^{-(B-C_2-2)s_\nu}\).

Choose \(A\) in Step~1 so large that the corresponding \(B\) satisfies
\(B>C_2+2\). Then
\(b_{q,\nu}^p-b_{q,\nu}\to0\) uniformly on
\(D_{12}\setminus E_\nu\). Since
\(\meas(E_\nu)\to0\) by \eqref{eq:smallvalues}, we obtain
\begin{equation}\label{eq:cramercomparison}
 b_{q,\nu}^p-b_{q,\nu}
 \convmeas0
 \qquad\text{locally on }D_{12}.
\end{equation}


\medskip\noindent
\textbf{\underline{Step 3. A circle separating the Wronskian zeros:}}
Write
\(P_\nu(z)=\prod_{i=1}^{m_\nu}(z-a_{\nu,i})\), with multiplicities.
For every \(c\geq0\),
\(\int_8^{10}\log^-|r-c|\,\dd r\leq2\).
Integrating the sum over \(i\) and then choosing one radius in
\((8,10)\), we obtain \(\eta_\nu\in(8,10)\) such that
\begin{equation}\label{eq:separating-radius}
 \sum_{i=1}^{m_\nu}
 \log^-|\eta_\nu-|a_{\nu,i}||
 \leq2m_\nu.
\end{equation}

For \(|z|=\eta_\nu\),
\(|z-a_{\nu,i}|
 \geq|\eta_\nu-|a_{\nu,i}||\).
Therefore \eqref{eq:separating-radius} gives
\begin{equation}\label{eq:P-circle-lower}
 |P_\nu(z)|
 \geq e^{-2m_\nu},
 \qquad |z|=\eta_\nu.
\end{equation}
Since \(m_\nu=o(s_\nu)\), equations
\eqref{eq:wronskiapprox} and \eqref{eq:P-circle-lower} imply
\(|R_\nu-P_\nu|<|P_\nu|\) on \(|z|=\eta_\nu\) for all sufficiently
large \(\nu\). Rouch\'e's theorem therefore gives
\begin{equation}\label{eq:rouche-root-count}
 n_{R_\nu}(\eta_\nu)
 =
 n_{P_\nu}(\eta_\nu)
 \leq m_\nu
 =
 o(s_\nu).
\end{equation}
At the same time,
\(\deg R_\nu=O(s_\nu)\) by \eqref{eq:wronskiapprox}.


\medskip\noindent
\textbf{\underline{Step 4. Interior and exterior roots:}}
Let
\(\zeta_{\nu,1},\ldots,\zeta_{\nu,d_\nu}\) be the zeros of \(R_\nu\),
repeated according to multiplicity, where
\(d_\nu=\deg R_\nu\). Define
\[
 \mathcal I_\nu
 :=
 \{r:|\zeta_{\nu,r}|<\eta_\nu\},
 \qquad
 \mathcal O_\nu
 :=
 \{1,\ldots,d_\nu\}\setminus\mathcal I_\nu,
\]
and
\[
 U_{\nu,\mathrm{in}}(z)
 :=
 \sum_{r\in\mathcal I_\nu}
 \frac1{|z-\zeta_{\nu,r}|},
 \qquad
 U_{\nu,\mathrm{out}}(z)
 :=
 \sum_{r\in\mathcal O_\nu}
 \frac1{|z-\zeta_{\nu,r}|}.
\]

There is a constant \(C_0\), independent of \(a\in\C\), such that
\(\||z-a|^{-1}\|_{L^{3/2}(D_6)}\leq C_0\).
By \eqref{eq:rouche-root-count},
\(\#\mathcal I_\nu=o(s_\nu)\), and Minkowski's inequality gives
\[
 \left\|
 \frac{U_{\nu,\mathrm{in}}}{s_\nu}
 \right\|_{L^{3/2}(D_6)}
 \leq
 C_0\frac{\#\mathcal I_\nu}{s_\nu}
 \longrightarrow0.
\]
If \(r\in\mathcal O_\nu\) and \(z\in D_6\), then
\(|z-\zeta_{\nu,r}|\geq\eta_\nu-6>2\).
Since \(d_\nu=O(s_\nu)\) by \eqref{eq:wronskiapprox},
\(\sup_{D_6}U_{\nu,\mathrm{out}}\leq C s_\nu\).
Thus
\begin{equation}\label{eq:rootbounds}
 \left\|
 \frac{U_{\nu,\mathrm{in}}}{s_\nu}
 \right\|_{L^{3/2}(D_6)}
 \longrightarrow0,
 \qquad
 \sup_{D_6}
 \frac{U_{\nu,\mathrm{out}}}{s_\nu}
 \leq C_3.
\end{equation}


\medskip\noindent
\textbf{\underline{Step 5. Coefficient contributions:}}
Apply Proposition~\ref{prop:polynomialcoeff} to the space spanned by
the components of \(\pv_\nu\). Multiplying
\(R_\nu\) by a nonzero constant does not change either its zero divisor
or the fundamental operator, so the monic normalization of \(R_\nu\)
may be used in that proposition.

For each \(q\), split the sum in \eqref{eq:productformula} into the
terms whose root set is contained in \(\mathcal O_\nu\) and the terms
whose root set meets \(\mathcal I_\nu\). This defines
\begin{equation}\label{eq:splitcoeff}
 \frac{b_{q,\nu}^p}{s_\nu^q}
 =
 B_{q,\nu}+F_{q,\nu}.
\end{equation}
The terms in \(B_{q,\nu}\) have poles only at exterior roots, hence
\(B_{q,\nu}\) is holomorphic on \(D_8\). By
\eqref{eq:productformula} and the second estimate in
\eqref{eq:rootbounds},
\begin{equation}\label{eq:B-bound}
 \sup_{D_6}|B_{q,\nu}|
 \leq C_q.
\end{equation}

Every term in \(F_{q,\nu}\) contains at least one interior root.
For \(q\ge2\), using \(|\gamma_{q,I}|\leq q!\) from
\eqref{eq:productformula} and estimating subsets by ordered
\(q\)-tuples gives
\begin{equation}\label{eq:F-bound}
 |F_{q,\nu}|
 \leq
 C_q
 \frac{U_{\nu,\mathrm{in}}}{s_\nu}
 \left(
 \frac{U_{\nu,\mathrm{in}}+U_{\nu,\mathrm{out}}}{s_\nu}
 \right)^{q-1}.
\end{equation}

Set \(A_\nu=U_{\nu,\mathrm{in}}/s_\nu\). For \(q=1\), the preceding
estimate gives
\[
 |F_{1,\nu}|\le C_1A_\nu,
 \qquad
 \int_{D_6}|F_{1,\nu}|^{3/2}\,\dA
 \le C_1^{3/2}\int_{D_6}A_\nu^{3/2}\,\dA
 \longrightarrow0
\]
by \eqref{eq:rootbounds}. Hence assume \(q\ge2\) below and set
\[
G_\nu=\frac{U_{\nu,\mathrm{in}}+U_{\nu,\mathrm{out}}}{s_\nu},
\qquad p=\frac{3}{2q}.
\]
Raising \eqref{eq:F-bound} to the power \(p\) and applying H\"older's
inequality with exponents \(q\) and \(q/(q-1)\) gives
\[
 \int_{D_6}|F_{q,\nu}|^p\,\dA
 \leq
 C_q
 \left(
 \int_{D_6}A_\nu^{3/2}\,\dA
 \right)^{1/q}
 \left(
 \int_{D_6}G_\nu^{3/2}\,\dA
 \right)^{(q-1)/q}.
\]
The first factor tends to zero by \eqref{eq:rootbounds}, while the
second remains bounded by the same equation. Hence
\begin{equation}\label{eq:errorvanish}
 \int_{D_6}|F_{q,\nu}|^{3/(2q)}\,\dA
 \longrightarrow0.
\end{equation}
When \(3/(2q)<1\), \eqref{eq:errorvanish} is only an estimate for an
integral of a positive power; no Banach-space assertion is being made.

By \eqref{eq:B-bound}, the families \(B_{q,\nu}\) are normal on
\(D_6\). Since there are only finitely many values of \(q\), Montel's
theorem gives a common subsequence on which all \(B_{q,\nu}\)
converge locally uniformly on \(D_6\). Equations
\eqref{eq:splitcoeff} and \eqref{eq:errorvanish} then show that
\(b_{q,\nu}^p/s_\nu^q\) converge locally in measure on \(D_4\) to
holomorphic limits. Finally, \eqref{eq:cramercomparison} gives
\[
 \frac{b_{q,\nu}^p-b_{q,\nu}}{s_\nu^q}
 \convmeas0
 \qquad\text{locally on }D_4,
\]
because \(s_\nu\to\infty\). Therefore
\(b_{q,\nu}/s_\nu^q\) have the same subsequential limits. This proves
the compactness and holomorphicity assertions of the proposition.
The bound on \(D_1\) follows from \eqref{eq:B-bound}; after the choice
of \(A\) in Step~2 has been fixed in terms of \({n+1}\) and \(C\), the
bound depends only on \({n+1}\) and \(C\).


\medskip\noindent
\textbf{\underline{Step 6. Passage to the canonical gauge:}}
By Lemma~\ref{lem:fundamental-operator}(iii) and
\(W(\gv_\nu)=P_\nu\),
\(b_{1,\nu}=-P_\nu'/P_\nu\).
Set
\(\ell_\nu=P_\nu'/((n+1)P_\nu)\).
If the zeros of \(P_\nu\) are listed with multiplicity, then
\[
 \ell_\nu(z)
 =
 \frac1{n+1}\sum_{i=1}^{m_\nu}\frac1{z-a_{\nu,i}},
 \qquad
 \ell_\nu^{(j)}(z)
 =
 \frac{(-1)^j j!}{{n+1}}
 \sum_{i=1}^{m_\nu}
 \frac1{(z-a_{\nu,i})^{j+1}}
 \quad(j\geq1).
\]

For \(j=0\), the uniform \(L^{3/2}\)-estimate used in Step~4 gives
\(\|\ell_\nu/s_\nu\|_{L^{3/2}(K)}
 \leq C_Km_\nu/s_\nu\to0\)
for every \(K\Subset D_4\).

For \(j\geq1\), choose
\(1/(j+1)<\alpha<\min\{1,2/(j+1)\}\).
Since \(\alpha<1\), subadditivity of \(x^\alpha\), together with local
integrability of
\(|z-a|^{-(j+1)\alpha}\), gives
\[
 \int_K
 \left|
 \frac{\ell_\nu^{(j)}}{s_\nu^{j+1}}
 \right|^\alpha
 \dA
 \leq
 C_{K,j}
 \frac{m_\nu}{s_\nu^{(j+1)\alpha}}.
\]
Here \(m_\nu=o(s_\nu)\) by \eqref{eq:localhyp}, while
\((j+1)\alpha>1\). Therefore
\begin{equation}\label{eq:gaugesmall}
 \frac{\ell_\nu^{(j)}}{s_\nu^{j+1}}
 \convmeas0
 \qquad\text{locally on }D_4,
 \qquad j\geq0.
\end{equation}

The fundamental operator of \(P_\nu^{-1/(n+1)}\gv_\nu\) is
\(P_\nu^{-1/(n+1)}L_\nu P_\nu^{1/{n+1}}\). Since
\(P_\nu^{-1/(n+1)}DP_\nu^{1/{n+1}}=D+\ell_\nu\), it is obtained from
\(L_\nu\) by replacing \(D\) with \(D+\ell_\nu\). The coefficient of
\(D^{n}\) is therefore
\(b_{1,\nu}+{n+1}\ell_\nu=0\), so
\(\widehat b_{1,\nu}=0\).

For \(q\geq2\), repeated use of
\(D\phi=\phi D+\phi'\) shows that
\(\widehat b_{q,\nu}-b_{q,\nu}\) is a finite sum of monomials of the
form
\begin{equation}\label{eq:gauge-monomial}
 C
 \prod_{r=1}^{q-1}b_{r,\nu}^{\,m_r}
 \prod_{j=0}^{q-1}
 \bigl(\ell_\nu^{(j)}\bigr)^{k_j},
 \qquad
 \sum_{r=1}^{q-1}r\,m_r
 +
 \sum_{j=0}^{q-1}(j+1)k_j
 =
 q,
 \qquad
 \sum_{j=0}^{q-1}k_j\geq1.
\end{equation}
The normalized factors \(b_{r,\nu}/s_\nu^r\) are bounded in measure
by Step~5, while every normalized factor
\(\ell_\nu^{(j)}/s_\nu^{j+1}\) tends to zero in measure by
\eqref{eq:gaugesmall}. Dividing each monomial in
\eqref{eq:gauge-monomial} by \(s_\nu^q\) therefore makes it tend to
zero in measure. Since only finitely many monomials occur, this proves
\eqref{eq:gaugecomparison}.
\end{proof}

The distinction between the two groups of roots in Steps~3--5 is
essential. Equation \eqref{eq:rouche-root-count} gives only
\(o(s_\nu)\) interior roots, and every term in \(F_{q,\nu}\) contains
at least one of them. The exterior roots are not discarded: their
number can be \(O(s_\nu)\), but their distance from \(D_6\) is bounded
below, and the corresponding terms form the holomorphic part
\(B_{q,\nu}\).


\subsection{A local representation of the curve}
\label{sec:representation}

We next show that the hypotheses of
Proposition~\ref{prop:localcompact} arise after rescaling the original
curve. The scalar gauge is chosen so that the local Wronskian becomes a
polynomial with \(o(s_\nu)\) zeros. We also need control of the value
of the gauge at the origin and of the negative part of
\(\log\norm{\gv_\nu}\).

\begin{prop}\label{prop:representation}
Let \(\fv\) be a curve as defined in Section~\ref{sec:intro}, and make
the fixed unitary coordinate change for which
\(f_j(0)\ne0\), \(0\leq j\leq n\).
Suppose that \(t_\nu,s_\nu\to\infty\) and, for a fixed \(C\),
\begin{equation}\label{eq:scalehyp}
 \log t_\nu=o(s_\nu),
 \qquad
 T(256t_\nu,\fv)\leq Cs_\nu,
 \qquad
 \NN_{W_{\fv}}(256t_\nu)=o(s_\nu).
\end{equation}
After omitting finitely many terms, there are holomorphic vectors
\(\gv_\nu\) on \(D_{64}\) and monic polynomials \(P_\nu\) such that
\begin{equation}\label{eq:representationbounds}
 \begin{gathered}
 W(\gv_\nu)=P_\nu,
 \qquad
 \deg P_\nu=o(s_\nu),\\
 \sup_{D_{32}}\log\norm{\gv_\nu}
 \leq(5C+1)s_\nu,\\
 s_\nu^{-1}\log|g_{\nu,j}(0)|
 \longrightarrow0,
 \qquad
 0\leq j\leq n.
 \end{gathered}
\end{equation}
For \(0<R<64\),
\begin{equation}\label{eq:meanidentity}
 \mean R{\log\norm{\gv_\nu}}
 =
 T(Rt_\nu,\fv)+c_\nu,
 \qquad
 c_\nu=o(s_\nu),
\end{equation}
where \(c_\nu\) is independent of \(R\).

The functions
\begin{equation}\label{eq:scaledcompact}
 \left(\frac{t_\nu}{s_\nu}\right)^qQ_q(t_\nu z),
 \qquad
 2\leq q\leq n+1,
\end{equation}
are simultaneously relatively compact locally in measure on \(D_4\).
Every subsequential limit is holomorphic, and its restriction to
\(D_1\) is bounded by a constant depending only on \({n+1}\) and \(C\).
\end{prop}

\begin{proof}

\medskip\noindent
\textbf{\underline{Step 1. Normalization of the local Wronskian divisor:}}
Recall that
\(\kappa=n(n+1)/2\), and write
\(W_{\fv}(z)=w_0z^{m_0}+O(z^{m_0+1})\), \(w_0\ne0\).

Let \(P_\nu\) be the monic polynomial whose zeros are the points
\(a/t_\nu\), where \(a\) runs through the zeros of \(W_{\fv}\) in
\(D_{64t_\nu}\), with the same multiplicities. Put
\(m_\nu=\deg P_\nu\). By \eqref{eq:countbound},
\[
 m_\nu\log2
 \leq
 \NN_{W_{\fv}}(128t_\nu)
 \leq
 \NN_{W_{\fv}}(256t_\nu).
\]
The last term is \(o(s_\nu)\) by \eqref{eq:scalehyp}; hence
\begin{equation}\label{eq:local-P-degree}
 m_\nu=o(s_\nu).
\end{equation}

The function
\(t_\nu^\kappa W_{\fv}(t_\nu z)/P_\nu(z)\)
is holomorphic and zero-free on \(D_{64}\). Since \(D_{64}\) is simply
connected, there is a holomorphic \(H_\nu\) on \(D_{64}\) such that
\begin{equation}\label{eq:localgauge}
 e^{(n+1)H_\nu(z)}
 =
 \frac{t_\nu^\kappa W_{\fv}(t_\nu z)}{P_\nu(z)},
 \qquad
 \gv_\nu(z)
 =
 e^{-H_\nu(z)}\fv(t_\nu z).
\end{equation}
The Wronskian multiplication and scaling formulas applied to
\eqref{eq:localgauge} give
\begin{equation}\label{eq:local-wronskian-normalization}
 W(\gv_\nu)=P_\nu.
\end{equation}


\medskip\noindent
\textbf{\underline{Step 2. The value of the gauge at the origin:}}
The zero of \(P_\nu\) at the origin has multiplicity \(m_0\). For every
nonzero zero \(a\) of \(W_{\fv}\) with \(|a|<64t_\nu\), the
corresponding factor of \(P_\nu\) contributes
\(\log(t_\nu/|a|)\) to the logarithm of the quotient in
\eqref{eq:localgauge} at the origin. Jensen's formula
\eqref{eq:zerojensen} therefore gives
\begin{equation}\label{eq:centeridentity}
 {n+1}\Ree H_\nu(0)
 =
 \kappa\log t_\nu
 +
 \NN_{W_{\fv}}(64t_\nu)
 -
 m_\nu\log64
 +
 \log|w_0|.
\end{equation}
By \eqref{eq:scalehyp} and \eqref{eq:local-P-degree}, every term on
the right-hand side of \eqref{eq:centeridentity}, except the fixed
constant \(\log|w_0|\), is \(o(s_\nu)\). Hence
\begin{equation}\label{eq:H-center-small}
 \Ree H_\nu(0)=o(s_\nu).
\end{equation}

Since \(f_j(0)\ne0\) is fixed,
\(\log|g_{\nu,j}(0)|
 =\log|f_j(0)|-\Ree H_\nu(0)\).
Equation \eqref{eq:H-center-small} therefore gives
\begin{equation}\label{eq:center-component-normalization}
 s_\nu^{-1}\log|g_{\nu,j}(0)|
 \longrightarrow0,
 \qquad
 0\leq j\leq n.
\end{equation}

Taking logarithms of the norm in \eqref{eq:localgauge} gives
\(\log\norm{\gv_\nu(z)}
 =\log\norm{\fv(t_\nu z)}-\Ree H_\nu(z)\).
Since \(\Ree H_\nu\) is harmonic, its circular mean equals
\(\Ree H_\nu(0)\). By the definition of the characteristic,
\[
 \mean R{\log\norm{\fv(t_\nu\cdot)}}
 =
 T(Rt_\nu,\fv)+\log\norm{\fv(0)}.
\]
Consequently \eqref{eq:meanidentity} holds with
\begin{equation}\label{eq:cnu-definition}
 c_\nu
 :=
 \log\norm{\fv(0)}-\Ree H_\nu(0)
 =
 o(s_\nu),
\end{equation}
where the last equality follows from \eqref{eq:H-center-small}; thus
the constant in \eqref{eq:meanidentity} is the \(c_\nu\) in
\eqref{eq:cnu-definition}.


\medskip\noindent
\textbf{\underline{Step 3. Control of the negative part of the norm:}}
For \(0\leq j\leq n\), let \(E_{\nu,j}\) be the set of points at
which
\(|g_{\nu,j}|=\max_\ell|g_{\nu,\ell}|\); ties may be resolved by taking
the smallest index. These sets cover \(D_{64}\) up to a set of area
zero.

Fix \(j\). On \(E_{\nu,j}\), put
\[
 h_\ell
 :=
 \frac{g_{\nu,\ell}}{g_{\nu,j}}
 =
 \frac{f_\ell(t_\nu z)}{f_j(t_\nu z)}.
\]
Then \(h_j=1\) and \(|h_\ell|\leq1\) on \(E_{\nu,j}\).
The Wronskian multiplication formula gives
\(P_\nu=g_{\nu,j}^{n+1}W(h_0,\ldots,h_{n})\).
Expanding the latter Wronskian along the column \(h_j=1\), and then
factoring \(h_\ell\) from the remaining columns, gives
\begin{equation}\label{eq:negativedeterminant}
 \frac{|P_\nu|}{\norm{\gv_\nu}^{n+1}}
 \leq
 \left|
 \det\left(
 \frac{h_\ell^{(i)}}{h_\ell}
 \right)_{
 1\leq i\leq n,\ \ell\ne j}
 \right|
 \qquad\text{on }E_{\nu,j}.
\end{equation}

For fixed \(j\ne\ell\), apply
Lemma~\ref{lem:logderivbound} to the meromorphic function
\(z\mapsto f_\ell(48t_\nu z)/f_j(48t_\nu z)\).
Its value at the origin is fixed and nonzero. By
\eqref{eq:quotientbound} and \eqref{eq:scalehyp},
\[
 \TN\left(
 4,
 \frac{f_\ell(48t_\nu z)}{f_j(48t_\nu z)}
 \right)
 \leq
 T(192t_\nu,\fv)+O(1)
 \leq
 Cs_\nu+O(1).
\]
After rescaling the derivatives by the fixed factor \(48\),
Lemma~\ref{lem:logderivbound} gives, for
\(1\leq i\leq n\),
\begin{equation}\label{eq:ratio-logderiv-48}
 m\left(
 48,\frac{h_\ell^{(i)}}{h_\ell}
 \right)
 \leq
 C_{n+1,C}(\log s_\nu+1)
 =
 o(s_\nu).
\end{equation}

Expanding the determinant in \eqref{eq:negativedeterminant} into its
finitely many permutation terms, and applying
\eqref{eq:ratio-logderiv-48} to every factor, gives
\begin{equation}\label{eq:determinant-proximity}
 \mean{48}{
 \log^+
 \left|
 \det\left(
 \frac{h_\ell^{(i)}}{h_\ell}
 \right)_{
 1\leq i\leq n,\ \ell\ne j}
 \right|
 }
 =
 o(s_\nu),
\end{equation}
uniformly in \(j\).

All roots of \(P_\nu\) lie in \(D_{64}\), and
\(\deg P_\nu=m_\nu\). Since
\(\sup_{a\in\C}\mean{48}{\log^-|z-a|}<\infty\),
subadditivity of \(\log^-\) gives
\begin{equation}\label{eq:P-negative-circle}
 \mean{48}{\log^-|P_\nu|}
 \leq
 C_4m_\nu
 =
 o(s_\nu),
\end{equation}
where \eqref{eq:local-P-degree} is used in the last equality.

Taking logarithms in \eqref{eq:negativedeterminant} gives on
\(E_{\nu,j}\)
\[
 {n+1}\log^-\norm{\gv_\nu}
 \leq
 \log^-|P_\nu|
 +
 \log^+
 \left|
 \det\left(
 \frac{h_\ell^{(i)}}{h_\ell}
 \right)
 \right|.
\]
Integrating this inequality over the sets \(E_{\nu,j}\cap\{|z|=48\}\),
summing over \(j\), and using
\eqref{eq:determinant-proximity} and
\eqref{eq:P-negative-circle}, we obtain
\begin{equation}\label{eq:negativepart}
 \mean{48}{\log^-\norm{\gv_\nu}}
 =
 o(s_\nu).
\end{equation}


\medskip\noindent
\textbf{\underline{Step 4. Uniform upper bound for the local frame:}}
At \(R=48\), equation \eqref{eq:meanidentity} and
\eqref{eq:scalehyp} give
\(\mean{48}{\log\norm{\gv_\nu}}\leq Cs_\nu+o(s_\nu)\).
Since
\(\log x=\log^+x-\log^-x\),
equation \eqref{eq:negativepart} yields
\begin{equation}\label{eq:positivepart}
 \mean{48}{\log^+\norm{\gv_\nu}}
 \leq
 Cs_\nu+o(s_\nu).
\end{equation}

The function \(\log\norm{\gv_\nu}\) is subharmonic. For
\(|z|\leq32\), the Poisson kernel of \(D_{48}\) is bounded by
\((48+32)/(48-32)=5\). Hence \eqref{eq:positivepart} gives
\[
 \sup_{D_{32}}\log\norm{\gv_\nu}
 \leq
 5Cs_\nu+o(s_\nu)
 \leq
 (5C+1)s_\nu
\]
for all sufficiently large \(\nu\). Together with
\eqref{eq:local-P-degree},
\eqref{eq:local-wronskian-normalization}, and
\eqref{eq:center-component-normalization}, this proves
\eqref{eq:representationbounds}.


\medskip\noindent
\textbf{\underline{Step 5. Compactness of the rescaled canonical coefficients:}}
By \eqref{eq:localgauge}, after multiplication by a constant
\((n+1)\)st root of unity,
\[
 P_\nu^{-1/(n+1)}\gv_\nu
 =
 t_\nu^{-\kappa/{n+1}}
 \bigl(W_{\fv}^{-1/(n+1)}\fv\bigr)(t_\nu z).
\]
Multiplication of all solutions by a nonzero constant does not change
their fundamental operator.

By \eqref{eq:representationbounds},
Proposition~\ref{prop:localcompact} applies to \(\gv_\nu\). Let
\(\widehat b_{q,\nu}\) be the coefficients of the fundamental operator
of \(P_\nu^{-1/(n+1)}\gv_\nu\). The scaling rule
\eqref{eq:scalecoeff} then gives
\begin{equation}\label{eq:canonical-coefficient-identification}
 \widehat b_{q,\nu}(z)
 =
 t_\nu^qQ_q(t_\nu z),
 \qquad
 2\leq q\leq n+1.
\end{equation}
Therefore
\[
 \frac{\widehat b_{q,\nu}}{s_\nu^q}
 =
 \left(\frac{t_\nu}{s_\nu}\right)^qQ_q(t_\nu z).
\]
The compactness and uniform bounds in
Proposition~\ref{prop:localcompact} now give exactly the assertion
\eqref{eq:scaledcompact}.
\end{proof}

The bound in the last assertion depends only on \({n+1}\) and \(C\), and
is uniform over all sequences satisfying \eqref{eq:scalehyp}. This
uniformity will later allow us to replace \(t_\nu\) by an arbitrary
fixed multiple of \(t_\nu\).


\subsection{Polynomial data retained from the approximation}
\label{sec:replacement}

Later we shall use the polynomial spaces constructed in the proof of
Proposition~\ref{prop:localcompact}, not only their differential
operators. The following lemma collects the properties needed later.

\begin{lem}\label{lem:replacement}
Under \eqref{eq:scalehyp}, let \(\gv_\nu\) be the vectors constructed
in Proposition~\ref{prop:representation}. They admit polynomial
approximants \(\pv_\nu\) of degree \(O(s_\nu)\) satisfying
\eqref{eq:taylorerror}, for any fixed sufficiently large \(A\).

If \(R_\nu=W(\pv_\nu)\), then there are
\(\eta_\nu\in(8,10)\) such that
\begin{equation}\label{eq:replacement-data}
 \begin{gathered}
 \deg R_\nu=O(s_\nu),
 \qquad
 n_{R_\nu}(\eta_\nu)=o(s_\nu),\\
 s_\nu^{-1}\log|R_\nu|
 \convmeas0
 \qquad\text{locally on }D_4.
 \end{gathered}
\end{equation}
If \(b_{q,\nu}^p\) are the fundamental coefficients of \(\pv_\nu\),
then
\begin{equation}\label{eq:polynomial-coefficients}
 \frac{b_{1,\nu}^p}{s_\nu}
 \convmeas0,
 \qquad
 \frac{b_{q,\nu}^p}{s_\nu^q}
 -
 \left(\frac{t_\nu}{s_\nu}\right)^qQ_q(t_\nu z)
 \convmeas0,
 \qquad
 2\leq q\leq n+1,
\end{equation}
locally on \(D_4\).
\end{lem}

\begin{proof}

\medskip\noindent
\textbf{\underline{Step 1. The normalized logarithm of \(P_\nu\):}}
All roots of \(P_\nu\) lie in \(D_{64}\), and
\(\deg P_\nu=o(s_\nu)\) by \eqref{eq:local-P-degree}.
Fix \(R>0\). Since \(P_\nu\) is monic,
\(\log^+|P_\nu(z)|
 \leq(\deg P_\nu)\log(R+64)\)
for \(|z|\leq R\). Moreover,
\eqref{eq:polynomial-negative-area}, with \(D_R\) in place of
\(D_{12}\), gives
\(\int_{D_R}\log^-|P_\nu|\,\dA
 \leq C_R\deg P_\nu\).
Consequently
\begin{equation}\label{eq:small-polynomial-log}
 s_\nu^{-1}\log|P_\nu|
 \longrightarrow0
 \qquad\text{in }L^1_{\loc}(\C).
\end{equation}


\medskip\noindent
\textbf{\underline{Step 2. Transfer from \(P_\nu\) to \(R_\nu\):}}
Equation \eqref{eq:wronskiapprox} gives
\[
 \sup_{D_{12}}|R_\nu-P_\nu|
 \leq e^{-Bs_\nu},
\]
where \(B>0\) is fixed.

Choose \(0<\epsilon<B\), and define
\(E_{\nu,\epsilon}
 :=\{z\in D_{12}:|P_\nu(z)|<e^{-\epsilon s_\nu}\}\).
Since
\(\log^-|P_\nu|>\epsilon s_\nu\) on this set,
\eqref{eq:polynomial-negative-area} and
\eqref{eq:local-P-degree} give
\begin{equation}\label{eq:replacement-exceptional-set}
 \meas(E_{\nu,\epsilon})
 \leq
 \frac{C\,\deg P_\nu}{\epsilon s_\nu}
 \longrightarrow0.
\end{equation}

On \(D_{12}\setminus E_{\nu,\epsilon}\),
\[
 \left|\frac{R_\nu}{P_\nu}-1\right|
 \leq
 e^{-(B-\epsilon)s_\nu}.
\]
It follows that
\(s_\nu^{-1}\log|R_\nu/P_\nu|\to0\) uniformly on
\(D_{12}\setminus E_{\nu,\epsilon}\). Equation
\eqref{eq:replacement-exceptional-set} and the convergence in
\eqref{eq:small-polynomial-log} therefore imply
\begin{equation}\label{eq:R-log-measure}
 s_\nu^{-1}\log|R_\nu|
 \convmeas0
 \qquad\text{locally on }D_4.
\end{equation}

The functions \(s_\nu^{-1}\log|R_\nu|\) are subharmonic and locally
uniformly bounded above: this follows directly from
\eqref{eq:wronskiapprox} and the pointwise upper bound for
\(\log^+|P_\nu|\) used in Step~1. Hence
Lemma~\ref{lem:subharmonic-compactness} and
\eqref{eq:R-log-measure} give
\[
 s_\nu^{-1}\log|R_\nu|
 \longrightarrow0
 \qquad\text{in }L^1_{\loc}(D_4).
\]

The degree estimate in \eqref{eq:replacement-data} is the first
assertion of \eqref{eq:wronskiapprox}, while the zero-counting estimate
is exactly \eqref{eq:rouche-root-count}. Thus all assertions in
\eqref{eq:replacement-data} are proved.


\medskip\noindent
\textbf{\underline{Step 3. Comparison of the differential coefficients:}}
Let \(b_{q,\nu}\) be the fundamental coefficients of \(\gv_\nu\), and
let \(\widehat b_{q,\nu}\) be those of
\(P_\nu^{-1/(n+1)}\gv_\nu\).

Equation \eqref{eq:cramercomparison} gives
\[
 \frac{b_{q,\nu}^p-b_{q,\nu}}{s_\nu^q}
 \convmeas0,
\]
and \eqref{eq:gaugecomparison} gives
\[
 \frac{b_{q,\nu}-\widehat b_{q,\nu}}{s_\nu^q}
 \convmeas0.
\]
For \(q=1\), Proposition~\ref{prop:localcompact} gives
\(\widehat b_{1,\nu}=0\). Hence
\(b_{1,\nu}^p/s_\nu\convmeas0\).

For \(2\leq q\leq n+1\), equation
\eqref{eq:canonical-coefficient-identification} gives
\[
 \frac{\widehat b_{q,\nu}}{s_\nu^q}
 =
 \left(\frac{t_\nu}{s_\nu}\right)^qQ_q(t_\nu z).
\]
Therefore
\[
 \frac{b_{q,\nu}^p}{s_\nu^q}
 -
 \left(\frac{t_\nu}{s_\nu}\right)^qQ_q(t_\nu z)
 =
 \frac{b_{q,\nu}^p-b_{q,\nu}}{s_\nu^q}
 +
 \frac{b_{q,\nu}-\widehat b_{q,\nu}}{s_\nu^q},
\]
and each term on the right tends to zero in measure by
\eqref{eq:cramercomparison} and \eqref{eq:gaugecomparison}. This is
\eqref{eq:polynomial-coefficients}.
\end{proof}
\section{Growth indices}\label{sec:indices}

To pass from selected radii to arbitrary radii, we use the full interval
of P\'olya peak orders. This is more convenient than working only with
the lower and upper orders.

\subsection{Drasin--Shea indices and P\'olya peaks}
\label{sec:polya-peaks}

Let \(T\) be a positive, unbounded, nondecreasing function. Its
Drasin--Shea indices \cite{DS72} are
\begin{equation}\label{eq:strong-indices}
\begin{split}
 \rho^*(T)
 &=
 \sup\left\{
 p\in\R:
 \limsup_{x,t\to\infty}
 \frac{T(xt)}{x^pT(t)}
 =
 \infty
 \right\},\\
 \rho_*(T)
 &=
 \inf\left\{
 p\in\R:
 \liminf_{x,t\to\infty}
 \frac{T(xt)}{x^pT(t)}
 =
 0
 \right\}.
\end{split}
\end{equation}
The limits in \eqref{eq:strong-indices} are joint limits as
\(x,t\to\infty\). For the characteristic of a holomorphic curve,
\[
0\leq
\rho_*(T)
\leq
\lambda(\fv)
\leq
\rho(\fv)
\leq
\rho^*(T).
\]

For later use, we need a form of the P\'olya-peak theorem that is
uniform for all fixed multipliers \(t>0\) and includes the endpoint
orders. The endpoint assertion is contained in the sentence following
\cite[p.~1203, equations~(16)--(19)]{Ere93}; see also
\cite[Section~2.2, equations~(17)--(19)]{Ere98}.

\begin{lem}[{\cite[p.~1203, equations~(16)--(19)]{Ere93};
\cite[Section~2.2, equations~(17)--(19)]{Ere98}}]
\label{lem:peaks}
Let \(T\) be a positive, unbounded, nondecreasing function. For every
finite \(\mu>0\) belonging to
\([\rho_*(T),\rho^*(T)]\) and every
\(0<\delta<\min\{\mu,1\}\), there exist
\(r_\nu\to\infty\) and a constant \(C_{T,\mu,\delta}>0\) such that,
with \(S_\nu=T(r_\nu)\),
\begin{equation}\label{eq:peak}
 \limsup_{\nu\to\infty}
 \frac{T(tr_\nu)}{S_\nu}
 \leq
 C_{T,\mu,\delta}
 \max\{t^{\mu-\delta},t^{\mu+\delta}\},
 \qquad t>0.
\end{equation}
\end{lem}

The P\'olya-peak theorem in the cited references gives a sequence
\(r_\nu\to\infty\), numbers \(\varepsilon_\nu\downarrow0\), and a fixed
\(\varepsilon_0\in(0,1)\) such that, for all sufficiently large \(\nu\),
\[
 \frac{T(tr_\nu)}{T(r_\nu)}
 \leq t^{\mu-\delta}
 \quad\text{whenever}\quad
 \varepsilon_\nu\leq t\leq\varepsilon_0,
\]
and
\[
 \frac{T(tr_\nu)}{T(r_\nu)}
 \leq t^{\mu+\delta}
 \quad\text{whenever}\quad
 \varepsilon_0^{-1}\leq t\leq\varepsilon_\nu^{-1}.
\]
The endpoint values of \(\mu\) are included in the P\'olya-peak
criterion stated after \cite[(19), p.~1203]{Ere93}.

For \(0<t\leq\varepsilon_0\), the first estimate applies for all
sufficiently large \(\nu\), because \(\varepsilon_\nu\to0\). For
\(t\geq\varepsilon_0^{-1}\), the second estimate applies. If
\(\varepsilon_0<t<1\), monotonicity gives
\[
 \frac{T(tr_\nu)}{T(r_\nu)}\leq1
 \leq
 \varepsilon_0^{-(\mu-\delta)}t^{\mu-\delta}.
\]
If \(1<t<\varepsilon_0^{-1}\), monotonicity and the second estimate at
the endpoint give
\[
 \frac{T(tr_\nu)}{T(r_\nu)}
 \leq
 \varepsilon_0^{-(\mu+\delta)}
 \leq
 \varepsilon_0^{-(\mu+\delta)}t^{\mu+\delta}.
\]
Thus \eqref{eq:peak} holds with
\[
 C_{T,\mu,\delta}
 =
 \max\{1,\varepsilon_0^{-(\mu-\delta)},
           \varepsilon_0^{-(\mu+\delta)}\}.
\]


\subsection{Collapse of the interval of growth indices}
\label{sec:index-collapse}

We now determine the interval under the small-ramification hypothesis:
$
[\rho_*(T),\rho^*(T)].
$

\begin{prop}\label{prop:indices}
Suppose that
\[
\lambda(\fv)<\infty,
\qquad
\NN_{W_{\fv}}(r)=o(T(r,\fv)),
\qquad
\log r=o(T(r,\fv)).
\]
Then, with \(T(r)=T(r,\fv)\),
\begin{equation}\label{eq:indices-collapse}
 \rho_*(T)
 =
 \lambda(\fv)
 =
 \rho(\fv)
 =
 \rho^*(T)
 \in
 \{0\}\cup\Gn.
\end{equation}
\end{prop}

\begin{proof}

If \(\rho^*(T)=0\), then
\[
0\leq
\rho_*(T)
\leq
\lambda(\fv)
\leq
\rho(\fv)
\leq
\rho^*(T)=0,
\]
and \eqref{eq:indices-collapse} follows. We therefore assume
\(\rho^*(T)>0\).

Fix
$
0<\mu\in[\rho_*(T),\rho^*(T)].
$
We shall prove that \(\mu\in\Gn\).
Assume, to the contrary, that \(\mu\notin\Gn\). Since
\(2\leq q\leq n+1\) takes only finitely many values, we may choose
\(0<\delta<\min\{\mu,1\}\) so small that
\begin{equation}\label{eq:no-integer-window}
 [q(\mu-1-\delta),q(\mu-1+\delta)]
 \cap\Z_{\geq0}
 =
 \varnothing,
 \qquad
 2\leq q\leq n+1.
\end{equation}
Let \(r_\nu\) be the sequence supplied by
Lemma~\ref{lem:peaks}, put \(S_\nu=T(r_\nu)\), and abbreviate
\(M(t)=M_{\mu,\delta}(t)\).


\medskip\noindent
\textbf{\underline{Step 1. The local compactness theorem applies at every fixed scale:}}
Fix \(R>0\). We apply Proposition~\ref{prop:representation} with
the scale \(Rr_\nu\) and the normalization \(M(R)S_\nu\).

First, since \(\log r=o(T(r))\),
\[
\frac{\log(Rr_\nu)}{M(R)S_\nu}\longrightarrow0.
\]
Next, \eqref{eq:peak} at the fixed multiplier \(256R\) gives, for all
sufficiently large \(\nu\),
\[
T(256Rr_\nu)
\leq
2C_{\mu,\delta}M(256R)S_\nu.
\]
Since
\(M(256R)\leq256^{\mu+\delta}M(R)\), we obtain
\begin{equation}\label{eq:peak-characteristic-scale}
 \frac{T(256Rr_\nu)}{M(R)S_\nu}
 \leq
 C_0,
 \qquad
 C_0:=2C_{\mu,\delta}256^{\mu+\delta}.
\end{equation}
The constant \(C_0\) is independent of \(R\).

Finally,
\[
\frac{\NN_{W_{\fv}}(256Rr_\nu)}{M(R)S_\nu}
=
\frac{\NN_{W_{\fv}}(256Rr_\nu)}
     {T(256Rr_\nu)}
\,
\frac{T(256Rr_\nu)}
     {M(R)S_\nu}.
\]
The first factor tends to zero by the small-ramification hypothesis,
while the second is bounded by \eqref{eq:peak-characteristic-scale}.
Hence
\begin{equation}\label{eq:peak-representation-hypotheses}
 \frac{\log(Rr_\nu)}{M(R)S_\nu}\longrightarrow0,
 \qquad
 \frac{\NN_{W_{\fv}}(256Rr_\nu)}
      {M(R)S_\nu}
 \longrightarrow0.
\end{equation}
Thus \eqref{eq:peak-representation-hypotheses} verifies all hypotheses
of Proposition~\ref{prop:representation},
with a characteristic bound independent of \(R\).


\medskip\noindent
\textbf{\underline{Step 2. Limits of the rescaled differential coefficients:}}
For \(2\leq q\leq n+1\), define
\begin{equation}\label{eq:peak-scaled-coefficients}
 A_{q,\nu}(z)
 :=
 \left(\frac{r_\nu}{S_\nu}\right)^q
 Q_q(r_\nu z).
\end{equation}

For fixed \(R>0\), Proposition~\ref{prop:representation}, applied as in
Step~1, gives relative compactness locally in measure of
\[
\left(
\frac{Rr_\nu}{M(R)S_\nu}
\right)^q
Q_q(Rr_\nu z)
\]
on \(D_4\). The relation with \eqref{eq:peak-scaled-coefficients} is
\begin{equation}\label{eq:peak-rescaling-relation}
 A_{q,\nu}(Rz)
 =
 \left(\frac{M(R)}{R}\right)^q
 \left(
 \frac{Rr_\nu}{M(R)S_\nu}
 \right)^q
 Q_q(Rr_\nu z).
\end{equation}

Apply Proposition~\ref{prop:representation} successively for
\(R=1,2,3,\ldots\), and take a diagonal subsequence. Then, for
\(2\leq q\leq n+1\),
\[
A_{q,\nu}\longrightarrow a_q
\]
locally in measure on \(\C\), where every \(a_q\) is entire.
The limits obtained on different disks agree because limits in measure
are unique.

The uniform bound in Proposition~\ref{prop:representation}, together
with \eqref{eq:peak-rescaling-relation}, gives
\begin{equation}\label{eq:peak-coefficient-bound}
 \sup_{|z|\leq R}|a_q(z)|
 \leq
 K_{n+1,\mu,\delta}
 \max\{
 R^{q(\mu-1-\delta)},
 R^{q(\mu-1+\delta)}
 \},
 \qquad
 R>0.
\end{equation}
Here \(K_{n+1,\mu,\delta}\) is independent of \(R\).


\medskip\noindent
\textbf{\underline{Step 3. The limiting differential operator is nontrivial:}}
We claim that
\[
(a_2,\ldots,a_{n+1})\not\equiv(0,\ldots,0).
\]
Suppose otherwise.

Apply Proposition~\ref{prop:representation} with
\(t_\nu=r_\nu\) and \(s_\nu=S_\nu\). This is the case \(R=1\) of
Step~1. Let
\(\gv_\nu=(g_{\nu,0},\ldots,g_{\nu,n})\)
be the resulting local frames, and put
\(u_{\nu,j}=S_\nu^{-1}\log|g_{\nu,j}|\).

By \eqref{eq:representationbounds}, the functions \(u_{\nu,j}\) are
locally uniformly bounded above on \(D_{32}\), and
\(u_{\nu,j}(0)\to0\). Lemma~\ref{lem:subharmonic-compactness} therefore
allows us, after passing to a further subsequence, to assume
\begin{equation}\label{eq:peak-component-limits}
 u_{\nu,j}
 \longrightarrow
 u_j
 \qquad\text{in }L^1_{\loc}(D_4),
 \qquad
 0\leq j\leq n,
\end{equation}
where no \(u_j\) is identically \(-\infty\).

Let \(\widehat b_{q,\nu}\) be the coefficients after the canonical
gauge. By the scaling rule used in
Proposition~\ref{prop:representation},
\begin{equation}\label{eq:peak-canonical-coefficients}
 \widehat b_{1,\nu}=0,
 \qquad
 \frac{\widehat b_{q,\nu}}{S_\nu^q}
 =
 A_{q,\nu},
 \qquad
 2\leq q\leq n+1.
\end{equation}
Under the assumption \(a_q\equiv0\), equations
\eqref{eq:peak-canonical-coefficients} and
\eqref{eq:gaugecomparison} give
\begin{equation}\label{eq:peak-unnormalized-coefficients}
 \frac{b_{q,\nu}}{S_\nu^q}
 \convmeas0
 \qquad\text{locally on }D_4,
 \qquad
 1\leq q\leq n+1.
\end{equation}

Each \(g_{\nu,j}\) satisfies the differential equation with coefficients
\(b_{q,\nu}\). Away from the zeros of \(g_{\nu,j}\), divide this
equation by \(S_\nu^{n+1}g_{\nu,j}\). We obtain
\begin{equation}\label{eq:normalized-peak-equation}
 \frac{g_{\nu,j}^{(n+1)}}{S_\nu^{n+1}g_{\nu,j}}
 +
 \sum_{q=1}^{n+1}
 \frac{b_{q,\nu}}{S_\nu^q}
 \frac{g_{\nu,j}^{({n+1}-q)}}
      {S_\nu^{{n+1}-q}g_{\nu,j}}
 =
 0.
\end{equation}
By Lemma~\ref{lem:logderivlimit} and
\eqref{eq:peak-component-limits},
the first term in \eqref{eq:normalized-peak-equation} converges in
measure to \((2\partial u_j)^{n+1}\), while each logarithmic-derivative
factor in the sum is bounded in measure. Equation
\eqref{eq:peak-unnormalized-coefficients} therefore makes every term
in the sum tend to zero in measure. Hence
\[
(2\partial u_j)^{n+1}=0
\qquad\text{a.e.\ on }D_4.
\]
By the Sobolev conclusion of Lemma~\ref{lem:logderivlimit},
\(\partial u_j=0\) almost everywhere, and thus every \(u_j\) is
constant on \(D_4\).

Set
\(U_\nu=S_\nu^{-1}\log\norm{\gv_\nu}\).
Since
\[
0\leq
S_\nu^{-1}\log\norm{\gv_\nu}
-
\max_j u_{\nu,j}
\leq
\frac{\log {n+1}}{2S_\nu},
\]
equation \eqref{eq:peak-component-limits} implies
\begin{equation}\label{eq:constant-norm-limit}
 U_\nu\longrightarrow c
 \qquad\text{in }L^1_{\loc}(D_4)
\end{equation}
for some constant \(c\).

By \eqref{eq:meanidentity},
\[
\mean t{U_\nu}
=
\frac{T(tr_\nu)}{S_\nu}+o(1),
\qquad
0<t<4,
\]
where the error is independent of \(t\). Put
\(F_\nu(t)=T(tr_\nu)/S_\nu\). Integrating over an annulus and using
\eqref{eq:constant-norm-limit} gives
\begin{equation}\label{eq:annular-characteristic-limit}
 \frac{2}{b^2-a^2}
 \int_a^b F_\nu(t)t\,\dd t
 \longrightarrow c,
 \qquad
 0<a<b<4.
\end{equation}

Take \(0<b<1\) and \(a=b/2\). Since \(F_\nu\) is nondecreasing,
\eqref{eq:annular-characteristic-limit} gives
\(c\leq\limsup_\nu F_\nu(b)\). By \eqref{eq:peak},
\[
c
\leq
C_{\mu,\delta}M(b)
=
C_{\mu,\delta}b^{\mu-\delta}.
\]
Since \(\mu-\delta>0\), letting \(b\downarrow0\) gives \(c\leq0\).

On the other hand, if \(2<a<b<3\), then
\(F_\nu(t)\geq F_\nu(1)=1\) for \(a<t<b\). Equation
\eqref{eq:annular-characteristic-limit} therefore gives \(c\geq1\).
This contradiction proves that at least one \(a_q\) is nonzero.


\medskip\noindent
\textbf{\underline{Step 4. The peak order belongs to \(\Gn\):}}
Write
\(a_q(z)=\sum_{m=0}^{\infty}c_{q,m}z^m\), and set
\[
\alpha_{q,-}=q(\mu-1-\delta),
\qquad
\alpha_{q,+}=q(\mu-1+\delta).
\]
Cauchy's estimate and \eqref{eq:peak-coefficient-bound} give
\begin{equation}\label{eq:peak-Taylor-coefficient-bound}
 |c_{q,m}|
 \leq
 \begin{cases}
 K_{n+1,\mu,\delta}R^{\alpha_{q,-}-m},
     &0<R\leq1,\\[2mm]
 K_{n+1,\mu,\delta}R^{\alpha_{q,+}-m},
     &R\geq1.
 \end{cases}
\end{equation}

If \(m<\alpha_{q,-}\), let \(R\downarrow0\) in the first inequality
of \eqref{eq:peak-Taylor-coefficient-bound}; then \(c_{q,m}=0\).
If \(m>\alpha_{q,+}\), let \(R\to\infty\) in the second inequality;
again \(c_{q,m}=0\).

By \eqref{eq:no-integer-window}, every nonnegative integer \(m\) lies
outside
\([\alpha_{q,-},\alpha_{q,+}]\). Hence every Taylor coefficient of
every \(a_q\) vanishes. Thus \(a_q\equiv0\) for all \(q\), contradicting
Step~3.

The contradiction proves
\begin{equation}\label{eq:positive-index-in-Rn}
 [\rho_*(T),\rho^*(T)]
 \cap(0,\infty)
 \subset\Gn.
\end{equation}


\medskip\noindent
\textbf{\underline{Step 5. Collapse of the index interval:}}
The set \(\Gn\) is locally finite. Indeed, on every bounded interval,
\(2\leq q\leq n+1\) gives only finitely many possible integers \(k\) in
\(1+k/q\).

The set
\([\rho_*(T),\rho^*(T)]\cap(0,\infty)\), however, is an interval.
If this interval contained two distinct finite points, then a nonempty
open subinterval would be contained in the locally finite set \(\Gn\) by
\eqref{eq:positive-index-in-Rn}, which is impossible. The same argument
rules out \(\rho^*(T)=+\infty\), because in that case the interval would
contain two distinct finite points. Hence
\([\rho_*(T),\rho^*(T)]\) is a singleton. Its value is finite, since
\[
\rho_*(T)\leq\lambda(\fv)<\infty.
\]

Finally,
\[
\rho_*(T)
\leq
\lambda(\fv)
\leq
\rho(\fv)
\leq
\rho^*(T),
\]
so all four quantities are equal. This proves
\eqref{eq:indices-collapse}.
\end{proof}


\subsection{Two-sided power bounds}
\label{sec:power-bounds}

The equality of the two Drasin--Shea indices gives uniform control of
\(T(tr)/T(r)\) at every fixed scale. Regular variation will require an
additional argument in Section~\ref{sec:regular}.

\begin{lem}\label{lem:power-bounds}
Suppose that
$
\rho_*(T)=\rho^*(T)=\rho\in(0,\infty).
$
For every \(0<\epsilon<\rho\), there are
\(C_\epsilon\geq1\) and \(r_\epsilon>0\) such that
\begin{equation}\label{eq:power-bounds}
 C_\epsilon^{-1}
 \min\{t^{\rho-\epsilon},t^{\rho+\epsilon}\}
 \leq
 \frac{T(tr)}{T(r)}
 \leq
 C_\epsilon
 \max\{t^{\rho-\epsilon},t^{\rho+\epsilon}\},
 \qquad
 r,tr\geq r_\epsilon.
\end{equation}
\end{lem}

\begin{proof}

Put
$
p_-=\rho-\epsilon,
$ and $
p_+=\rho+\epsilon.
$
Since \(p_+>\rho^*(T)\), the definition of \(\rho^*(T)\) in
\eqref{eq:strong-indices} gives
\[
\limsup_{x,r\to\infty}
\frac{T(xr)}{x^{p_+}T(r)}
<\infty.
\]
Since \(p_-<\rho_*(T)\), the definition of \(\rho_*(T)\) gives
\[
\liminf_{x,r\to\infty}
\frac{T(xr)}{x^{p_-}T(r)}
>0.
\]
Consequently there are \(c,C>0\) and \(X,r_0\geq1\) such that
\begin{equation}\label{eq:large-multiplier-power-bounds}
 c\,x^{\rho-\epsilon}
 \leq
 \frac{T(xr)}{T(r)}
 \leq
 C\,x^{\rho+\epsilon},
 \qquad
 x\geq X,\quad r\geq r_0.
\end{equation}

We first extend \eqref{eq:large-multiplier-power-bounds} to
\(1\leq x\leq X\). By monotonicity,
\(T(xr)\leq T(Xr)\), and applying the upper bound in
\eqref{eq:large-multiplier-power-bounds} with multiplier \(X\) gives
\[
\frac{T(xr)}{T(r)}
\leq
CX^{\rho+\epsilon}.
\]
Since \(x\geq1\), this is bounded by
\(CX^{\rho+\epsilon}x^{\rho+\epsilon}\).
For the lower bound, monotonicity gives
\(T(xr)/T(r)\geq1\), and
\(1\geq X^{-(\rho-\epsilon)}x^{\rho-\epsilon}\).
After enlarging one constant, there are
\(C_1\geq1\) and \(r_1\geq r_0\) such that
\begin{equation}\label{eq:multiplier-at-least-one}
 C_1^{-1}x^{\rho-\epsilon}
 \leq
 \frac{T(xr)}{T(r)}
 \leq
 C_1x^{\rho+\epsilon},
 \qquad
 x\geq1,\quad r\geq r_1.
\end{equation}

Now let \(0<x<1\), and assume \(r,xr\geq r_1\). Apply
\eqref{eq:multiplier-at-least-one} at the base radius \(xr\) with
multiplier \(1/x\). This gives
\[
C_1^{-1}x^{-(\rho-\epsilon)}
\leq
\frac{T(r)}{T(xr)}
\leq
C_1x^{-(\rho+\epsilon)}.
\]
After inversion,
\begin{equation}\label{eq:multiplier-less-than-one}
 C_1^{-1}x^{\rho+\epsilon}
 \leq
 \frac{T(xr)}{T(r)}
 \leq
 C_1x^{\rho-\epsilon},
 \qquad
 0<x<1,\quad r,xr\geq r_1.
\end{equation}

For \(x\geq1\), the minimum and maximum in
\eqref{eq:power-bounds} are respectively
\(x^{\rho-\epsilon}\) and \(x^{\rho+\epsilon}\);
for \(0<x<1\), they are respectively
\(x^{\rho+\epsilon}\) and \(x^{\rho-\epsilon}\).
Thus \eqref{eq:multiplier-at-least-one} and
\eqref{eq:multiplier-less-than-one} give
\eqref{eq:power-bounds}, with
\(C_\epsilon=C_1\) and \(r_\epsilon=r_1\).
\end{proof}


\section{The case of order less than one}\label{sec:smallorder}

Proposition~\ref{prop:indices} leaves open the possibility that the
common growth index is zero. In fact, the argument in this section
applies in the larger range \(\rho(\fv)<1\) and does not require
P\'olya peaks.

The main point is that Proposition~\ref{prop:initial-basis}, which was
proved for polynomial solution spaces, has an analogue for entire
solution spaces of order less than one. Here the order of a
nonconstant entire function \(h\) is
\[
 \rho(h)
 :=
 \limsup_{r\to\infty}
 \frac{\log^+\log^+M(r,h)}{\log r}.
\]


\subsection{An entire version of the initial basis estimate}

We first pass from polynomial spaces to finite-dimensional spaces of
entire functions of order less than one.

\begin{lem}\label{lem:entire-majorant}
Let \(\yv=(y_0,\ldots,y_{n})\) be a vector of entire functions of
orders less than one, and suppose that
\(y_j^{(i)}(0)=\delta_{ij}\) for \(0\leq i,j\leq n\).
Let \((a_\ell)\) be the zeros of \(W(\yv)\), repeated according to
multiplicity. Then
\begin{equation}\label{eq:entire-majorant}
 |y_j(z)|
 \leq
 \frac{|z|^j}{j!}
 \prod_\ell
 \left(1+\frac{|z|}{|a_\ell|}\right),
 \qquad
 0\leq j\leq n.
\end{equation}
The product in \eqref{eq:entire-majorant} converges because
\(\sum_\ell|a_\ell|^{-1}<\infty\).

If
\[
 H_{\yv}(r)
 :=
 \log\max_{0\leq j\leq n}M(r,y_j),
 \qquad
 N_{\yv}(r)
 :=
 \sum_\ell\log^+\frac{r}{|a_\ell|},
\]
then
\begin{equation}\label{eq:convolution}
 H_{\yv}(r)
 \leq
 n\log^+r
 +
 r\int_0^\infty
 \frac{N_{\yv}(t)}{(r+t)^2}\,\dd t.
\end{equation}
\end{lem}

\begin{proof}

\medskip\noindent
\textbf{\underline{Step 1. Taylor polynomial approximations:}}
Choose \(\sigma\in(0,1)\) larger than the orders of all the
\(y_j\). Then, after increasing a constant if necessary,
\begin{equation}\label{eq:small-order-component-growth}
 \log M(r,y_j)\leq C r^\sigma,
 \qquad
 r\geq1,\quad 0\leq j\leq n.
\end{equation}

Write \(y_j(z)=\sum_{m=0}^\infty c_{j,m}z^m\), and let
\(p_{j,N}(z)=\sum_{m=0}^Nc_{j,m}z^m\), where \(N\geq n\).
Cauchy's estimate for the coefficients on \(|z|=2t\) gives
\[
 M(t,p_{j,N})
 \leq
 M(2t,y_j)\sum_{m=0}^N2^{-m}
 \leq
 2M(2t,y_j).
\]
Applying the preceding estimate with \(2t\) in place of \(t\), and
then applying Cauchy's derivative estimate on \(D_{2t}\), gives
\[
 M(t,p_{j,N}^{(k)})
 \leq
 C_k t^{-k}M(2t,p_{j,N})
 \leq
 2C_k t^{-k}M(4t,y_j),
 \qquad 0\leq k\leq n.
\]
Hence, by \eqref{eq:small-order-component-growth},
\[
 \log^+M(t,p_{j,N}^{(k)})
 \leq C_1t^\sigma,
 \qquad t\geq1,
\]
where \(C_1\) is independent of \(N\).

Put
\(W_N=W(p_{0,N},\ldots,p_{n,N})\).
Expanding this determinant into its \({n+1}!\) terms gives
\begin{equation}\label{eq:WN-growth}
 \log M(t,W_N)\leq C_2t^\sigma,
 \qquad
 t\geq1,
\end{equation}
with \(C_2\) independent of \(N\).

The normalization at the origin is preserved by Taylor truncation:
\(p_{j,N}^{(i)}(0)=\delta_{ij}\) for
\(0\leq i,j\leq n\). Hence \(W_N(0)=1\).


\medskip\noindent
\textbf{\underline{Step 2. Uniform control of the zeros of \(W_N\):}}
Let \(n_{W_N}(t)\) denote the number of zeros of \(W_N\) in \(D_t\),
counted with multiplicity. Jensen's formula applied at radius \(2t\),
together with \(W_N(0)=1\) and \eqref{eq:WN-growth}, gives
\begin{equation}\label{eq:WN-zero-count}
 n_{W_N}(t)
 \leq
 \frac{\log M(2t,W_N)}{\log2}
 \leq
 C_3t^\sigma,
 \qquad
 t\geq1,
\end{equation}
where \(C_3\) is independent of \(N\).

Let \((a_{N,\ell})\) be the zeros of \(W_N\), repeated according to
multiplicity. Integration by parts for the counting function gives
\[
 \sum_{|a_{N,\ell}|>R}\frac1{|a_{N,\ell}|}
 \leq
 \int_R^\infty
 \frac{n_{W_N}(t)}{t^2}\,\dd t.
\]
Using \eqref{eq:WN-zero-count}, we obtain the uniform tail estimate
\begin{equation}\label{eq:reciprocal-tail}
 \sum_{|a_{N,\ell}|>R}\frac1{|a_{N,\ell}|}
 \leq
 C_4R^{\sigma-1},
 \qquad
 R\geq1.
\end{equation}


\medskip\noindent
\textbf{\underline{Step 3. Passage to the entire system:}}
The vectors \(\pv_N=(p_{0,N},\ldots,p_{n,N})\) are linearly
independent. Their first \(n+1\) derivative rows at the origin form the
identity matrix. Proposition~\ref{prop:initial-basis}, applied at
\(a=0\), therefore gives
\begin{equation}\label{eq:taylor-section-majorant}
 |p_{j,N}(z)|
 \leq
 \frac{|z|^j}{j!}
 \prod_\ell
 \left(1+\frac{|z|}{|a_{N,\ell}|}\right).
\end{equation}

As \(N\to\infty\), the functions \(p_{j,N}\) and all their derivatives
converge locally uniformly to \(y_j\). Hence \(W_N\to W(\yv)\) locally uniformly. Since \(W(\yv)(0)=1\), none of the zeros of \(W(\yv)\) is zero, and no zero of \(W_N\) approaches the origin. If \(R>0\) is chosen so that
\(W(\yv)\) has no zero on \(|z|=R\), Rouch\'e's theorem shows that the
zeros of \(W_N\) in \(D_R\) converge, with multiplicities, to the zeros
of \(W(\yv)\) in \(D_R\).

For fixed \(r>0\), the contribution of the zeros of \(W_N\) outside
\(D_R\) satisfies, by \(\log(1+x)\leq x\) and
\eqref{eq:reciprocal-tail},
\begin{equation}\label{eq:product-tail}
 \sum_{|a_{N,\ell}|>R}
 \log\left(1+\frac{r}{|a_{N,\ell}|}\right)
 \leq
 C_4rR^{\sigma-1}.
\end{equation}
First let \(N\to\infty\) in \eqref{eq:taylor-section-majorant} with
\(R\) fixed, using the local convergence of the zeros. Then let
\(R\to\infty\) in \eqref{eq:product-tail}. Since \(\sigma<1\), the
right-hand side of \eqref{eq:product-tail} tends to zero. This proves
\eqref{eq:entire-majorant}.

The same argument, or Fatou's lemma applied to
\eqref{eq:reciprocal-tail}, gives
\begin{equation}\label{eq:entire-reciprocal-sum}
 \sum_\ell\frac1{|a_\ell|}<\infty.
\end{equation}


\medskip\noindent
\textbf{\underline{Step 4. The convolution estimate:}}
Taking the maximum over \(j\) in \eqref{eq:entire-majorant} gives
\[
 H_{\yv}(r)
 \leq
 n\log^+r
 +
 \sum_\ell
 \log\left(1+\frac{r}{|a_\ell|}\right).
\]
For every \(a\ne0\), direct integration gives
\begin{equation}\label{eq:kernel-identity}
 r\int_0^\infty
 \frac{\log^+(t/|a|)}{(r+t)^2}\,\dd t
 =
 \log\left(1+\frac r{|a|}\right).
\end{equation}
By \eqref{eq:entire-reciprocal-sum}, the sum on the right is finite.
Tonelli's theorem \cite[Theorem~2.37(a), p.~67]{Fol99} applied to
\eqref{eq:kernel-identity} therefore gives
\[
 \sum_\ell
 \log\left(1+\frac r{|a_\ell|}\right)
 =
 r\int_0^\infty
 \frac{N_{\yv}(t)}{(r+t)^2}\,\dd t.
\]
Substitution into the estimate for \(H_{\yv}(r)\) proves
\eqref{eq:convolution}.
\end{proof}


\subsection{A sequence of useful radii}

The convolution estimate \eqref{eq:convolution} will be applied at
radii for which \(H(3t)\) can be compared with \(H(r)\). The following
elementary lemma provides such radii.

\begin{lem}\label{lem:envelope}
Let \(H\) be continuous, nonnegative, and nondecreasing, and suppose
that \(H(r)>0\) for all sufficiently large \(r\). If
\(H(r)=o(r^\alpha)\) for some \(0<\alpha<1\), then there are
\(r_\nu\to\infty\) such that
\begin{equation}\label{eq:envelope}
 r_\nu\int_0^\infty
 \frac{H(3t)}{(r_\nu+t)^2}\,\dd t
 \leq
 I_\alpha H(r_\nu),
 \qquad
 I_\alpha
 :=
 \int_0^\infty
 \frac{\max\{1,(3u)^\alpha\}}{(1+u)^2}\,\dd u.
\end{equation}
Moreover, \(I_\alpha<\infty\) for \(0<\alpha<1\), and
\(I_\alpha\to1\) as \(\alpha\downarrow0\).
\end{lem}

\begin{proof}

\medskip\noindent
\textbf{\underline{Step 1. Choice of the radii:}}
Choose any sequence \(R_\nu\to\infty\). Since
\(H(t)/t^\alpha\to0\), the continuous function
\(H(t)/t^\alpha\) attains its maximum on \([R_\nu,\infty)\).
Choose \(r_\nu\geq R_\nu\) at which this maximum is attained. Then
\(r_\nu\to\infty\), and
\begin{equation}\label{eq:envelope-maximizer}
 \frac{H(t)}{t^\alpha}
 \leq
 \frac{H(r_\nu)}{r_\nu^\alpha},
 \qquad
 t\geq r_\nu.
\end{equation}

If \(3u\leq1\), monotonicity gives
\(H(3r_\nu u)\leq H(r_\nu)\). If \(3u\geq1\), apply
\eqref{eq:envelope-maximizer} with \(t=3r_\nu u\). Thus, for every
\(u>0\),
\begin{equation}\label{eq:envelope-pointwise}
 H(3r_\nu u)
 \leq
 H(r_\nu)\max\{1,(3u)^\alpha\}.
\end{equation}


\medskip\noindent
\textbf{\underline{Step 2. Integration:}}
In the left-hand side of \eqref{eq:envelope}, make the change of
variables \(t=r_\nu u\). Equation
\eqref{eq:envelope-pointwise} then gives exactly
\eqref{eq:envelope}.

For large \(u\), the integrand defining \(I_\alpha\) is
\(O(u^{\alpha-2})\), so \(I_\alpha<\infty\) because \(\alpha<1\).
Fix \(0<\alpha_0<1\). For \(0<\alpha\leq\alpha_0\), the integrand is
bounded by
\(\max\{1,(3u)^{\alpha_0}\}/(1+u)^2\), which is integrable. Since the
integrand tends pointwise to \((1+u)^{-2}\) as \(\alpha\downarrow0\),
dominated convergence gives
\(I_\alpha\to\int_0^\infty(1+u)^{-2}\,\dd u=1\).
\end{proof}


\subsection{Classification when the upper order is less than one}

We now apply Lemmas~\ref{lem:entire-majorant} and
\ref{lem:envelope} to the curve. We use the following form of a
rational normal curve:
\begin{equation}\label{eq:rationalnormal}
 [\fv(z)]
 =
 [(1,z,\ldots,z^{n})A],
 \qquad
 A\in\operatorname{GL}_{n+1}(\C).
\end{equation}
Thus \eqref{eq:rationalnormal} is a statement about the projective
curve and not about a particular reduced representation.

\begin{prop}\label{prop:small-order}
Let \(f=[\fv]\) be linearly non-degenerate. If \(\rho(\fv)<1\) and
\(N_1(r,\fv)=o(T(r,\fv))\), then \([\fv]\) has the form
\eqref{eq:rationalnormal}.
\end{prop}
\begin{proof}

\medskip\noindent
\textbf{\underline{Step 1. A representation whose components have order less than one:}}
After a constant change of homogeneous coordinates, assume
\(f_0(0)\ne0\). Fix
\(\beta\) with \(\rho(\fv)<\beta<1\). By the definition of
\(\rho(\fv)\), \(T(r,\fv)=O(r^\beta)\). Jensen's formula and the
standard coordinate estimate give
\[
 \NN_{f_0}(r)
 \leq
 T(r,\fv)+O(1)
 =
 O(r^\beta).
\]
If \(n_{f_0}(r)\) denotes the number of zeros of \(f_0\) in \(D_r\),
then
\(n_{f_0}(r)\log2\leq\NN_{f_0}(2r)\), so
\(n_{f_0}(r)=O(r^\beta)\). Since \(\beta<1\), this implies
\[
 \sum_{f_0(a)=0}\frac1{|a|}<\infty.
\]
Thus the genus-zero canonical product
\(P(z)=\prod_{f_0(a)=0}(1-z/a)\) converges.

For \(|a|\leq r\),
\(\log(1+r/|a|)
 \leq\log2+\log(r/|a|)\), while for \(|a|>r\),
\(\log(1+r/|a|)\leq r/|a|\). Using
\(\NN_{f_0}(r)=O(r^\beta)\) for the first group and
\(n_{f_0}(t)=O(t^\beta)\) for the second gives
\begin{equation}\label{eq:canonical-product-growth}
 \log M(r,P)\leq C_\beta r^\beta,
 \qquad
 r\geq1.
\end{equation}

Since \(P\) and \(f_0\) have the same zeros with the same
multiplicities, \(f_0/P\) is zero-free. Hence
\(f_0=Pe^G\) for an entire function \(G\). Put
\(\gv=e^{-G}\fv\). Then \(g_0=P\), and for every \(j\),
\(g_j=P(f_j/f_0)\) is entire.

By the product inequality for the Nevanlinna characteristic,
\eqref{eq:quotientbound}, and \eqref{eq:canonical-product-growth},
\[
 \TN(r,g_j)
 \leq
 \TN(r,P)
 +
 \TN\left(r,\frac{f_j}{f_0}\right)
 +O(1)
 =
 O(r^\beta).
\]
Applying the Poisson estimate to \(\log|g_j|\) on \(D_{2r}\) gives
\begin{equation}\label{eq:gauge-component-growth}
 \log M(r,g_j)
 \leq
 3\TN(2r,g_j)
 =
 O(r^\beta).
\end{equation}
Since \(\beta>\rho(\fv)\) was arbitrary, every \(g_j\) has order at
most \(\rho(\fv)\), and hence order less than one.

The transformation \(\fv\mapsto e^{-G}\fv\) does not change the
projective curve. Moreover,
\(W(\gv)=e^{-(n+1)G}W(\fv)\), so it does not change the zero divisor of
the Wronskian.


\medskip\noindent
\textbf{\underline{Step 2. Normalization at an ordinary point:}}
Choose \(b\in\C\) such that \(W(\gv)(b)\ne0\), and put \(B=|b|\).
Since the matrix
\((g_j^{(i)}(b))_{0\leq i,j\leq n}\) is invertible, a constant
change of basis in the components of \(\gv(b+z)\) gives an entire
vector
\(\yv=(y_0,\ldots,y_{n})\) satisfying
\[
 y_j^{(i)}(0)=\delta_{ij},
 \qquad
 0\leq i,j\leq n.
\]
Constant changes of basis and translations preserve the property of
having order less than one, so Lemma~\ref{lem:entire-majorant} applies
to \(\yv\).

Let \(H_{\yv}\) and \(N_{\yv}\) be as in
Lemma~\ref{lem:entire-majorant}. Norm equivalence under the fixed
change of basis, together with
\(|z-b|\leq s+B\) for \(|z|=s\), gives
\begin{equation}\label{eq:small-order-T}
 T(s,\fv)
 \leq
 H_{\yv}(s+B)+C,
 \qquad
 s\geq1.
\end{equation}
Notice also that \(H_{\yv}(r)\geq0\), because \(y_0(0)=1\). Since
\(y_1'(0)=1\), Cauchy's estimate gives \(M(r,y_1)\geq r\), and hence
\begin{equation}\label{eq:H-lower-log}
 H_{\yv}(r)\geq\log r,
 \qquad
 r\geq1.
\end{equation}


\medskip\noindent
\textbf{\underline{Step 3. Translated Wronskian zeros:}}
The Wronskian \(W(\gv)\) also has order less than one. Its determinant
expansion involves only derivatives and finite products of the functions
\(g_j\), all of which have order less than one by
\eqref{eq:gauge-component-growth}. Jensen's formula therefore gives
\begin{equation}\label{eq:wronskian-reciprocal-sum}
 \sum_{\substack{W(\gv)(a)=0\\ a\ne0}}
 \frac1{|a|}
 <\infty.
\end{equation}

The zeros of \(W(\yv)\) are the numbers \(a-b\), where \(a\) runs
through the zeros of \(W(\gv)\), with multiplicities preserved.
For \(|a|>2B\),
\[
 \log^+\frac{t}{|a-b|}
 \leq
 \log^+\frac{t+B}{|a|}
 +
 \log^+\frac{|a|}{|a-b|}
 \leq
 \log^+\frac{t+B}{|a|}+\frac{2B}{|a|}.
\]
The last terms are summable by
\eqref{eq:wronskian-reciprocal-sum}. For each of the finitely many nonzero zeros with
\(0<|a|\leq2B\), put \(c_a=|a-b|>0\). The elementary inequality
\[
 \log^+\frac{t}{c_a}
 \leq
 \log^+\frac{t+B}{|a|}
 +\log^+\frac{|a|}{c_a}
\]
shows that its contribution is bounded by the corresponding term in
\(\NN_{W_{\fv}}(t+B)\) plus a constant. A zero at the origin is
handled using the standard \(n(0)\log t\) term in the Nevanlinna
counting function. Since the gauge in Step~1
does not change the Wronskian zero divisor, we obtain
\begin{equation}\label{eq:translated-count}
 N_{\yv}(t)
 \leq
 \NN_{W_{\fv}}(t+B)+C_b,
 \qquad
 t\geq1.
\end{equation}

By \(N_1(r,\fv)=\NN_{W_{\fv}}(r)+O(1)\) and the hypothesis of the
proposition,
\(\NN_{W_{\fv}}(r)=o(T(r,\fv))\). For sufficiently large \(t\),
\(t+2B\leq3t\), and \eqref{eq:small-order-T} gives
\[
 T(t+B,\fv)
 \leq
 H_{\yv}(t+2B)+C
 \leq
 H_{\yv}(3t)+C.
\]
Equations \eqref{eq:translated-count} and \eqref{eq:H-lower-log}
therefore imply
\[
 N_{\yv}(t)
 =
 o(H_{\yv}(3t)).
\]
Consequently, for every \(\varepsilon>0\), there is a constant
\(C_\varepsilon\) such that
\begin{equation}\label{eq:small-order-divisor-majorant}
 N_{\yv}(t)
 \leq
 \varepsilon H_{\yv}(3t)+C_\varepsilon,
 \qquad
 t>0.
\end{equation}


\medskip\noindent
\textbf{\underline{Step 4. Absorption of the Wronskian contribution:}}
Choose \(\alpha\in(0,1)\) larger than the orders of all the \(y_j\).
Then \(H_{\yv}(r)=o(r^\alpha)\), so
Lemma~\ref{lem:envelope} gives a sequence \(r_\nu\to\infty\)
satisfying \eqref{eq:envelope}.

Substitute \eqref{eq:small-order-divisor-majorant} into
\eqref{eq:convolution}. The term containing \(H_{\yv}(3t)\) is
bounded by \(\varepsilon I_\alpha H_{\yv}(r_\nu)\) by
\eqref{eq:envelope}, while
\(r_\nu\int_0^\infty(r_\nu+t)^{-2}\,\dd t=1\).
Thus
\begin{equation}\label{eq:absorption}
 (1-\varepsilon I_\alpha)H_{\yv}(r_\nu)
 \leq
 n\log r_\nu+C_\varepsilon.
\end{equation}

Choose \(\varepsilon>0\) so that
\(0<\varepsilon I_\alpha<1/{n+1}\), and put
\(\gamma=n/(1-\varepsilon I_\alpha)\). Then \(\gamma<{n+1}\), and
\eqref{eq:absorption} gives
\begin{equation}\label{eq:small-order-polynomial-growth}
 M(r_\nu,y_j)
 \leq
 C r_\nu^\gamma,
 \qquad
 0\leq j\leq n.
\end{equation}

For every integer \(k\geq n+1\), Cauchy's estimate and
\eqref{eq:small-order-polynomial-growth} give
\[
 \frac{|y_j^{(k)}(0)|}{k!}
 \leq
 \frac{M(r_\nu,y_j)}{r_\nu^k}
 \leq
 C r_\nu^{\gamma-k}
 \longrightarrow0.
\]
Hence \(y_j^{(k)}(0)=0\) for every \(k\geq n+1\). Thus each \(y_j\) is
a polynomial of degree at most \(n\). The normalization
\(y_j^{(i)}(0)=\delta_{ij}\) now gives
\(y_j(z)=z^j/j!\).

Undoing the translation, the constant changes of basis, and the
zero-free gauge proves \eqref{eq:rationalnormal}.
\end{proof}


\subsection{A ramification lower bound in order zero}

The same argument gives a quantitative conclusion for transcendental
curves of order zero, without assuming small ramification. Recall that
\(N_1(r,\fv)=\NN_{W_{\fv}}(r)+O(1)\). Since
\(T(r,\fv)\to\infty\) for a linearly nondegenerate curve, replacing
\(N_1(r,\fv)\) by \(\NN_{W_{\fv}}(r)\) does not change the following
limsup.

\begin{prop}\label{prop:zero-order-ramification}
Let \(\fv\) be a linearly nondegenerate transcendental curve of order
zero. Then
\begin{equation}\label{eq:zero-order-ramification}
 \limsup_{r\to\infty}
 \frac{N_1(r,\fv)}{T(r,\fv)}
 \geq1.
\end{equation}
\end{prop}

\begin{proof}

\medskip\noindent
\textbf{\underline{Step 1. Reduction to the normalized entire system:}}
Suppose, to the contrary, that the left-hand side of
\eqref{eq:zero-order-ramification} is less than one. Choose \(c\) such
that
\[
 \limsup_{r\to\infty}
 \frac{N_1(r,\fv)}{T(r,\fv)}
 <c<1.
\]
Since \(N_1(r,\fv)=\NN_{W_{\fv}}(r)+O(1)\) and
\(T(r,\fv)\to\infty\), there is \(r_0\) such that
\begin{equation}\label{eq:zero-order-counting-assumption}
 \NN_{W_{\fv}}(r)
 \leq
 c\,T(r,\fv),
 \qquad
 r\geq r_0.
\end{equation}

Apply the constructions in Steps~1 and~2 of the proof of
Proposition~\ref{prop:small-order}. These steps use only
\(\rho(\fv)<1\), not the small-ramification hypothesis. Since now
\(\rho(\fv)=0\), the resulting normalized functions
\(y_0,\ldots,y_{n}\) all have order zero. In particular,
\(H_{\yv}(r)=o(r^\alpha)\) for every \(\alpha>0\).

The estimates \eqref{eq:small-order-T} and
\eqref{eq:translated-count} remain valid. Applying
\eqref{eq:zero-order-counting-assumption} at \(t+B\), and then using
\eqref{eq:small-order-T}, gives, after increasing a constant to cover
bounded \(t\),
\begin{equation}\label{eq:zero-order-divisor-majorant}
 N_{\yv}(t)
 \leq
 c\,H_{\yv}(3t)+C,
 \qquad
 t>0.
\end{equation}


\medskip\noindent
\textbf{\underline{Step 2. The convolution estimate forces polynomial growth:}}
By Lemma~\ref{lem:envelope},
\(I_\alpha\to1\) as \(\alpha\downarrow0\). Since \(c<1\), choose
\(\alpha\in(0,1)\) so small that
\(cI_\alpha<1\). Because every \(y_j\) has order zero,
\(H_{\yv}(r)=o(r^\alpha)\), and Lemma~\ref{lem:envelope} supplies
\(r_\nu\to\infty\).

Substituting \eqref{eq:zero-order-divisor-majorant} into
\eqref{eq:convolution}, using \eqref{eq:envelope}, and using
\(r_\nu\int_0^\infty(r_\nu+t)^{-2}\,\dd t=1\), gives
\begin{equation}\label{eq:zero-order-absorption}
 (1-cI_\alpha)H_{\yv}(r_\nu)
 \leq
 n\log r_\nu+C.
\end{equation}
By \eqref{eq:zero-order-absorption}, with
\(\Gamma=n/(1-cI_\alpha)<\infty\),
\begin{equation}\label{eq:zero-order-polynomial-growth}
 M(r_\nu,y_j)
 \leq
 C_1r_\nu^\Gamma,
 \qquad
 0\leq j\leq n.
\end{equation}

Choose an integer \(k>\Gamma\). Cauchy's estimate and
\eqref{eq:zero-order-polynomial-growth} give
\[
 \frac{|y_j^{(k)}(0)|}{k!}
 \leq
 C_1r_\nu^{\Gamma-k}
 \longrightarrow0.
\]
The same argument applies to every integer larger than \(\Gamma\), so
each \(y_j\) is a polynomial. Therefore the projective curve
\([\fv]\) is rational, contradicting the assumption that it is
transcendental.

This proves \eqref{eq:zero-order-ramification}.
\end{proof}

\section{Regular variation}\label{sec:regular}

Throughout this section we assume that
\(\NN_{W_{\fv}}(r)=o(T(r,\fv))\) and
\[
\rho_*(T)=\rho^*(T)=\rho\in\Gn,
\qquad
T(r):=T(r,\fv).
\]
Since
\(\rho_*(T)\leq\lambda(\fv)\leq\rho(\fv)\leq\rho^*(T)\),
both the lower and upper orders of \(\fv\) are equal to \(\rho>0\).
In particular, \(\log r=o(T(r))\).

The proof has three parts. We begin with an arbitrary sequence of
radii and construct limiting differential coefficients and a limiting
norm profile. We then show that this profile is homogeneous of degree
\(\rho\). Finally, the homogeneity determines every subsequential limit
of \(T(cr)/T(r)\), which yields regular variation.


\subsection{Limits along arbitrary radii}

Fix an arbitrary sequence \(r_\nu\to\infty\), and put
\(s_\nu=T(r_\nu)\).

\medskip\noindent
\textbf{\underline{Step 1. Limits of the differential coefficients:}}
Fix \(0<\epsilon<\rho\), and write
\[
M_\epsilon(R)
:=
\max\{R^{\rho-\epsilon},R^{\rho+\epsilon}\},
\qquad R>0.
\]
Lemma~\ref{lem:power-bounds} gives constants
\(C_\epsilon\geq1\) and \(r_\epsilon>0\) such that
\begin{equation}\label{eq:regular-power-upper}
 \frac{T(Rr)}{T(r)}
 \leq
 C_\epsilon M_\epsilon(R)
 \qquad
 (r,Rr\geq r_\epsilon).
\end{equation}

Fix \(R>0\). We apply Proposition~\ref{prop:representation} with
\(t_\nu=Rr_\nu\) and
\(s_{\nu,R}=M_\epsilon(R)s_\nu\).
The logarithmic condition follows from
\(\log r=o(T(r))\). Moreover, applying
\eqref{eq:regular-power-upper} with multiplier \(256R\) gives
\[
\frac{T(256Rr_\nu)}{M_\epsilon(R)s_\nu}
\leq
C_\epsilon
\frac{M_\epsilon(256R)}{M_\epsilon(R)}
\leq
C_\epsilon256^{\rho+\epsilon}.
\]
Thus the characteristic bound required in
Proposition~\ref{prop:representation} is independent of \(R\).
Furthermore,
\[
\frac{\NN_{W_{\fv}}(256Rr_\nu)}
     {M_\epsilon(R)s_\nu}
=
\frac{\NN_{W_{\fv}}(256Rr_\nu)}
     {T(256Rr_\nu)}
\frac{T(256Rr_\nu)}
     {M_\epsilon(R)s_\nu},
\]
and the first factor tends to zero while the second is bounded by the
estimate just obtained. Hence the ramification hypothesis of
Proposition~\ref{prop:representation} also holds.

For \(2\leq q\leq n+1\), set
\begin{equation}\label{eq:arbitrary-scaled-coefficients}
 A_{q,\nu}(z)
 :=
 \left(\frac{r_\nu}{s_\nu}\right)^q
 Q_q(r_\nu z).
\end{equation}
At the scale \(Rr_\nu\), the coefficient appearing in
Proposition~\ref{prop:representation} is related to
\eqref{eq:arbitrary-scaled-coefficients} by
\begin{equation}\label{eq:arbitrary-scaling-relation}
 A_{q,\nu}(Rz)
 =
 \left(\frac{M_\epsilon(R)}{R}\right)^q
 \left(
 \frac{Rr_\nu}{M_\epsilon(R)s_\nu}
 \right)^q
 Q_q(Rr_\nu z).
\end{equation}

Apply Proposition~\ref{prop:representation} successively with
\(R=1,2,3,\ldots\), and take a diagonal subsequence. Equation
\eqref{eq:arbitrary-scaling-relation} gives entire functions
\(a_q\), \(2\leq q\leq n+1\), such that
\begin{equation}\label{eq:arbitrary-coefficients}
 A_{q,\nu}\longrightarrow a_q
 \qquad\text{locally in measure on }\C.
\end{equation}
The uniqueness of limits in measure shows that the limits obtained on
overlapping disks agree.

The uniform bound in
Proposition~\ref{prop:representation}, applied in
\eqref{eq:arbitrary-scaling-relation}, gives
\begin{equation}\label{eq:arbitrary-coefficient-growth}
 \sup_{|z|\leq R}|a_q(z)|
 \leq
 K_{{n+1},\rho,\epsilon}
 \max\{
 R^{q(\rho-1-\epsilon)},
 R^{q(\rho-1+\epsilon)}
 \},
 \qquad R>0.
\end{equation}


\medskip\noindent
\textbf{\underline{Step 2. The limiting coefficients are monomials:}}
Although \(\epsilon\) was fixed above to obtain compactness, the limit
\(a_q\) is independent of that choice. Indeed, once
\eqref{eq:arbitrary-coefficients} holds, for any
\(0<\epsilon'<\rho\) the same application of
Proposition~\ref{prop:representation} at a fixed scale \(R\), together with
uniqueness of limits in measure, yields
\eqref{eq:arbitrary-coefficient-growth} with \(\epsilon'\) in place of
\(\epsilon\). Thus the latter estimate may be used with arbitrarily small
positive \(\epsilon\).

Write
\(a_q(z)=\sum_{m=0}^{\infty}c_{q,m}z^m\).
For \(R\geq1\), Cauchy's estimate and
\eqref{eq:arbitrary-coefficient-growth} give
\[
|c_{q,m}|
\leq
K_{{n+1},\rho,\epsilon}
R^{q(\rho-1+\epsilon)-m}.
\]
If \(m>q(\rho-1)\), choose \(\epsilon>0\) so small that
\(m>q(\rho-1+\epsilon)\), and let \(R\to\infty\). Then
\(c_{q,m}=0\).

For \(0<R\leq1\), equation
\eqref{eq:arbitrary-coefficient-growth} instead gives
\[
|c_{q,m}|
\leq
K_{{n+1},\rho,\epsilon}
R^{q(\rho-1-\epsilon)-m}.
\]
If \(m<q(\rho-1)\), choose \(\epsilon>0\) so small that
\(m<q(\rho-1-\epsilon)\), and let \(R\downarrow0\). Again
\(c_{q,m}=0\).

Thus a Taylor coefficient can be nonzero only when
\(m=q(\rho-1)\). Consequently
\begin{equation}\label{eq:monomial-coefficients}
 a_q(z)
 =
 c_qz^{q(\rho-1)},
 \qquad
 c_q=0
 \quad\text{if }q(\rho-1)\notin\Z_{\geq0}.
\end{equation}
The constants \(c_q\) may depend on the chosen subsequence.


\medskip\noindent
\textbf{\underline{Step 3. Construction of the limiting norm profile:}}
Apply Proposition~\ref{prop:representation} with
\(t_\nu=r_\nu\) and \(s_\nu=T(r_\nu)\), and let
\(\gv_\nu=(g_{\nu,0},\ldots,g_{\nu,n})\) and \(P_\nu\) be the
resulting local data.

The bounds in \eqref{eq:representationbounds} imply that the
subharmonic functions
\(s_\nu^{-1}\log|g_{\nu,j}|\) are locally uniformly bounded above.
Their values at the origin remain bounded from below by
\eqref{eq:representationbounds}. After passing to a subsequence,
Lemma~\ref{lem:subharmonic-compactness} therefore gives
\[
s_\nu^{-1}\log|g_{\nu,j}|
\longrightarrow u_j
\qquad\text{in }L^1_{\loc}(D_4),
\]
where no \(u_j\) is identically \(-\infty\).

Since
\[
0\leq
\log\norm{\gv_\nu}
-
\max_j\log|g_{\nu,j}|
\leq
\frac12\log {n+1},
\]
we also obtain
\begin{equation}\label{eq:norm-limit}
 U_\nu
 :=
 s_\nu^{-1}\log\norm{\gv_\nu}
 \longrightarrow
 U:=\max_j u_j
 \qquad\text{in }L^1_{\loc}(D_4).
\end{equation}

All roots of \(P_\nu\) lie in \(D_{64}\), and
\(\deg P_\nu=o(s_\nu)\). Hence
\eqref{eq:small-polynomial-log} gives
\(s_\nu^{-1}\log|P_\nu|\to0\) in \(L^1_{\loc}\).
Lemma~\ref{lem:sum}, applied to the limits \(u_j\), yields
\(\sum_j u_j\geq0\) almost everywhere. Since
\(U=\max_j u_j\), it follows that
\begin{equation}\label{eq:U-nonnegative}
 U\geq0
 \qquad\text{on }D_4.
\end{equation}

Dividing \eqref{eq:meanidentity} by \(s_\nu=T(r_\nu)\) gives
\begin{equation}\label{eq:radial-mean}
 \mean t{U_\nu}
 =
 \frac{T(tr_\nu)}{T(r_\nu)}
 +
 o(1),
 \qquad
 0<t<4,
\end{equation}
where the error term is independent of \(t\).


\medskip\noindent
\textbf{\underline{Step 4. Behavior of \(U\) at the origin:}}
Fix \(0<\epsilon<\rho\). We first estimate the area integral of \(U\)
over small disks.

Let \(0<R<1/2\). By \eqref{eq:radial-mean},
\[
\int_{D_R}U_\nu\,\dA
=
2\pi
\int_0^R
\frac{T(tr_\nu)}{T(r_\nu)}t\,\dd t
+
o(1).
\]
For \(tr_\nu\geq r_\epsilon\),
Lemma~\ref{lem:power-bounds} gives
\(T(tr_\nu)/T(r_\nu)\leq
C_\epsilon t^{\rho-\epsilon}\).
For \(tr_\nu<r_\epsilon\), monotonicity gives
\(T(tr_\nu)\leq T(r_\epsilon)\), and
\(T(r_\epsilon)/T(r_\nu)\to0\).
Passing to the limit in \eqref{eq:norm-limit} therefore gives
\begin{equation}\label{eq:origin-area-bound}
 \int_{D_R}U\,\dA
 \leq
 C_\epsilon R^{2+\rho-\epsilon},
 \qquad
 0<R<\frac12.
\end{equation}

Let \(0<|z|<1/4\). The submean inequality, the nonnegativity
\eqref{eq:U-nonnegative}, and
\(D(z,|z|)\subset D_{2|z|}\) give
\[
U(z)
\leq
\frac1{\pi|z|^2}
\int_{D_{2|z|}}U\,\dA.
\]
Using \eqref{eq:origin-area-bound} with \(R=2|z|\) yields
\[
U(z)\leq C_\epsilon|z|^{\rho-\epsilon}.
\]
The same estimate at \(z=0\), followed by \(R\downarrow0\), gives
\(U(0)=0\). Thus
\begin{equation}\label{eq:origin-bound}
 0\leq
 U(z)
 \leq
 C_\epsilon|z|^{\rho-\epsilon},
 \qquad
 |z|<\frac14.
\end{equation}


\medskip\noindent
\textbf{\underline{Step 5. Polynomial approximation of the norm profile:}}
Let \(\pv_\nu\) be the polynomial approximants from
Lemma~\ref{lem:replacement}, and put
\(R_\nu=W(\pv_\nu)\).
We claim that
\begin{equation}\label{eq:polynomial-norm}
 s_\nu^{-1}\log\norm{\pv_\nu}
 \longrightarrow U
 \qquad\text{in }L^1_{\loc}(D_4).
\end{equation}

Fix a compact set \(K\Subset D_4\) and \(\eta>0\). By
\eqref{eq:norm-limit} and \eqref{eq:U-nonnegative},
\[
\meas\{
z\in K:\norm{\gv_\nu(z)}<e^{-\eta s_\nu}
\}
\longrightarrow0.
\]
Choose the constant \(A\) in \eqref{eq:taylorerror} with \(A>\eta\).
Outside this exceptional set,
\[
\frac{\norm{\pv_\nu-\gv_\nu}}
     {\norm{\gv_\nu}}
\leq
C e^{-(A-\eta)s_\nu},
\]
and hence
\[
s_\nu^{-1}
 \bigl|\log\norm{\pv_\nu}-\log\norm{\gv_\nu}\bigr|\longrightarrow0
\]
uniformly there. It follows from \eqref{eq:norm-limit} that
\[
s_\nu^{-1}\log\norm{\pv_\nu}\longrightarrow U
\]
in measure on \(K\).

The functions \(s_\nu^{-1}\log\norm{\pv_\nu}\) are subharmonic and
locally uniformly bounded above by \eqref{eq:taylorerror} and
\eqref{eq:representationbounds}. Lemma~\ref{lem:subharmonic-compactness},
together with uniqueness of the limit in measure, proves
\eqref{eq:polynomial-norm}.


\medskip\noindent
\textbf{\underline{Step 6. Choice of good centers:}}
For the Wronskians \(R_\nu\), Lemma~\ref{lem:replacement} gives
\begin{equation}\label{eq:R-log-good}
 \frac{\log|R_\nu|}{s_\nu}
 \convmeas0
 \qquad\text{locally on }D_4.
\end{equation}
After passing to a subsequence, the convergence in
\eqref{eq:R-log-good} holds almost everywhere.

Let
\[
B_\nu(a)
:=
\sum_{R_\nu(\zeta)=0}\frac1{|a-\zeta|},
\]
with zeros repeated according to multiplicity.
Choose the separating radius from Lemma~\ref{lem:replacement} and split
\(B_\nu\) into the contributions of the interior and exterior
roots. By \eqref{eq:rootbounds}, the
interior part divided by \(s_\nu\) tends to zero in
\(L^{3/2}(D_4)\), while the exterior part divided by \(s_\nu\) is
uniformly bounded on \(D_4\). After one further extraction, the
interior part tends to zero almost everywhere.

We may therefore choose a set
\(E\subset D_2\setminus\{0\}\) of full measure such that
\begin{equation}\label{eq:good-centers}
 \frac{\log|R_\nu(a)|}{s_\nu}\longrightarrow0,
 \qquad
 \limsup_{\nu\to\infty}
 \frac{B_\nu(a)}{s_\nu}<\infty,
 \qquad
 a\in E.
\end{equation}
We also remove from \(E\) the countable union of the zero sets of the
polynomials \(R_\nu\), so that \(R_\nu(a)\ne0\) for every
\(a\in E\) and every \(\nu\).


\subsection{Choice of basis at a point}

At a fixed good center we now choose a basis whose limiting component
logarithms are balanced at that point. The unitary nature of the basis
change is important because it leaves the Euclidean norm unchanged.

\begin{lem}\label{lem:basis-at-point}
For every \(a\in E\), there are a further subsequence and constant
unitary matrices \(V_\nu\) such that, with
\(\hv_\nu=\pv_\nu V_\nu\),
\[
s_\nu^{-1}\log|h_{\nu,j}|
\longrightarrow v_j
\qquad\text{in }L^1_{\loc}(D_4),
\qquad
0\leq j\leq n,
\]
where the \(v_j\) are subharmonic and not identically \(-\infty\), and
\begin{equation}\label{eq:point-balance}
 U=\max_jv_j,
 \qquad
 \sum_{j=0}^{n}v_j\geq0,
 \qquad
 \sum_{j=0}^{n}v_j(a)=0.
\end{equation}
The first two relations in \eqref{eq:point-balance} hold on \(D_4\),
with the sum inequality understood almost everywhere.
\end{lem}

\begin{proof}

\medskip\noindent
\textbf{\underline{Step 1. Orthogonalize the normalized initial derivatives:}}
Define the normalized jet matrix
\[
J_\nu(a)
:=
\bigl(
s_\nu^{-i}p_{\nu,j}^{(i)}(a)
\bigr)_{0\leq i,j\leq n}.
\]
Choose a unitary matrix \(V_\nu\) by singular value decomposition so
that the columns of \(J_\nu(a)V_\nu\) are orthogonal. Denote their
lengths by
\(\sigma_{\nu,0},\ldots,\sigma_{\nu,n}>0\).

The local upper bound for \(\pv_\nu\), obtained from
\eqref{eq:taylorerror} and \eqref{eq:representationbounds}, and
Cauchy's derivative estimate give
\(\sigma_{\nu,j}\leq e^{Cs_\nu}\) for a constant \(C\) independent of
\(\nu\) and \(j\).

Since \(V_\nu\) is unitary,
\[
\prod_{j=0}^{n}\sigma_{\nu,j}
=
|\det J_\nu(a)|
=
s_\nu^{-\kappa}|R_\nu(a)|,
\qquad
\kappa=\frac{n(n+1)}2.
\]
By \eqref{eq:good-centers},
\begin{equation}\label{eq:singular-product}
 \frac1{s_\nu}
 \log
 \prod_{j=0}^{n}\sigma_{\nu,j}
 \longrightarrow0.
\end{equation}
The upper bounds
\(\sigma_{\nu,j}\leq e^{Cs_\nu}\), together with
\eqref{eq:singular-product}, give a finite exponential lower bound for
each \(\sigma_{\nu,j}\). After passing to a further subsequence, we may
therefore assume
\begin{equation}\label{eq:singular-exponents}
 \frac{\log\sigma_{\nu,j}}{s_\nu}
 \longrightarrow
 \lambda_j\in\R,
 \qquad
 \sum_{j=0}^{n}\lambda_j=0.
\end{equation}


\medskip\noindent
\textbf{\underline{Step 2. Limits of the new component logarithms:}}
Put \(\hv_\nu=\pv_\nu V_\nu\).
Since \(V_\nu\) is unitary,
\(\norm{\hv_\nu}=\norm{\pv_\nu}\). Hence
\eqref{eq:polynomial-norm} gives
\[
s_\nu^{-1}\log\norm{\hv_\nu}
\longrightarrow U
\qquad\text{in }L^1_{\loc}(D_4).
\]

For each \(j\), the functions
\(s_\nu^{-1}\log|h_{\nu,j}|\) are locally uniformly bounded above.
They cannot converge locally uniformly to \(-\infty\): if they did,
Cauchy's derivative estimates at \(a\) would imply that every entry in
the \(j\)-th column of \(J_\nu(a)V_\nu\) has logarithmic size divided
by \(s_\nu\) tending to \(-\infty\), contradicting
\eqref{eq:singular-exponents}.

Lemma~\ref{lem:subharmonic-compactness} therefore gives, after another
subsequence,
\[
s_\nu^{-1}\log|h_{\nu,j}|
\longrightarrow v_j
\qquad\text{in }L^1_{\loc}(D_4).
\]
Since
\[
0\leq
\log\norm{\hv_\nu}
-
\max_j\log|h_{\nu,j}|
\leq
\frac12\log {n+1},
\]
equation \eqref{eq:polynomial-norm} gives
\(U=\max_jv_j\).

Moreover,
\(W(\hv_\nu)=(\det V_\nu)R_\nu\), so
\(|W(\hv_\nu)|=|R_\nu|\). Equation
\eqref{eq:R-log-good}, together with Lemma~\ref{lem:sum}, gives
\(\sum_jv_j\geq0\) almost everywhere.


\medskip\noindent
\textbf{\underline{Step 3. Upper bound for the values \(v_j(a)\):}}
By construction of the singular values,
\[
|h_{\nu,j}^{(i)}(a)|
\leq
s_\nu^i\sigma_{\nu,j},
\qquad
0\leq i\leq n.
\]
Apply \eqref{eq:initial-value-bound} to \(h_{\nu,j}\) at the point
\(a\). Since \(B_\nu(a)\) is defined in
\eqref{eq:good-centers}, we obtain
\begin{equation}\label{eq:point-initial-upper}
 |h_{\nu,j}(a+z)|
 \leq
 \sigma_{\nu,j}
 \left(
 \sum_{i=0}^{n}
 \frac{(s_\nu|z|)^i}{i!}
 \right)
 e^{|z|B_\nu(a)}.
\end{equation}

Let
\(C_a>\limsup_\nu B_\nu(a)/s_\nu\).
Taking logarithms in \eqref{eq:point-initial-upper}, dividing by
\(s_\nu\), and using \eqref{eq:singular-exponents}, we obtain in the
limit
\[
v_j(a+z)\leq\lambda_j+C_a|z|
\]
for almost every sufficiently small \(z\). The subharmonic
representative satisfies the same inequality by the submean property.
At \(z=0\), this gives
\begin{equation}\label{eq:point-value-upper}
 v_j(a)\leq\lambda_j.
\end{equation}


\medskip\noindent
\textbf{\underline{Step 4. Reverse inequality at \(a\):}}
For fixed \(r>0\) with
\(\overline{D(a,2r)}\Subset D_4\), Cauchy's estimate gives
\[
|h_{\nu,j}^{(i)}(a)|
\leq
\frac{i!}{r^i}
M(r,h_{\nu,j}(a+\cdot)).
\]
Using the definition of \(\sigma_{\nu,j}\), we obtain
\begin{equation}\label{eq:singular-value-cauchy}
 \sigma_{\nu,j}
 \leq
 M(r,h_{\nu,j}(a+\cdot))
 \left(
 \sum_{i=0}^{n}
 \frac{(i!)^2}{(s_\nu r)^{2i}}
 \right)^{1/2}.
\end{equation}
For fixed \(r\), the logarithm of the last factor in
\eqref{eq:singular-value-cauchy}, divided by \(s_\nu\), tends to zero.
The upper-bound part of
Lemma~\ref{lem:subharmonic-compactness} gives
\[
\limsup_{\nu\to\infty}
\sup_{|z-a|\leq r}
s_\nu^{-1}\log|h_{\nu,j}(z)|
\leq
\sup_{|z-a|<2r}v_j(z).
\]
Using this inequality in \eqref{eq:singular-value-cauchy}, together
with \eqref{eq:singular-exponents}, gives
\[
\lambda_j
\leq
\sup_{|z-a|<2r}v_j(z).
\]
Letting \(r\downarrow0\) and using upper semicontinuity of \(v_j\)
gives \(\lambda_j\leq v_j(a)\). Together with
\eqref{eq:point-value-upper}, this proves
\begin{equation}\label{eq:point-values}
 v_j(a)=\lambda_j,
 \qquad
 0\leq j\leq n.
\end{equation}
Finally, summing \eqref{eq:point-values} and using
\eqref{eq:singular-exponents} gives
\(\sum_jv_j(a)=0\), completing the proof of
\eqref{eq:point-balance}.
\end{proof}

The matrices \(V_\nu\) and the extracted subsequence may depend on the
chosen point \(a\). The limit \(U\), however, is fixed by
\eqref{eq:norm-limit} and is independent of this choice.


\subsection{Homogeneity of the norm limit}

We use the following consequence of
Bergkvist--Rullg{\aa}rd \cite[Corollary~1]{BR02}.

\begin{lem}[\cite{BR02}]\label{lem:finite-gradient-convex}
Let \(G\subset\C\) be convex and let \(A\subset\C\) be finite. If
\(v\) is subharmonic on \(G\) and the measurable representative of
\(2\partial v\) belongs to \(A\) almost everywhere, then \(v\) is
convex.
\end{lem}

\begin{prop}\label{prop:homogeneity}
The limit \(U\) in \eqref{eq:norm-limit} satisfies
\begin{equation}\label{eq:homogeneity}
 U(tz)=t^\rho U(z)
 \qquad
 (t>0,\ z,tz\in D_2).
\end{equation}
\end{prop}

\begin{proof}

\medskip\noindent
\textbf{\underline{Step 1. The limiting first-order equation:}}
Let \(a\in E\), and let \(v\) be any one of the component limits
obtained from Lemma~\ref{lem:basis-at-point}. The corresponding
component \(h_\nu\) of \(\hv_\nu\) satisfies the polynomial
fundamental equation
\[
h_\nu^{(n+1)}
+
\sum_{q=1}^{n+1}
b_{q,\nu}^p h_\nu^{({n+1}-q)}
=
0.
\]
After division by \(s_\nu^{n+1}h_\nu\), this becomes
\begin{equation}\label{eq:normalized-polynomial-equation}
 \frac{h_\nu^{(n+1)}}{s_\nu^{n+1}h_\nu}
 +
 \sum_{q=1}^{n+1}
 \frac{b_{q,\nu}^p}{s_\nu^q}
 \frac{h_\nu^{({n+1}-q)}}{s_\nu^{{n+1}-q}h_\nu}
 =
 0.
\end{equation}

For \(q=1\), \eqref{eq:polynomial-coefficients} gives
\(b_{1,\nu}^p/s_\nu\convmeas0\). For \(q\geq2\),
\eqref{eq:polynomial-coefficients},
\eqref{eq:arbitrary-coefficients}, and
\eqref{eq:monomial-coefficients} give
\[
\frac{b_{q,\nu}^p}{s_\nu^q}
\convmeas
c_qz^{q(\rho-1)}.
\]
Lemma~\ref{lem:logderivlimit} gives, for
\(0\leq k\leq n+1\),
\[
\frac{h_\nu^{(k)}}{s_\nu^kh_\nu}
\convmeas
(2\partial v)^k.
\]
Passing to the limit in
\eqref{eq:normalized-polynomial-equation} yields
\begin{equation}\label{eq:gradient-equation}
 (2\partial v)^{n+1}
 +
 \sum_{q=2}^{n+1}
 c_qz^{q(\rho-1)}
 (2\partial v)^{{n+1}-q}
 =
 0
 \qquad\text{a.e.}
\end{equation}


\medskip\noindent
\textbf{\underline{Step 2. Continuity of \(U\) away from the origin:}}
Fix a compact set
\(K\Subset D_4\setminus\{0\}\).
The coefficients of the monic polynomial in
\eqref{eq:gradient-equation} are uniformly bounded on \(K\).
The standard root bound for a monic polynomial therefore gives a
constant \(C_K\) such that
\[
|2\partial v(z)|\leq C_K
\qquad\text{for a.e. }z\in K.
\]
Since \(v\) is real-valued, its full weak gradient is locally bounded.
Thus \(v\in W^{1,\infty}_{\loc}(D_4\setminus\{0\})\), and its
Sobolev representative is locally Lipschitz there.

Choose one point \(a_0\in E\) and the corresponding functions
\(v_0,\ldots,v_{n}\) from
Lemma~\ref{lem:basis-at-point}. By the conclusion just obtained, all
the \(v_j\) are locally Lipschitz on
\(D_4\setminus\{0\}\). Equation \eqref{eq:point-balance} gives
\(U=\max_jv_j\) almost everywhere. The function \(U\) is subharmonic,
and the right-hand side is continuous. Since the two representatives agree
almost everywhere, their upper-semicontinuous representatives agree
everywhere. Hence
\[
U=\max_jv_j
\qquad\text{on }D_4\setminus\{0\},
\]
and \(U\) is continuous there. Equation \eqref{eq:origin-bound} gives
\(U(z)\to0=U(0)\) as \(z\to0\). Thus \(U\) is continuous on \(D_4\).


\medskip\noindent
\textbf{\underline{Step 3. Straightening the radial coefficients:}}
Fix \(a\in E\), and choose the functions
\(v_0,\ldots,v_{n}\) supplied by
Lemma~\ref{lem:basis-at-point} at this point.

Choose a sufficiently narrow sector \(S\subset D_2\) with vertex at
the origin and containing \(a\), such that a branch of
\[
w=\frac{z^\rho}{\rho}
\]
is single-valued and injective on \(S\), and such that its image
\(\Omega\) is convex. Define
\(V_j(w)=v_j(z(w))\).

Since the change of variables is conformal, each \(V_j\) is
subharmonic. The weak chain rule gives
\(2\partial_zv_j=z^{\rho-1}2\partial_wV_j\).
Substitution into \eqref{eq:gradient-equation} gives
\begin{equation}\label{eq:constant-gradient-equation}
 (2\partial_wV_j)^{n+1}
 +
 \sum_{q=2}^{n+1}
 c_q(2\partial_wV_j)^{{n+1}-q}
 =
 0
 \qquad\text{a.e.\ on }\Omega.
\end{equation}
Indeed, every term contains the common factor
\(z^{{n+1}(\rho-1)}\). Since the calculation is performed on the fixed
sector \(S\), where the chosen branch of \(z^\rho\) is single-valued,
this common factor may be cancelled there. No further integrality
assertion about the individual exponents is needed.

Let \(A\) be the finite set of roots of
\[
\xi^{n+1}+\sum_{q=2}^{n+1}c_q\xi^{{n+1}-q}.
\]
Equation \eqref{eq:constant-gradient-equation} says that
\(2\partial_wV_j\in A\) almost everywhere. Therefore
Lemma~\ref{lem:finite-gradient-convex} shows that every \(V_j\) is
convex on \(\Omega\).


\medskip\noindent
\textbf{\underline{Step 4. Equality on radial segments:}}
From \(U=\max_jv_j\) and
\(\sum_jv_j\geq0\), we have, wherever the latter inequality holds,
\[
v_j
\geq
-\sum_{k\ne j}v_k
\geq
-nU.
\]
By continuity away from the origin, this extends throughout the
sector. Hence
\begin{equation}\label{eq:component-sandwich}
 -nU(z)
 \leq
 v_j(z)
 \leq
 U(z).
\end{equation}
Equation \eqref{eq:origin-bound} and
\eqref{eq:component-sandwich} show that
\(v_j(z)\to0\) as \(z\to0\) inside \(S\). Thus \(V_j\) extends
continuously to the vertex \(w=0\), with \(V_j(0)=0\).

Put \(w_0=a^\rho/\rho\). By convexity,
\[
V_j(tw_0)
\leq
tV_j(w_0),
\qquad
0<t<1.
\]
Equation \eqref{eq:point-balance} gives
\(\sum_jV_j(w_0)=0\). The inequality
\(\sum_jv_j\geq0\), together with continuity, also gives
\(\sum_jV_j(tw_0)\geq0\). Therefore
\[
0
\leq
\sum_jV_j(tw_0)
\leq
t\sum_jV_j(w_0)
=
0.
\]
Thus equality holds in every convexity inequality:
\(V_j(tw_0)=tV_j(w_0)\). Taking the maximum over \(j\) and using
\(U=\max_jv_j\) gives
\[
U(t^{1/\rho}a)=tU(a),
\qquad
0<t<1.
\]
Writing \(s=t^{1/\rho}\), we obtain
\begin{equation}\label{eq:homogeneity-on-E}
 U(sa)=s^\rho U(a),
 \qquad
 a\in E,\quad 0<s<1.
\end{equation}

The set \(E\) has full measure and is therefore dense in \(D_2\).
Continuity of \(U\) extends
\eqref{eq:homogeneity-on-E} to every \(a\in D_2\).
For \(s>1\) with \(a,sa\in D_2\), apply the already proved relation
with multiplier \(1/s\) at the point \(sa\). This proves
\eqref{eq:homogeneity}.
\end{proof}


\subsection{Convergence of the characteristic ratios}

We now show that the homogeneous norm profile determines the
characteristic ratio.

\begin{prop}\label{prop:regular-variation}
Under the assumptions of this section,
\begin{equation}\label{eq:ratio-limit}
 \frac{T(cr)}{T(r)}
 \longrightarrow
 c^\rho
 \qquad(r\to\infty),
\end{equation}
locally uniformly for \(c\in(0,\infty)\).
\end{prop}

\begin{proof}

\medskip\noindent
\textbf{\underline{Step 1. Annular averages of the characteristic ratios:}}
Let \(r_\nu\to\infty\) be arbitrary, and perform the extractions made
above. Put
\(F_\nu(t)=T(tr_\nu)/T(r_\nu)\).

Fix \(0<a<b<2\). Integrating \eqref{eq:radial-mean} with weight
\(t\,\dd t\), and using that the error in
\eqref{eq:radial-mean} is independent of \(t\), gives
\[
\int_a^bF_\nu(t)t\,\dd t
=
\frac1{2\pi}
\int_{\{a<|z|<b\}}U_\nu(z)\,\dA
+
o(1).
\]
By \eqref{eq:norm-limit}, the area integral converges to the
corresponding integral of \(U\). By
\eqref{eq:homogeneity}, if
\(K:=\mean1U\), then \(\mean tU=Kt^\rho\). Hence
\begin{equation}\label{eq:annular-ratio-limit}
 \int_a^bF_\nu(t)t\,\dd t
 \longrightarrow
 K\int_a^bt^{\rho+1}\,\dd t,
 \qquad
 0<a<b<2.
\end{equation}


\medskip\noindent
\textbf{\underline{Step 2. Pointwise convergence for \(0<t<2\):}}
Fix \(t\in(0,2)\), and choose \(h>0\) so small that
\(0<t-h<t<t+h<2\).
Since each \(F_\nu\) is nondecreasing,
\[
\frac{\int_{t-h}^tF_\nu(s)s\,\dd s}
     {\int_{t-h}^ts\,\dd s}
\leq
F_\nu(t)
\leq
\frac{\int_t^{t+h}F_\nu(s)s\,\dd s}
     {\int_t^{t+h}s\,\dd s}.
\]
Using \eqref{eq:annular-ratio-limit} in both quotients gives
\[
K
\frac{\int_{t-h}^ts^{\rho+1}\,\dd s}
     {\int_{t-h}^ts\,\dd s}
\leq
\liminf_{\nu\to\infty}F_\nu(t)
\leq
\limsup_{\nu\to\infty}F_\nu(t)
\leq
K
\frac{\int_t^{t+h}s^{\rho+1}\,\dd s}
     {\int_t^{t+h}s\,\dd s}.
\]
Letting \(h\downarrow0\) gives
\begin{equation}\label{eq:subsequence-ratio-limit}
 F_\nu(t)\longrightarrow Kt^\rho,
 \qquad
 0<t<2.
\end{equation}
At \(t=1\), \(F_\nu(1)=1\) for every \(\nu\). Hence
\eqref{eq:subsequence-ratio-limit} gives \(K=1\), and therefore
\[
F_\nu(t)\longrightarrow t^\rho,
\qquad
0<t<2.
\]


\medskip\noindent
\textbf{\underline{Step 3. Removal of the subsequence:}}
The original sequence \(r_\nu\to\infty\) was arbitrary. Suppose that
for some \(t\in(0,2)\) the full limit
\(T(tr)/T(r)\to t^\rho\) failed. Then there would be a sequence
\(r_\nu\to\infty\) and an \(\eta>0\) such that
\[
\left|
\frac{T(tr_\nu)}{T(r_\nu)}
-
t^\rho
\right|
\geq\eta
\]
for all \(\nu\). Applying Steps~1--2 to this sequence would produce a
subsequence on which the ratio converges to \(t^\rho\), a
contradiction. Hence
\begin{equation}\label{eq:ratio-below-two}
 \frac{T(tr)}{T(r)}
 \longrightarrow t^\rho,
 \qquad
 0<t<2.
\end{equation}


\medskip\noindent
\textbf{\underline{Step 4. All positive multipliers and local uniformity:}}
Let \(c\geq2\). Choose an integer \(k\geq1\) such that
\(b=c^{1/k}\in(1,2)\). Then
\[
\frac{T(cr)}{T(r)}
=
\prod_{j=0}^{k-1}
\frac{T(b^{j+1}r)}{T(b^jr)}.
\]
For every fixed \(j\), the base radius \(b^jr\) tends to infinity.
Applying \eqref{eq:ratio-below-two} with multiplier \(b\) to each
factor gives the limit \(b^\rho\). Hence the product tends to
\(b^{k\rho}=c^\rho\). Together with
\eqref{eq:ratio-below-two}, this proves \eqref{eq:ratio-limit} for
every \(c>0\).

It remains to prove local uniformity. Let
\([a,b]\Subset(0,\infty)\) and \(\eta>0\). Since \(c^\rho\) is
uniformly continuous on \([a,b]\), choose a partition
\(a=c_0<c_1<\cdots<c_m=b\) so fine that
\(c_{j+1}^\rho-c_j^\rho<\eta\).
Pointwise convergence at the finitely many points \(c_j\) is uniform
in \(j\) for all sufficiently large \(r\). If
\(c_j\leq c\leq c_{j+1}\), monotonicity of \(T(cr)\) in \(c\) gives
\[
\frac{T(c_jr)}{T(r)}
\leq
\frac{T(cr)}{T(r)}
\leq
\frac{T(c_{j+1}r)}{T(r)}.
\]
Using the convergence at \(c_j,c_{j+1}\) and the choice of the
partition proves uniform convergence on \([a,b]\). This establishes
the local uniformity asserted in \eqref{eq:ratio-limit}.
\end{proof}


We can now complete the proof of Theorem~\ref{thm:main}.

\begin{proof}[Proof of Theorem~\ref{thm:main}]
If \(T(r,\fv)=O(\log r)\), then \(\rho(\fv)=0<1\).
Proposition~\ref{prop:small-order}, together with the small-ramification
hypothesis, would then imply that \([\fv]\) is a rational normal curve,
contrary to the transcendence assumption. Hence the second alternative
in Lemma~\ref{lem:logdichotomy} holds:
\[
\log r=o(T(r,\fv)).
\]

Proposition~\ref{prop:indices} now gives
\[
\rho_*(T)
=
\lambda(\fv)
=
\rho(\fv)
=
\rho^*(T)
\in
\{0\}\cup\Gn.
\]
The value \(0\) is excluded once more by
Proposition~\ref{prop:small-order} and transcendence. Thus the common
value \(\rho\) belongs to \(\Gn\).

Proposition~\ref{prop:regular-variation} gives
\(T(cr,\fv)/T(r,\fv)\to c^\rho\), locally uniformly for
\(c\in(0,\infty)\). Defining
\(\ell(r)=T(r,\fv)/r^\rho\), we therefore obtain
\[
\frac{\ell(cr)}{\ell(r)}
=
c^{-\rho}
\frac{T(cr,\fv)}{T(r,\fv)}
\longrightarrow1
\]
locally uniformly for \(c\in(0,\infty)\). Thus \(\ell\) is slowly
varying, and
\(T(r,\fv)=r^\rho\ell(r)\), which proves \eqref{reg}.
\end{proof}


\section{Sharpness and radial area}\label{sec:sharpness}

We first show that every positive order allowed by
Theorem~\ref{thm:main} actually occurs. The examples come from linear
differential equations with a single polynomial coefficient. The
construction is explicit and does not require asymptotic integration.

We then consider the rational normal curve, which realizes the
order-zero alternative, and finally derive the radial area asymptotics
from the regular variation proved in Section~\ref{sec:regular}.


\subsection{Realization of the admissible positive orders}

\begin{thm}\label{thm:realization}
For every \(\rho\in\Gn\), there is a reduced vector
\(\fv=(f_0,\ldots,f_{n})\) of entire functions with
\(W_{\fv}\equiv1\) such that \([\fv]\) is transcendental and
\begin{equation}\label{eq:realization-T-growth}
 c r^\rho
 \leq
 T(r,\fv)
 \leq
 C r^\rho
\end{equation}
for all sufficiently large \(r\), where \(c,C>0\).
\end{thm}

\begin{proof}

\medskip\noindent
\textbf{\underline{Step 1. The realizing differential equation:}}
Since \(\rho\in\Gn\), there are \(2\leq q\leq n+1\) and
\(m\in\Z_{\geq0}\) such that
\(\rho=1+m/q\). Put
\(\ell=m+q\), so that \(\rho=\ell/q\), and consider
\begin{equation}\label{eq:realizing-equation}
 y^{(n+1)}
 =
 z^m y^{({n+1}-q)},
 \qquad
 f_j^{(i)}(0)=\delta_{ij},
 \qquad
 0\leq i,j\leq n.
\end{equation}
The coefficient \(z^m\) is entire, so the normalized solutions
\(f_0,\ldots,f_{n}\) extend to entire functions on \(\C\).

The coefficient of \(y^{(n)}\) in
\eqref{eq:realizing-equation} is zero. Abel's identity therefore shows
that \(W_{\fv}\) is constant. The initial conditions in
\eqref{eq:realizing-equation} give \(W_{\fv}(0)=1\), and hence
\begin{equation}\label{eq:realization-wronskian}
 W_{\fv}\equiv1.
\end{equation}
By \eqref{eq:realization-wronskian}, the components are linearly independent and have no
common zero, so \(\fv\) is reduced.


\medskip\noindent
\textbf{\underline{Step 2. Power series for the reduced \(q\)-th order equation:}}
Put \(p={n+1}-q\). If \(u=y^{(p)}\), then
\eqref{eq:realizing-equation} becomes
\[
u^{(q)}=z^m u.
\]
For \(0\leq j\leq q-1\), let \(u_j\) be the solution satisfying
\(u_j^{(i)}(0)=\delta_{ij}\), \(0\leq i\leq q-1\).

Write \(u_j(z)=\sum_{s\geq0}d_sz^s\). Comparing coefficients in
\(u_j^{(q)}=z^m u_j\) shows that the only nonzero coefficients occur
at exponents \(j+k\ell\), \(k\geq0\), and
\begin{equation}\label{eq:realizing-series}
 [z^{j+k\ell}]u_j(z)
 =
 \frac1{j!}
 \prod_{t=1}^k
 \frac{(j+t\ell-q)!}{(j+t\ell)!},
 \qquad
 k\geq0.
\end{equation}
Here the product is equal to \(1\) when \(k=0\).

For \(t\geq1\),
\[
\frac{(j+t\ell-q)!}{(j+t\ell)!}
=
\frac1{
 (j+t\ell-q+1)\cdots(j+t\ell)
 }.
\]
The \(q\) factors in the denominator are bounded above and below by
positive constant multiples of \(t\), with constants depending only on
\(m,q,j\). Consequently there are constants
\(c_0,C_0>0\) such that
\begin{equation}\label{eq:realizing-coefficient-bounds}
 \frac{c_0^k}{j!(k!)^q}
 \leq
 [z^{j+k\ell}]u_j(z)
 \leq
 \frac{C_0^k}{j!(k!)^q},
 \qquad
 k\geq0.
\end{equation}
In particular, all nonzero coefficients are positive.


\medskip\noindent
\textbf{\underline{Step 3. Growth of the solutions \(u_j\):}}
We first prove the upper bound. By
\eqref{eq:realizing-coefficient-bounds},
\[
M(r,u_j)
\leq
\frac{r^j}{j!}
\sum_{k=0}^\infty
\frac{C_0^kr^{k\ell}}{(k!)^q}.
\]
Put \(x=C_0^{1/q}r^{\ell/q}=C_0^{1/q}r^\rho\). Since
\[
\sum_{k=0}^\infty\frac{x^{qk}}{(k!)^q}
\leq
\left(
\sum_{k=0}^\infty\frac{x^k}{k!}
\right)^q
=
e^{qx},
\]
we obtain \(\log M(r,u_j)\leq C_1r^\rho\) for all sufficiently large
\(r\).

For the lower bound, put
\(x=c_0^{1/q}r^\rho\) and
\(k=\lfloor x/2\rfloor\). Since the coefficients in
\eqref{eq:realizing-series} are positive,
\[
M(r,u_j)
\geq
\frac{r^j}{j!}
\frac{x^{qk}}{(k!)^q}.
\]
For large \(r\), \(k\leq x/2\), and \(k!\leq k^k\). Hence
\(x^k/k!\geq(x/k)^k\geq2^k\). Since
\(k\geq x/3\) for all sufficiently large \(x\), it follows that
\(\log M(r,u_j)\geq c_1r^\rho\).

We have therefore proved
\begin{equation}\label{eq:series-growth}
 c_1r^\rho
 \leq
 \log M(r,u_j)
 \leq
 C_1r^\rho,
 \qquad
 0\leq j\leq q-1,
\end{equation}
for all sufficiently large \(r\). The upper estimate also proves
directly that the series in \eqref{eq:realizing-series} converges for
every \(z\in\C\).


\medskip\noindent
\textbf{\underline{Step 4. Growth of the solutions of the original equation:}}
Recall that \(p={n+1}-q\). The normalized solutions of
\eqref{eq:realizing-equation} have the following form. For
\(0\leq j<p\),
\[
f_j(z)=\frac{z^j}{j!}.
\]
For \(0\leq j<q\), the solution \(f_{p+j}\) satisfies
\(f_{p+j}^{(p)}=u_j\) and
\(f_{p+j}^{(i)}(0)=0\) for \(0\leq i<p\).

If \(y^{(p)}=u\) and \(y^{(i)}(0)=0\) for \(0\leq i<p\), repeated
integration along the radial segment gives
\(M(r,y)\leq r^pM(r,u)/p!\).
Conversely, Cauchy's estimate for the \(p\)-th derivative on
\(D_{2r}\) gives
\(M(r,u)\leq p!M(2r,y)/r^p\). Thus
\begin{equation}\label{eq:antiderivative-comparison}
 M(r,y)
 \leq
 \frac{r^p}{p!}M(r,u),
 \qquad
 M(r/2,u)
 \leq
 \frac{2^pp!}{r^p}M(r,y).
\end{equation}

Applying \eqref{eq:antiderivative-comparison} to
\(y=f_{p+j}\) and \(u=u_j\), and then using
\eqref{eq:series-growth}, gives
\begin{equation}\label{eq:nonpolynomial-component-growth}
 c_2r^\rho
 \leq
 \log M(r,f_{p+j})
 \leq
 C_2r^\rho,
 \qquad
 0\leq j<q,
\end{equation}
for all sufficiently large \(r\). Indeed, the factors \(r^p\) in
\eqref{eq:antiderivative-comparison} contribute only \(O(\log r)\) to
the logarithm, whereas \(\rho\geq1\).

The remaining \(p\) components are polynomials. Hence
\eqref{eq:nonpolynomial-component-growth} also implies
\begin{equation}\label{eq:vector-maximum-growth}
 c_3r^\rho
 \leq
 \log\max_{0\leq j\leq n}M(r,f_j)
 \leq
 C_3r^\rho.
\end{equation}


\medskip\noindent
\textbf{\underline{Step 5. Passage from component growth to the projective characteristic:}}
For this step only, make a fixed unitary change of homogeneous
coordinates so that all components are nonzero at the origin. This
does not change \(\norm{\fv}\) or \(T(r,\fv)\), and it changes the
Wronskian only by a constant of modulus one. We continue to denote the
resulting vector by \(\fv\).

Equation \eqref{eq:vector-maximum-growth} implies
\(\log M(r,f_j)=O(r^\rho)\) for every component. For \(i=0\), the
ratio \(f_j^{(0)}/f_j\) is identically \(1\). For each fixed
\(1\leq i\leq n\), set \(h_{j,r}(z):=f_j(rz)\). Since
\(T(4,h_{j,r})=O(r^\rho)\) and \(h_{j,r}(0)\ne0\),
Lemma~\ref{lem:logderivbound} gives
\[
 m\left(1,\frac{h_{j,r}^{(i)}}{h_{j,r}}\right)=O(\log r).
\]
Moreover,
\[
\frac{h_{j,r}^{(i)}(z)}{h_{j,r}(z)}
=r^i\frac{f_j^{(i)}(rz)}{f_j(rz)},
\]
and hence
\begin{equation}\label{eq:realization-logderiv}
 m\left(r,\frac{f_j^{(i)}}{f_j}\right)
 =m\left(1,r^{-i}\frac{h_{j,r}^{(i)}}{h_{j,r}}\right)
 \le m\left(1,\frac{h_{j,r}^{(i)}}{h_{j,r}}\right)
 =O(\log r).
\end{equation}


Since \(|W_{\fv}|=1\),
\[
1
=
\left|\prod_{j=0}^{n}f_j\right|
\left|
\det
\left(
\frac{f_j^{(i)}}{f_j}
\right)_{0\leq i,j\leq n}
\right|.
\]
Using
\(\prod_j|f_j|\leq\norm{\fv}^{\,{n+1}}\), we get almost everywhere
\[
\log^-\norm{\fv}
\leq
\frac1{n+1}
\log^+
\left|
\det
\left(
\frac{f_j^{(i)}}{f_j}
\right)
\right|.
\]
Expanding the determinant and applying
\eqref{eq:realization-logderiv} to its finitely many factors gives
\begin{equation}\label{eq:realization-negative-part}
 \mean r{\log^-\norm{\fv}}
 =
 O(\log r).
\end{equation}

The function \(\log\norm{\fv}\) is subharmonic. The Poisson kernel
from \(D_{2r}\) to \(D_r\) is bounded by \(3\). Since
\(\mean{2r}{\log\norm{\fv}}
 =T(2r,\fv)+O(1)\), equation
\eqref{eq:realization-negative-part} gives
\begin{equation}\label{eq:realization-poisson}
 \max_{|z|\leq r}\log\norm{\fv(z)}
 \leq
 3T(2r,\fv)+O(\log r).
\end{equation}

By \eqref{eq:vector-maximum-growth},
\[
\max_{|z|\leq r}\log\norm{\fv(z)}
\geq
c_3r^\rho.
\]
Substitution into \eqref{eq:realization-poisson} gives
\(T(2r,\fv)\geq c_4r^\rho\) for large \(r\), and hence
\begin{equation}\label{eq:realization-T-lower}
 T(r,\fv)\geq c_5r^\rho.
\end{equation}

For the opposite inequality,
\(\log\norm{\fv}\leq
\log\sqrt {n+1}+\max_j\log^+|f_j|\).
Therefore \eqref{eq:vector-maximum-growth} gives
\begin{equation}\label{eq:realization-T-upper}
 T(r,\fv)\leq C_5r^\rho
\end{equation}
for all sufficiently large \(r\).

Equations \eqref{eq:realization-T-lower} and
\eqref{eq:realization-T-upper} prove
\eqref{eq:realization-T-growth}. If \([\fv]\) were rational, then a
reduced representation would have the form
\[
 \fv=e^g(p_0,\ldots,p_n),
\]
where \(g\) is entire and the \(p_j\) are polynomials without a common
zero. Since
\[
 W_{\fv}=e^{(n+1)g}W(p_0,\ldots,p_n)\equiv1,
\]
the function \(g\) must be constant. Consequently
\(T(r,\fv)=O(\log r)\), contradicting
\eqref{eq:realization-T-lower}. Thus \([\fv]\) is transcendental.
The two-sided
estimate also gives
\(\lambda(\fv)=\rho(\fv)=\rho\).
\end{proof}


\subsection{The rational normal example}

The order-zero case is realized by the rational normal curve
\begin{equation}\label{eq:rationalnormal-example}
 \fv_0(z)
 =
 (1,z,\ldots,z^{n}).
\end{equation}
Its Wronskian is the nonzero constant
\(\prod_{j=0}^{n}j!\), and
\begin{equation}\label{eq:rationalnormal-growth}
 T(r,\fv_0)
 =
 \frac12
 \log\left(
 \sum_{j=0}^{n}r^{2j}
 \right)
 =
 n\log r+O(1).
\end{equation}
A constant change of basis with the appropriate determinant changes
the Wronskian to \(1\) without changing the projective nature of the
example.

Thus Theorem~\ref{thm:realization} realizes every positive exponent in
\(\Gn\), while \eqref{eq:rationalnormal-example} realizes the
order-zero alternative of Proposition~\ref{prop:small-order}. Hence
the set of possible orders in Theorem~\ref{thm:main} cannot be
reduced.


\subsection{Radial area}

Regular variation also determines the asymptotic radial area of the
curve. We specify the normalization explicitly.

Since a reduced vector has no common zero,
\(\log\norm{\fv}\) is smooth and subharmonic. Define
\begin{equation}\label{eq:radial-area-definition}
 \mathcal A_{\fv}(r)
 :=
 \frac1{2\pi}
 \int_{D_r}
 \Delta\log\norm{\fv(z)}\,\mathrm dA(z).
\end{equation}
If \(\fv\) is multiplied by a common zero-free entire factor
\(e^g\), then \(\log\norm{\fv}\) changes by the harmonic function
\(\Ree g\). Thus \(\mathcal A_{\fv}\) is independent of this choice of
reduced representation.

Green's formula applied to
\eqref{eq:radial-area-definition} gives
\begin{equation}\label{eq:area-derivative}
 \mathcal A_{\fv}(r)
 =
 rT'(r,\fv).
\end{equation}
In particular, \(\mathcal A_{\fv}\) is nondecreasing.

\begin{cor}\label{cor:area}
Under the hypotheses of Theorem~\ref{thm:main}, a transcendental curve
of order \(\rho\) satisfies
\begin{equation}\label{eq:area-asymptotics}
 \mathcal A_{\fv}(r)
 \sim
 \rho T(r,\fv),
 \qquad
 \frac{\mathcal A_{\fv}(cr)}
      {\mathcal A_{\fv}(r)}
 \longrightarrow
 c^\rho
 \quad(c>0).
\end{equation}
For a curve of the form \eqref{eq:rationalnormal},
\[
\mathcal A_{\fv}(r)\longrightarrow n.
\]
\end{cor}

\begin{proof}

\medskip\noindent
\textbf{\underline{Step 1. Comparison of the area with difference quotients of \(T\):}}
Fix \(c>1\). By \eqref{eq:area-derivative},
\[
T(r)-T(r/c)
=
\int_{r/c}^{r}
\frac{\mathcal A_{\fv}(t)}{t}\,\dd t,
\qquad
T(cr)-T(r)
=
\int_r^{cr}
\frac{\mathcal A_{\fv}(t)}{t}\,\dd t.
\]
Since \(\mathcal A_{\fv}\) is nondecreasing,
\begin{equation}\label{eq:area-squeeze}
 \frac{T(r)-T(r/c)}{\log c}
 \leq
 \mathcal A_{\fv}(r)
 \leq
 \frac{T(cr)-T(r)}{\log c}.
\end{equation}


\medskip\noindent
\textbf{\underline{Step 2. Asymptotic size of the area:}}
Divide \eqref{eq:area-squeeze} by \(T(r)\).
Proposition~\ref{prop:regular-variation} gives
\[
\frac{T(cr)}{T(r)}\longrightarrow c^\rho,
\qquad
\frac{T(r/c)}{T(r)}\longrightarrow c^{-\rho}.
\]
Hence
\begin{equation}\label{eq:area-ratio-squeeze}
 \frac{1-c^{-\rho}}{\log c}
 \leq
 \liminf_{r\to\infty}
 \frac{\mathcal A_{\fv}(r)}{T(r)}
 \leq
 \limsup_{r\to\infty}
 \frac{\mathcal A_{\fv}(r)}{T(r)}
 \leq
 \frac{c^\rho-1}{\log c}.
\end{equation}
Letting \(c\downarrow1\) in
\eqref{eq:area-ratio-squeeze} gives
\begin{equation}\label{eq:area-to-T}
 \frac{\mathcal A_{\fv}(r)}{T(r)}
 \longrightarrow\rho.
\end{equation}


\medskip\noindent
\textbf{\underline{Step 3. Regular variation of the area:}}
For fixed \(c>0\), write
\[
\frac{\mathcal A_{\fv}(cr)}
     {\mathcal A_{\fv}(r)}
=
\frac{\mathcal A_{\fv}(cr)}
     {T(cr)}
\,
\frac{T(cr)}{T(r)}
\,
\frac{T(r)}
     {\mathcal A_{\fv}(r)}.
\]
The first and third factors tend to
\(\rho\) and \(1/\rho\), respectively, by
\eqref{eq:area-to-T}; the middle factor tends to \(c^\rho\) by
\eqref{eq:ratio-limit}. Therefore
\[
\frac{\mathcal A_{\fv}(cr)}
     {\mathcal A_{\fv}(r)}
\longrightarrow c^\rho.
\]
This proves \eqref{eq:area-asymptotics}.


\medskip\noindent
\textbf{\underline{Step 4. The rational normal curve:}}
For the rational normal curve,
\eqref{eq:rationalnormal-growth} gives
\[
T(e^x,\fv)=n x+O(1).
\]
The function \(x\mapsto T(e^x,\fv)\) is convex, and by
\eqref{eq:area-derivative} its derivative is
\(\mathcal A_{\fv}(e^x)\). A convex function differing from
\(n x\) by a bounded quantity has derivative tending to \(n\).
Hence
\[
\mathcal A_{\fv}(r)\longrightarrow n.
\]
\end{proof}

Thus the radial area asymptotics require no additional selection of
radii: they follow directly from regular variation and the monotonicity
of \(\mathcal A_{\fv}\).

Finally, the pointwise Wronskian zero count is also negligible. By
\eqref{eq:countbound},
\(n_{W_{\fv}}(r)\log2\leq
 \NN_{W_{\fv}}(2r)\).
The small-ramification hypothesis gives
\(\NN_{W_{\fv}}(2r)=o(T(2r,\fv))\), while
\eqref{eq:ratio-limit} gives
\(T(2r,\fv)/T(r,\fv)\to2^\rho\). Therefore
\[
n_{W_{\fv}}(r)=o(T(r,\fv)).
\]
Neither this conclusion nor Corollary~\ref{cor:area} requires
uniqueness of the angular subsequential profiles. In particular, no
asymptotic statement in individual sectors is needed.




\renewcommand{\bibliofont}{\fontsize{10}{11}\selectfont}
\begin{thebibliography}{Ere89a}
\raggedright

\bibitem[BR02]{BR02}
T.~Bergkvist and H.~Rullg\aa rd,
\emph{On polynomial eigenfunctions for a class of differential operators},
Math. Res. Lett. \textbf{9} (2002), no.~2--3, 153--171.
\doi{10.4310/MRL.2002.v9.n2.a3}.

\bibitem[Car33]{Car33}
H.~Cartan,
\emph{Sur les z\'eros des combinaisons lin\'eaires de $p$ fonctions
holomorphes donn\'ees},
Mathematica (Cluj) \textbf{7} (1933), 5--31.

\bibitem[Dra81]{Dra81}
D.~Drasin,
\emph{Quasi-conformal modifications of functions having deficiency sum two},
Ann.\ of Math. (2) \textbf{114} (1981), no.~3, 493--518.
\doi{10.2307/1971300}.

\bibitem[Dra87]{Dra87}
D.~Drasin,
\emph{Proof of a conjecture of F.~Nevanlinna concerning functions which
have deficiency sum two},
Acta Math. \textbf{158} (1987), no.~1--2, 1--94.
\doi{10.1007/BF02392256}.

\bibitem[DS72]{DS72}
D.~Drasin and D.~F.~Shea,
\emph{P\'olya peaks and the oscillation of positive functions},
Proc. Amer. Math. Soc. \textbf{34} (1972), no.~2, 403--411.
\doi{10.1090/S0002-9939-1972-0294580-X}.

\bibitem[Ere89a]{Ere89a}
A.~Eremenko,
\emph{A new proof of the Drasin theorem on meromorphic functions of finite
order with a maximal sum of defects.~I},
Teor. Funktsii Funktsional. Anal.\ i Prilozhen. No.~51 (1989), 107--116
(Russian); English translation, J. Soviet Math. \textbf{52} (1990),
no.~6, 3522--3529.

\bibitem[Ere89b]{Ere89b}
A.~Eremenko,
\emph{A new proof of the Drasin theorem on meromorphic functions of finite
order with a maximal sum of defects.~II},
Teor. Funktsii Funktsional. Anal.\ i Prilozhen. No.~52 (1989), 69--77
(Russian); English translation, J. Soviet Math. \textbf{52} (1990),
no.~5, 3397--3403.
\doi{10.1007/BF01099906}.

\bibitem[Ere93]{Ere93}
A.~Eremenko,
\emph{Meromorphic functions with small ramification},
Indiana Univ. Math. J. \textbf{42} (1993), no.~4, 1193--1218.
\doi{10.1512/iumj.1993.42.42055}.

\bibitem[Ere98]{Ere98}
A.~Eremenko,
\emph{Extremal holomorphic curves for defect relations},
J. Anal. Math. \textbf{74} (1998), 307--323.
\doi{10.1007/BF02819454}.

\bibitem[Ere02]{Ere02}
A.~Eremenko,
\emph{Value distribution and potential theory},
Proceedings of the International Congress of Mathematicians
(Beijing, 2002), vol.~II, Higher Education Press, Beijing, 2002,
pp.~681--690.

\bibitem[Ere14]{Ere14}
A.~Eremenko,
\emph{On the second main theorem of Cartan},
Ann. Acad. Sci. Fenn. Math. \textbf{39} (2014), no.~2, 859--871.
\doi{10.5186/aasfm.2014.3937}.

\bibitem[Ere15a]{Ere15a}
A.~Eremenko,
\emph{Correction to the paper ``On the second main theorem of Cartan''},
Ann. Acad. Sci. Fenn. Math. \textbf{40} (2015), no.~1, 503--506.
\doi{10.5186/aasfm.2015.4030}.

\bibitem[Ere15b]{Ere15b}
A.~Eremenko,
\emph{Holomorphic curves with few inflection points},
unpublished problem note, April~4, 2015, 2~pp.
\url{https://www.math.purdue.edu/~eremenko/dvi/wronsk.pdf}.

\bibitem[ES91]{ES91}
A.~Eremenko and M.~Sodin,
\emph{Meromorphic functions of finite order with a maximum sum of defects},
Teor. Funktsii Funktsional. Anal.\ i Prilozhen. No.~55 (1991), 84--95
(Russian); English translation, \emph{Meromorphic functions of finite
order with maximal deficiency sum}, J. Soviet Math. \textbf{59} (1992),
no.~1, 643--651.
\doi{10.1007/BF01102487}.

\bibitem[ES92]{ES92}
A.~Eremenko and M.~Sodin,
\emph{Distribution of values of meromorphic functions and meromorphic
curves from the standpoint of potential theory},
St. Petersburg Math. J. \textbf{3} (1992), no.~1, 109--136;
Russian original, Algebra i Analiz \textbf{3} (1991), no.~1, 131--164.

\bibitem[Fol99]{Fol99}
G.~B.~Folland,
\emph{Real analysis: Modern techniques and their applications},
2nd ed., Wiley-Interscience, New York, 1999.

\bibitem[Fre53]{Fre53}
M.~Frei,
\emph{Sur l'ordre des solutions enti\`eres d'une \'equation diff\'erentielle
lin\'eaire},
C. R. Acad. Sci. Paris \textbf{236} (1953), 38--40.

\bibitem[GH04]{GH04}
G.~G.~Gundersen and W.~K.~Hayman,
\emph{The strength of Cartan's version of Nevanlinna theory},
Bull. London Math. Soc. \textbf{36} (2004), no.~4, 433--454.
\doi{10.1112/S0024609304003418}.

\bibitem[Hay64]{Hay64}
W.~K.~Hayman,
\emph{Meromorphic functions},
Clarendon Press, Oxford, 1964.

\bibitem[Hor03]{Hor03}
L.~H\"ormander,
\emph{The analysis of linear partial differential operators.~I:
Distribution theory and Fourier analysis},
Classics in Mathematics, Springer-Verlag, Berlin, 2003,
reprint of the second (1990) edition.
\doi{10.1007/978-3-642-61497-2}.

\bibitem[KP26]{KP26}
S.~N.~Karp and K.~Purbhoo,
\emph{Universal Pl\"ucker coordinates for the Wronski map and positivity
in real Schubert calculus},
J.~Amer.~Math.~Soc. (2026), in press.
\doi{10.1090/jams/1087}.

\bibitem[MTV13]{MTV13}
E.~Mukhin, V.~Tarasov, and A.~Varchenko,
\emph{Bethe subalgebras of the group algebra of the symmetric group},
Transform. Groups \textbf{18} (2013), no.~3, 767--801.
\doi{10.1007/s00031-013-9232-y}.

\bibitem[Nev30]{Nev30}
F.~Nevanlinna,
\emph{\"Uber eine Klasse meromorpher Funktionen},
Septi\`eme Congr\`es des Math\'ematiciens Scandinaves
(Oslo, 1929), A.~W.~Br\o ggers Boktrykkeri, Oslo, 1930, pp.~81--83.

\bibitem[Noc83]{Noc83}
E.~I.~Nochka,
\emph{On the theory of meromorphic functions},
Dokl. Akad. Nauk SSSR \textbf{269} (1983), no.~3, 547--552 (Russian);
English translation under the title \emph{On the theory of meromorphic curves},
Soviet Math. Dokl. \textbf{27} (1983), no.~2, 377--381.

\bibitem[Pet84]{Pet84}
V.~P.~Petrenko,
\emph{Entire curves},
Vyshcha Shkola, Kharkov, 1984 (Russian).

\bibitem[Pfl46]{Pfl46}
A.~Pfluger,
\emph{Zur Defektrelation ganzer Funktionen endlicher Ordnung},
Comment. Math. Helv. \textbf{19} (1946), 91--104.
\doi{10.1007/BF02565950}.

\bibitem[Pur23]{Pur23}
K.~Purbhoo,
\emph{An identity in the Bethe subalgebra of $\C[\mathfrak S_n]$},
Proc. Lond. Math. Soc. (3) \textbf{127} (2023), no.~5, 1247--1267.
\doi{10.1112/plms.12560}.

\bibitem[Ru21]{Ru21}
M.~Ru,
\emph{Nevanlinna theory and its relation to Diophantine approximation},
2nd ed., World Scientific Publishing Co., Hackensack, NJ, 2021.
\doi{10.1142/12188}.

\bibitem[She85]{She85}
D.~F.~Shea,
\emph{On the frequency of multiple values of a meromorphic function of
small order},
Michigan Math. J. \textbf{32} (1985), no.~1, 109--116.
\doi{10.1307/mmj/1029003137}.

\bibitem[Voj97]{Voj97}
P.~Vojta,
\emph{On Cartan's theorem and Cartan's conjecture},
Amer. J. Math. \textbf{119} (1997), no.~1, 1--17.
\doi{10.1353/ajm.1997.0009}.

\bibitem[Was65]{Was65}
W.~Wasow,
\emph{Asymptotic expansions for ordinary differential equations},
John Wiley \& Sons, New York, 1965.

\bibitem[Wei69]{Wei69}
A.~Weitsman,
\emph{Meromorphic functions with maximal deficiency sum and a conjecture
of F.~Nevanlinna},
Acta Math. \textbf{123} (1969), 115--139.
\doi{10.1007/BF02392387}.

\end{thebibliography}
\appendix
\section{A differential polynomial identity}\label{app:differential-polynomial}
The following differential-polynomial identity is used in the proof of Lemma~\ref{lem:logderivbound}.

\begin{lem}\label{lem:differential-polynomial}
Let \(h\) be meromorphic and put
$
L={h'}/{h}.
$
Then, for every integer \(k\ge1\),
\begin{equation}\label{eq:bell-polynomial}
P_k:=\frac{h^{(k)}}{h}=
\sum_{\substack{m_0,\ldots,m_{k-1}\ge0\\
\sum_{j=0}^{k-1}(j+1)m_j=k}}
C_{\mathbf m}
\prod_{j=0}^{k-1}
\bigl(L^{(j)}\bigr)^{m_j},
\end{equation}
where
$
C_{\mathbf m}
=
{k!}/
{\Big(\displaystyle\prod_{j=0}^{k-1}
m_j!\,((j+1)!)^{m_j}\Big)}.
$
The multi-index is \(\mathbf m=(m_0,\ldots,m_{k-1})\).
\end{lem}

\begin{proof}
Fix a point \(z\) which is neither a zero nor a pole of \(h\), and choose
a local branch \(G=\log h\) near \(z\). Then
$$
G'=L,\qquad
G^{(j+1)}=L^{(j)},\qquad j\ge0.
$$
For sufficiently small \(t\),
$$
\frac{h(z+t)}{h(z)}
=
\exp\bigl(G(z+t)-G(z)\bigr).
$$
Expanding \(G\) at \(z\), we obtain
$$
G(z+t)-G(z)
=
\sum_{j\ge0}
\frac{L^{(j)}(z)}{(j+1)!}\,t^{j+1},
$$
and therefore
\begin{equation}\label{eq:generating-differential-polynomial}
\frac{h(z+t)}{h(z)}
=
\prod_{j\ge0}
\exp\left(
 \frac{L^{(j)}(z)}{(j+1)!}\,t^{j+1}
\right).
\end{equation}
For each \(j\),
$$
\exp\left(
\frac{L^{(j)}(z)}{(j+1)!}\,t^{j+1}
\right)
=
\sum_{m_j\ge0}
\frac{\bigl(L^{(j)}(z)\bigr)^{m_j}}
{m_j!\,((j+1)!)^{m_j}}
t^{(j+1)m_j}.
$$
Hence the coefficient of \(t^k\) on the right-hand side of
\eqref{eq:generating-differential-polynomial} is
$$
\sum_{\substack{m_0,\ldots,m_{k-1}\ge0\\
\sum_{j=0}^{k-1}(j+1)m_j=k}}
\frac{\displaystyle
\prod_{j=0}^{k-1}
\bigl(L^{(j)}(z)\bigr)^{m_j}}
{\displaystyle
\prod_{j=0}^{k-1}
m_j!\,((j+1)!)^{m_j}}.
$$

On the other hand, Taylor expansion of the left-hand side gives
$$
\frac{h(z+t)}{h(z)}
=
\sum_{k\ge0}
\frac{h^{(k)}(z)}{k!\,h(z)}t^k.
$$
Comparing the coefficients of \(t^k\) and multiplying by \(k!\) proves
\eqref{eq:bell-polynomial} at every point which is neither a zero nor a
pole of \(h\). Since both sides of \eqref{eq:bell-polynomial} are
meromorphic, the identity holds everywhere by meromorphic continuation.
\end{proof}

\end{document}
